% concatenated from main.tex, xuelin_style/xuelin_macros.tex, symbols.tex, sections/0-abstract.tex, sections/1-intro.tex, sections/2-method.tex, sections/3-alg.tex, sections/3-input_alg_population.tex, sections/4-consistency.tex, sections/5-conformal_prediction.tex, sections/5-conformal_classification.tex, sections/6-experiment.tex, sections/6-alg_empirical.tex, sections/3-input_alg_empirical.tex, artifacts/task1_multivariate_mc_summary.tex, sections/related_works.tex, sections/7-discussion.tex, sections/appendix_organization.tex, sections/appendix/appendix_proofs_organization.tex, sections/proofs/proofs_3-alg/proof_thm_algorithm_population.tex, sections/proofs/proofs_4-consistency/proof_thm_consistency_prediction.tex, sections/proofs/proofs_4-consistency/proof_lem_sample-size-interval.tex, sections/proofs/proofs_4-consistency/proof_lem_dist-regularity.tex, sections/proofs/proofs_4-consistency/proof_lem_conditional-dist-shift.tex, sections/proofs/proofs_5-conformal/proof_thm_tp.tex, sections/proofs/proofs_5-conformal/proof_thm_mc_error.tex, sections/proofs/proofs_5-conformal/proof_thm_mc_optimality.tex, sections/proofs/proofs_5-conformal/proof_thm_mc_same_sample.tex, sections/appendix/appendix_experiments_baselines.tex, sections/appendix/appendix_experiments_metrics.tex, sections/appendix/appendix_task1_setup.tex, sections/appendix/appendix_task2_setup.tex, sections/appendix/appendix_ver_thm5.2.tex, sections/appendix/appendix_algorithm.tex
\documentclass[11pt]{article}
% Xuelin's single-spaced paper style.
% Load this file immediately after \documentclass.

\ProvidesFile{xuelin_macros.tex}[2026/08/10 Xuelin paper macros and style]

% Page, language, and typography
\usepackage[
  letterpaper,
  left=1.5in,
  right=1.5in,
  top=1in,
  bottom=1in,
  footskip=0.6in
]{geometry}
\usepackage[english]{babel}
\usepackage[T1]{fontenc}
\usepackage{newtxtext}
\usepackage{microtype}
\usepackage{setspace}
\usepackage{titlesec}
\usepackage{etoolbox}

% Mathematics and theorem environments
\usepackage{amsmath,amssymb,amsfonts,amsthm,mathtools}
\usepackage{newtxmath}
\usepackage{bbm}

% Figures, tables, and algorithms
\usepackage{graphicx}
\usepackage{caption}
\usepackage{subcaption}
\usepackage{booktabs}
\usepackage{multirow}
\usepackage{float}
\usepackage{algorithm}
\usepackage{algpseudocode}

% Utilities and references
\usepackage{comment}
\usepackage{xcolor}
\usepackage{xspace}
\usepackage{natbib}
\usepackage{authblk}
\usepackage[
  colorlinks=false,
  linkbordercolor={1 0 0},
  citebordercolor={0 1 0},
  urlbordercolor={0 1 1},
  pdfborder={0 0 1}
]{hyperref}
\usepackage{cleveref}

% Single-spaced text with a balanced statistics-journal page.
\setstretch{1}
\AtBeginDocument{%
  \setlength{\abovedisplayskip}{6pt plus 2pt minus 1pt}%
  \setlength{\belowdisplayskip}{6pt plus 2pt minus 1pt}%
  \setlength{\abovedisplayshortskip}{3pt plus 1pt}%
  \setlength{\belowdisplayshortskip}{4pt plus 1pt minus 1pt}%
  \allowdisplaybreaks[1]%
}

% Left-aligned, restrained heading hierarchy.
\titleformat{\section}
  {\large\bfseries}{\thesection.}{0.65em}{}
\titleformat{\subsection}
  {\normalsize\bfseries}{\thesubsection.}{0.55em}{}
\titleformat{\subsubsection}
  {\normalsize\itshape}{\thesubsubsection.}{0.55em}{}
\titlespacing*{\section}{0pt}{2.4ex plus 0.6ex minus 0.2ex}
  {1.0ex plus 0.2ex}
\titlespacing*{\subsection}{0pt}{1.9ex plus 0.4ex minus 0.2ex}
  {0.7ex plus 0.2ex}
\titlespacing*{\subsubsection}{0pt}{1.5ex plus 0.3ex minus 0.2ex}
  {0.5ex plus 0.1ex}

% Conventional theorem styling with the manuscript's shared numbering.
\newtheoremstyle{statplain}
  {6pt}{6pt}{\itshape}{}% space, body font, and indent
  {\bfseries}{.}{0.5em}{}% head font, punctuation, and spacing
\newtheoremstyle{statdefinition}
  {6pt}{6pt}{\normalfont}{}% space, body font, and indent
  {\bfseries}{.}{0.5em}{}% head font, punctuation, and spacing

\theoremstyle{statplain}
\newtheorem{theorem}{Theorem}[section]
\newtheorem{lemma}[theorem]{Lemma}
\newtheorem{proposition}[theorem]{Proposition}
\newtheorem{corollary}[theorem]{Corollary}
\newtheorem{claim}[theorem]{Claim}
\newtheorem{result}[theorem]{Result}

\theoremstyle{statdefinition}
\newtheorem{definition}[theorem]{Definition}
\newtheorem{assumption}[theorem]{Assumption}
\newtheorem{example}[theorem]{Example}
\newtheorem{remark}[theorem]{Remark}
\newtheorem{comments}[theorem]{Comment}

\newenvironment{proofsketch}
  {\begin{proof}[Proof sketch]\renewcommand{\qedsymbol}{}}
  {\end{proof}}

% Reusable mathematical notation
% Keep manuscript-specific notation in the paper's symbols.tex.
\DeclareMathOperator*{\argmin}{arg\,min}
\DeclareMathOperator*{\argmax}{arg\,max}
\DeclareMathOperator{\poly}{poly}
\DeclareMathOperator{\polylog}{polylog}
\DeclareMathOperator{\vecz}{vec}
\DeclareMathOperator{\diag}{diag}
\DeclareMathOperator{\rank}{rank}
\DeclareMathOperator{\fed}{fed}
\DeclareMathOperator{\Var}{Var}
\DeclareMathOperator{\Cov}{Cov}
\DeclareMathOperator{\acc}{acc}
\DeclareMathOperator{\test}{test}
\DeclareMathOperator{\train}{train}
\DeclareMathOperator{\tr}{tr}
\DeclareMathOperator{\KL}{KL}
\DeclareMathOperator{\TV}{TV}
\DeclareMathOperator{\sgn}{sgn}
\DeclareMathOperator{\obj}{obj}
\DeclareMathOperator{\eff}{eff}
\DeclareMathOperator{\peff}{pure}
\DeclareMathOperator{\supp}{supp}

% Common estimands in causal inference.
\DeclareMathOperator{\ATE}{ATE}
\DeclareMathOperator{\ATT}{ATT}
\DeclareMathOperator{\CATE}{CATE}
\DeclareMathOperator{\LATE}{LATE}

% Backward-compatible shorthand used in existing manuscripts.
\newcommand{\var}{\Var}
\newcommand{\cov}{\Cov}
\newcommand{\mc}{\mathrm{mc}}
\newcommand{\tp}{\mathrm{tp}}
\newcommand{\sg}{\mathrm{SG}}
\newcommand{\poa}{\mathrm{POA}}
\newcommand{\stat}{\mathrm{stat}}
\newcommand{\sumin}{\sum_{i=1}^n}
\newcommand{\ra}{\rightarrow}
\newcommand{\la}{\leftarrow}

% Number systems, probability, and independence.
\newcommand{\Z}{\mathbb{Z}}
\newcommand{\R}{\mathbb{R}}
\newcommand{\E}{\mathbb{E}}
\newcommand{\PP}{\mathbb{P}}
\newcommand{\indep}{\perp\!\!\!\perp}
\newcommand{\indicator}[1]{\mathbbm{1}\!\left\{#1\right\}}
\newcommand{\define}{\coloneqq}

% Frequently used bold symbols. Bold Greek letters are defined explicitly;
% unlike guido_macros.tex, \mathbf retains its standard LaTeX meaning.
\newcommand{\bX}{\mathbf{X}}
\newcommand{\bY}{\mathbf{Y}}
\newcommand{\bZ}{\mathbf{Z}}
\newcommand{\bG}{\mathbf{G}}
\newcommand{\bR}{\mathbf{R}}
\newcommand{\M}{\mathbf{M}}
\newcommand{\V}{\mathbf{V}}
\newcommand{\0}{\mathbf{0}}
\newcommand{\bx}{\mathbf{x}}
\newcommand{\by}{\mathbf{y}}
\newcommand{\bz}{\mathbf{z}}
\newcommand{\bw}{\mathbf{w}}
\newcommand{\btheta}{\boldsymbol{\theta}}
\newcommand{\bbeta}{\boldsymbol{\beta}}

% Frequently used calligraphic symbols.
\newcommand{\cA}{\mathcal{A}}
\newcommand{\contC}{\mathcal{C}}
\newcommand{\D}{\mathcal{D}}
\newcommand{\cI}{\mathcal{I}}
\newcommand{\I}{\mathcal{I}}
\newcommand{\J}{\mathcal{J}}
\newcommand{\mec}{\mathcal{M}}
\newcommand{\N}{\mathcal{N}}
\newcommand{\cP}{\mathcal{P}}
\newcommand{\Q}{\mathcal{Q}}
\newcommand{\cS}{\mathcal{S}}
\newcommand{\W}{\mathcal{W}}
\newcommand{\cX}{\mathcal{X}}
\newcommand{\X}{\mathcal{X}}

% Norms and general manuscript conveniences.
\newcommand{\norm}[1]{\left\lVert #1 \right\rVert}
\newcommand{\normsm}[1]{\lVert #1 \rVert}
\newcommand{\normtwo}[1]{\norm{#1}_2}
\newcommand{\normtwosm}[1]{\normsm{#1}_2}
\newcommand{\normF}[1]{\left\lVert #1 \right\rVert_{\mathrm{F}}}
\newcommand{\normast}[1]{\left\lVert #1 \right\rVert_{\ast}}
\newcommand{\wt}{\widetilde}
\newcommand{\property}{\mathbf{P}}
\newcommand{\groups}{\mathcal{G}}
\newcommand{\group}{g}
\newcommand{\capgroup}{G}
\newcommand{\groupalt}{j}
\newcommand{\risk}{\mathcal{R}}
\newcommand{\groupfxn}{r}
\newcommand{\ie}{\textit{i.e.}}
\newcommand{\eg}{\textit{e.g.}}

% Publication-style figures, tables, and captions.
\renewcommand{\arraystretch}{1.08}
\captionsetup{
  font=small,
  labelfont=bf,
  labelsep=period,
  justification=justified,
  singlelinecheck=false
}
\captionsetup[sub]{font=footnotesize,labelfont=bf,labelsep=period}
\captionsetup[table]{position=top,skip=5pt}
\captionsetup[figure]{
  position=bottom,
  skip=5pt,
  justification=justified,
  singlelinecheck=true
}

% A compact but readable title block and bibliography.
\renewcommand{\Authfont}{\normalfont}
\renewcommand{\Affilfont}{\normalfont}
\setlength{\affilsep}{1em}
\setlength{\bibsep}{2pt plus 0.3ex}
\AtBeginEnvironment{thebibliography}{\small\setstretch{1}}
\urlstyle{same}

% Venue-neutral BibTeX default. Redefine this in a manuscript when a venue
% requires its own bibliography style (for example, imsart-nameyear for AOS).
\providecommand{\xuelinbibstyle}{plainnat}

% Coauthor notes.
\definecolor{myGreen}{RGB}{80,180,0}
\definecolor{myGold}{rgb}{0.75,0.6,0.12}
\definecolor{myBlue}{rgb}{0.12,0.45,0.75}

\newcommand{\Baihe}[1]{{\color{myGreen}[Baihe: #1]}}
\newcommand{\Xuelin}[1]{{\color{myGold}[Xuelin: #1]}}
\newcommand{\Yilong}[1]{{\color{myBlue}[Yilong: #1]}}
\newcommand{\Mike}[1]{{\color{myGold}[Mike: #1]}}
\newcommand{\Guido}[1]{{\color{myGold}[Guido: #1]}}

% PICPI-specific notation for this manuscript.
\newcommand{\bigO}{\mathcal{O}}
\newcommand{\interval}{I} % single interval
\newcommand{\intervalset}{\I} % set of interval
\newcommand{\numInterval}{J} % number of total intervals
\newcommand{\iteratorInterval}{j} 

\newcommand{\numClass}{G} % number of total classes
\newcommand{\iteratorClass}{g} 

\newcommand{\mcThreshold}{\tau} % threshold for multi class
\newcommand{\numData}{n} % number of data

\newcommand{\failureProb}{\delta} % failure probability

\newcommand{\intervalMaxLen}{\Delta} % max length among intervals

\newcommand{\discreteNum}{K} % discretization number of [0,1]

\newcommand{\tre}{W} % treatment
\newcommand{\littletre}{w}

\newcommand{\estModel}{p} % estimated model

% title
\newcommand{\titleName}{Prediction-Interval-Conditional Prediction Interval\xspace}
\newcommand{\titleShort}{PICPI\xspace}
\newcommand{\titleNames}{Prediction-Interval-Conditional Prediction Intervals\xspace}
\newcommand{\titleShorts}{PICPIs\xspace}


% bandit
\newcommand{\expectedValue}{R}



\title{PICPIs: Prediction-Interval-Conditional Prediction Intervals}


\author[{}]{%
  Xuelin Yang\textsuperscript{1,}\thanks{Equal contribution. Imbens' work is supported by the Office of Naval Research under Grant N00014-17-1-2131; Jordan’s work is supported by European Union under Grant ERC-2022-SYG-OCEAN-101071601. Email correspondence: \texttt{xuelin@berkeley.edu}.  Code: \texttt{https://github.com/xyang23/picpi}.}\qquad
  Baihe Huang\textsuperscript{1,}\protect\footnotemark[1]\qquad
  Yilong Hou\textsuperscript{1,}\protect\footnotemark[1]\authorcr
  Guido Imbens\textsuperscript{2} \qquad
  Michael I. Jordan\textsuperscript{1,}\textsuperscript{3}\\
  {\Affilfont
    \textsuperscript{1}University of California, Berkeley\\
    \textsuperscript{2}Stanford University\\
    \textsuperscript{3}Inria Paris}
}
\date{}
\begin{document}

\maketitle
\begin{abstract}
A classical question in statistics is which observable quantities to
condition on when drawing inferences about unobservable targets.
For conformal prediction in nonparametric uncertainty quantification, standard marginal validity offers limited resolution at the prediction values on which decisions are based, and fully conditional guarantees with respect to the covariates are provably unattainable. 
We address this gap by introducing a prediction-based conditioning framework that we refer to as \titleNames (\titleShorts). Formally,  a \titleShort is an interval $I$ satisfying a self-consistency condition:
$$ \E [ Y\mid p(X) \in I] \in I,$$
for predictive model $p$, contextual covariate $X$, and outcome $Y$.
Thus, an interval simultaneously defines a stratum of prediction values and certifies that the mean outcome in that stratum lies in the same interval.
This self-consistency condition yields data-adaptive strata without altering the original prediction.
Such intervals can be constructed using practical algorithms.  
Under regularity of the prediction distribution, the constructed intervals cover all but an arbitrarily small fraction of prediction values and have widths that decrease at rate $n^{-1/3}$, up to logarithmic factors and the prediction error.
Moreover, identifying these locally calibrated intervals can, in turn, inform downstream decision-making.
We derive inference procedures for \titleShorts in probabilistic prediction and multi-class classification, accompanied by theoretical guarantees.
Empirical results are provided that compare \titleShorts with existing interval-based baselines.

% details on downstream:
%In particular, we show how \titleShorts yield prediction intervals for the true conditional probability, with marginal miscoverage controlled by within-stratum variation and uncovered prediction mass. For multiclass classification, we derive worst-case class-conditional coverage certificates and show that the label set constructed from these intervals attains set-size optimality given prescribed coverage.

%We introduce a prediction-conditioned uncertainty-quantification framework.
%Point predictions of event probabilities are routinely used for decision-making, but their reliability is commonly assessed with fixed bins for the scores or enforced by post-hoc recalibration.
%, and bandit problems.
% We establish theoretical coverage certificates for \titleShorts over the prediction range and non-asymptotic bounds on their widths, which vanish as the sample size grows under regularity conditions and base predictor consistency.
\end{abstract}
\noindent\textbf{Keywords:} Prediction-conditioned inference; conformal prediction; uncertainty quantification; calibration; distribution-free inference

\newpage
\section{Introduction}

A classical problem in statistical inference is that of
\emph{conditioning}---on what observable quantities should one condition in order to make inferences about unobservable quantities~\citep{fisher1934statistical,cox1958some,  cox1988some, reid1995}? 
This problem has received renewed attention in recent work on nonparametric uncertainty quantification,
specifically the area of \emph{conformal prediction}, which aims to produce confidence intervals for nonparametric prediction~\citep{vovk2005algorithmic,shafer2008tutorial}. Consider for example the following formal model for binary prediction. Let $X\in\mathcal{X}$ denote a vector of covariates, let $Y\in\{0,1\}$ denote an outcome, and let $p:\mathcal{X}\to[0,1]$ be a fitted model whose output $p(X)$ is intended to approximate the conditional probability $p^*(X):=\PP(Y=1\mid X)$. Assuming exchangeability for the observable data, $\{(X_i, Y_i)\}_{i=1}^n$, but making no other structural or probabilistic assumptions, conformal prediction yields distribution-free prediction sets that come with a marginal guarantee of the form:
\begin{align}\label{eq:intro_conformal_marginal}
\PP\Big(Y\in C(X)\Big)\ge 1-\alpha,
\end{align}
where $C(X)$ is a prediction set constructed from $p(X)$, and $\alpha \in (0,1)$ is a user-specified failure probability. As has been emphasized in recent years, this marginal guarantee is unconditional with respect to the covariate $X$, and thus does not provide direct assessment of uncertainty about a specific prediction; moreover, fully conditional coverage with respect to $X$ is generally unattainable without strong distributional assumptions.  More formally,
\citet{lei2014distribution} and \citet{barber2021limits} demonstrate that any prediction interval $C(X)$ satisfying the strict conditional validity guarantee
\begin{align}\label{eq:conditional_conformal}
    \PP\Big(Y \in C(x)\mid X=x\Big) \geq 1-\alpha \quad \text{for almost all } x,
\end{align}
must be vacuous.  A significant effort has been made to temper the message in this result by enforcing conditional or group-wise validity via coarsening of the covariate~\citep{ding2023class, martinezgil2024groups, gibbs2025conformal,duchi2025few}. 

We present an alternative perspective on the problem of obtaining conditional estimates of uncertainty. Rather than conditioning on the covariate $X$, we suggest conditioning on the prediction $p(X)$. The idea of prediction-based conditioning is not new---indeed, it is the foundation of the literature on probabilistic calibration~\citep{dawid1982wellcalibrated,foster1998asymptotic}.  Recall that a calibrated model is defined as one that satisfies the following conditional criterion:
\begin{align}\label{eq:intro_cali}
    \PP\Big(Y=1\mid p(X)=v\Big) = v.
\end{align} 
That is, among all instances where the predicted probability is approximately $v$, the outcome should occur with frequency approximately $v$. 

Our approach to conditional uncertainty quantification is built on a similar ``self-consistency'' property. Specifically, we ask that the probability of a positive outcome, conditioning on the model's prediction falling within a specific interval $I$, also lies within $I$. Formally, we have:
\begin{align}\label{eq:intro_picpi}
\PP \Big(Y=1 \mid p(X)\in I\Big)\in I.
\end{align}
Effectively, we are using the model's predictions to stratify risk. We refer to this approach as \emph{\titleNames (\titleShorts)}.  
As illustrated in Figure~\ref{fig:teaser}, each \titleShort serves a dual role: it simultaneously defines the stratum of predictions being evaluated and provides a calibrated uncertainty statement for that specific stratum. 
In this way, uncertainty quantification is aligned with prediction, producing a summary of model reliability across the prediction distribution. In contrast to the intervals defined in Eq.~\eqref{eq:conditional_conformal},we show (in Theorem~\ref{thm:consistency-prediction}) that the width of the \titleShort intervals converges to zero under suitable assumptions, providing increasingly precise characterizations of risk. 
%under consistency of the base predictor and suitable regularity conditions



\begin{figure}[tb]
    \centering
    \includegraphics[width=0.45\linewidth]{Figures/picpi_teaser_green.pdf}
    \caption{Illustration of the \titleShort methodology. 
    Each blue interval $I_j$ is both a stratum of predicted probabilities and a certificate that its empirical outcome frequency, shown in red, lies in $I_j$.
    The dashed diagonal line corresponds to perfect calibration. This example is generated under model misspecification (quadratic ground truth with a linear fitted model).}
    \label{fig:teaser}
\end{figure}

We make the following two distinctions in comparison with standard methodology in calibration.
First, calibration is usually assessed via binning. 
Binning schemes are typically based on \emph{fixed, external} partitions of the prediction space, with overall calibration error computed as an aggregation of the errors within each bin.
While useful for general model assessment, fixed bins need not align with the decision boundaries that arise when optimizing a downstream risk function. 
The proposed \titleShorts provide a local, data-adaptive diagnostics for  calibration.
Rather than forcing the model to conform to pre-defined bins through recalibration, we characterize regions of the prediction space where the model is  ``self-consistent'' and thus reliable. 

Second, calibration error computed from binning schemes is often used to guide \emph{post-hoc} recalibration procedures such as Platt scaling, isotonic regression, or self-calibration \citep{platt1999probabilistic,  van2024self,berta2024classifier}. 
Recalibration, however, can be undesirable \citep{guo2017calibration, pleiss2017fairness, carriero2025harms}. 
%MJ Reference?
With \titleShorts, a practitioner can instead identify regions of the prediction space that are directly actionable. We further demonstrate how the knowledge of these locally calibrated intervals can inform downstream decision-making, including probabilistic prediction and classification.
An extended discussion of these issues is provided in Section~\ref{sec:related_work}.



Our specific contributions are as follows: 
\begin{enumerate}
    \item We propose \titleShorts as a novel approach to conditional uncertainty quantification, we design algorithms to construct \titleShorts, and we provide theoretical characterizations of these intervals. In particular, we show that  the \titleShorts constructed via our algorithm cover all but a user-specified prediction mass have widths controlled according to a $\Tilde{\bigO}(n^{-1/3})$ term plus the prediction error (Theorem~\ref{thm:consistency-prediction}). Consequently, these widths vanish when the base predictor is consistent.
    \item We provide inference procedures for probabilistic prediction and multi-class classification, and present both theoretical guarantees and empirical results. Specifically, we show that probabilistic prediction using \titleShorts has marginal miscoverage that is bounded by the within-stratum variation of $p^*(X)$ together with the prediction mass not covered by the \titleShorts (Theorem~\ref{thm:tp}).
%Third, we showcase an application of \titleShorts in bandit problems, where the regret depends on the interval width.
Multi-class classification using \titleShorts has worst-case class-conditional coverage certificates (Theorem~\ref{thm:mc_error}), a set-size optimality result given a prescribed coverage threshold (Theorem~\ref{thm:mc_optimality}), which allows data re-use when constructing intervals (Theorem~\ref{thm:mc_same_sample}).
\end{enumerate}
 


The remainder of the paper is organized as follows: Section~\ref{sec:def-alg} introduces the \titleShort definition and presents its algorithmic instantiation.  
Section~\ref{sec:properties} derives theoretical properties of \titleShorts. 
Section~\ref{sec:tasks} presents an inference framework using \titleShorts for probabilistic prediction and classification. 
Section~\ref{sec:experiment} discusses the empirical performance of the \titleShort approach in comparison with several standard interval-based baselines.
%MJ Need to mention the Related WOrk section.
Section~\ref{sec:related_work} presents related work and we present our conclusions in Section~\ref{sec:conclusions}.



\section{Definition, Algorithmic Construction, and Validity}\label{sec:def-alg}

Let $Y\in\{0,1\}$, let $X\in\X$, and let $p:\X\to[0,1]$ be a fitted
model for $p^*(x):=\PP(Y=1\mid X=x)$. The model may be trained on a separate sample or supplied as an
off-the-shelf predictor.  A PICPI is an interval in the prediction space that both groups similar predictions and contains the
group's true underlying outcome probability.

\begin{definition}[\titleName (\titleShort)]\label{def:picpi}
A collection $\intervalset_p=\{\interval_j\}_{j=1}^{J}$ is called a
collection of \emph{\titleNames (\titleShorts)} with respect to $p$ if,
for every $j\in[J]$ with $\PP(p(X)\in\interval_j)>0$,
\begin{align}
\PP(Y=1\mid p(X)\in\interval_j)\in\interval_j.
\label{eq:picpi-definition}
\end{align}
\end{definition}

For simplicity of exposition we focus on binary outcomes in this section. We note that the \titleShort notion extends directly to regression with a bounded real-valued
outcome, as we will formalize in Remark~\ref{rmk:extension_regression}.


%\section{Constructing \titleShorts}\label{sec:alg}


\begin{algorithm}[bth] % changed form {algorithm*} (two columns) 
   \caption{\titleShort Construction (population mode with theoretical validity)} 
   \label{alg:calibration_population}
   \small
\begin{algorithmic}[1]
\Procedure{\textsc{Calibration}}{data points $\{(x_i,y_i)\}_{i=1}^\numData,$ predictor $ p$, failure probability $\delta$}

\Comment{Obtain a set of \titleShorts}
\State Let $\{0,1/\discreteNum,2/\discreteNum,\cdots,1\}$ be a discretization of $[0,1]$.
\State Initialize \titleShorts $\intervalset_p \la \emptyset$ 
% \State Let $\epsilon \leftarrow $
\For {$0\leq i<j \leq\discreteNum$} 
\State Let $[a,b]\la [i/\discreteNum, j/\discreteNum]$ \Comment{Evaluate the current interval}
\If{$[a,b]$ satisfies $\sum_{i=1}^n \mathbbm{1}(p(x_i) \in [a,b]) > 0$ and 
\begin{align*}
    \frac{\sum_{i=1}^n y_i \cdot \mathbbm{1}(p(x_i) \in [a,b])}{\sum_{i=1}^n \mathbbm{1}(p(x_i) \in [a,b])} \in \left[a + 2\sqrt{\frac{\log (\discreteNum^2/\delta)}{\sum_{i=1}^n \mathbbm{1}(p(x_i) \in [a,b])}}  ,b - 2\sqrt{\frac{\log (\discreteNum^2/\delta)}{\sum_{i=1}^n \mathbbm{1}(p(x_i) \in [a,b])}}\right]
\end{align*}} \label{line:empirical-interval}

\State $\intervalset_p \la \intervalset_p \cup \{[a, b]\}$ 

% \State Find all intervals $[a,b]$ in $\mathcal{H}$ satisfying 
% \begin{align*}
%     \frac{\sum_{i=1}^n y_i \cdot \mathbbm{1}(p(x_i) \in [a,b])}{\sum_{i=1}^n \mathbbm{1}(p(x_i) \in [a,b])} \in \left(a + 2\sqrt{\frac{\log (K/\delta)}{\sum_{i=1}^n \mathbbm{1}(p(x_i) \in [a,b])}} ,b - 2\sqrt{\frac{\log (K/\delta)}{\sum_{i=1}^n \mathbbm{1}(p(x_i) \in [a,b])}}\right)
% \end{align*}
% \Else
\EndIf
\EndFor
\State {\bf Return} \titleShorts $\intervalset_p$
\EndProcedure
\\
\Procedure{\textsc{Inference}}{covariate input $x,$ predictor $ p$, failure probability $\delta$}
\Comment{Find the shortest \titleShort that contains $p(x)$}
\State $\intervalset_p \leftarrow \textsc{Calibration}\left(\{(x_i,y_i)\}_{i=1}^n, p, \delta \right)$
\State Let $[a,b]$ be the shortest interval in $\intervalset_p$ that contains $p(x)$
\State {\bf Return} $[a,b]$
\EndProcedure
\end{algorithmic}
\end{algorithm}



We now present an algorithm to construct \titleShorts with an accompanying validity statement. 
The population-level algorithm, presented in pseudocode as Algorithm~\ref{alg:calibration_population}, comprises two procedures.  \textsc{Calibration} returns
a collection of intervals that we will show to satisfy the \titleShort condition (Theorem~\ref{thm:algorithm_population}).  Once this collection has been computed,
\textsc{Inference} reports the shortest returned interval containing a new
prediction $p(x)$. If a disjoint summary is preferred, we select a disjoint subset with maximum covered length through the \textsc{DisjointCalibration} procedure provided in Appendix~\ref{app:alg}.

 

The theoretical \titleShort validity of the produced intervals---that is, whether they satisfy Definition~\ref{def:picpi} defined on the population---follows from a concentration argument. In particular, for a candidate
interval $I=[a,b]$, the \titleShort condition asks whether the unknown
conditional mean
\(
\PP(Y=1\mid p(X)\in I)
\)
lies between $a$ and $b$.  We estimate this quantity by the average label
among calibration observations whose predictions fall in $I$. Because this
average could be noisy, it is not enough for it merely to lie in $I$: we require it
to be separated from both endpoints by a concentration margin (Line~\ref{line:empirical-interval}).  The margin is
larger for intervals containing fewer calibration observations, precisely
where the empirical average is less reliable.
%Algorithm~\ref{alg:calibration_population} proceeds as follows: in a grid of $\discreteNum+1$ points in $[0,1]$, it examines every interval whose endpoints are grid points and retains an interval only when its empirical label mean lies inside the interval with the concerntration margin (Line~\ref{line:empirical-interval}).  
This ensures that every interval admitted by the algorithm satisfies the population-level \titleShort property.
We formalize this statement of validity in the following theorem:
\begin{theorem}\label{thm:algorithm_population}

Suppose $\{(x_i,y_i)\}_{i=1}^{\numData}$ are drawn i.i.d. from the
distribution of $(X,Y)$ and $p$ is fixed independently of these observations.
Then, with probability at least $1-\failureProb$ over the calibration sample,
every interval returned by \textsc{Calibration} in
Algorithm~\ref{alg:calibration_population} is a \titleShort with respect to
$p$.

\end{theorem}

\begin{proofsketch}
We use Hoeffding's inequality on $\PP(Y=1\mid p(X)\in I)$ and a union bound over $\discreteNum^2$ discretized candidate intervals. For the full proof see Appendix~\ref{proof_thm:algorithm_population}.
\end{proofsketch}



Theorem~\ref{thm:algorithm_population} guarantees the PICPI property under the true underlying distribution. An empirical version of the algorithm that we will discuss in a later section (Algorithm~\ref{alg:calibration_empirical}, Section~\ref{sec:alg_emp}) produces PICPIs with respect to the empirical distribution of the calibration sample. 
Conceptually, their relationship is analogous to the distinction between a population CDF and an empirical CDF: while Theorem~\ref{thm:algorithm_population} enforces the PICPI property with respect to the  underlying population measure, Algorithm~\ref{alg:calibration_empirical} satisfies it under the plug-in sample measure.





\iffalse
Algorithm~\ref{alg:calibration_empirical}  presents an algorithm that constructs \titleShorts and performs predictive inference with a partition of $[0,1]$ as an empirical version of \titleShort.

\begin{theorem}
Algorithm~\ref{alg:calibration_empirical} returns an (empirically) \titleShort inference on prediction $p$ and is optimal in the following partial order: $\{\I_j\}_{j=1}^m \succeq \{\I_j'\}_{j=1}^{m'}$ if $\exists \{\I_j''\}_{j=1}^{m'}$ such that
\begin{itemize}
    \item There is a one-to-one map between $\{\I_j'\}_{j=1}^{m'}$ and $\{\I_j''\}_{j=1}^{m'}$ such that the corresponding bins have the same set of predictions.
    \item $\{\I_j\}_{j=1}^m$ is a refinement of $\{\I_j''\}_{j=1}^{m'}$.
\end{itemize}
\end{theorem}

\fi



% \begin{theorem}
% Algorithm~\ref{alg:calibration} returns a \titleShort on prediction $p$ and is optimal in the following partial order: $\{\I_j\}_{j=1}^m \succeq \{\I_j'\}_{j=1}^{m'}$ if $\exists \{\I_j''\}_{j=1}^{m'}$ such that
% \begin{itemize}
%     \item There is a one-to-one map between $\{\I_j'\}_{j=1}^{m'}$ and $\{\I_j''\}_{j=1}^{m'}$ such that the corresponding bins have the same set of predictions.
%     \item $\{\I_j\}_{j=1}^m$ is a refinement of $\{\I_j''\}_{j=1}^{m'}$.
% \end{itemize}
% \end{theorem}






% \section{calibrated predictive inference}

% Consider the same prediction task under distributional assumptions:
% \begin{itemize}
%     \item Train data: $(x_i,y_i) \in \mathcal{X} \times [0,1]$;
%     \item Predictor: $p: \mathcal{X} \to [0,1]$
%     \item IID assumption: $(x,y), (x_i,y_i) \sim_{i.i.d.} \mathcal{D}$.
% \end{itemize}

% We introduce the following concept to generalize the notion of conformal prediction.

% \begin{definition}[Calibrated predictive inference]
% We say that a set of bins $\{B_j\}_{j=1}^m$ is $\alpha$-calibrated predictive inference on prediction $p$, if
% \begin{itemize}
%     \item $\{B_j\}_{j=1}^m$ forms a partition of $[0,1]$;
%     \item For any $j \in [m]$, we have $\Pr \left(y \in B_j \big | p(x) \in B_j\right) \geq 1-\alpha$.
% \end{itemize}
% \end{definition}

% \begin{remark}[Relationship with conformal prediction]
% Conformal prediction~\citep{vovk2005algorithmic} achieves a marginal guarantee in the form of $\Pr \left(y \in [p(x)-l, p(x) + l] \right) \geq 1-\alpha$. 
% Calibrated predictive inference formalizes the predictive internals such as $[p(x)-l, p(x) + l]$ into disjoint bins $\{B_j\}_{j=1}^m$ and extends to a conditional notion on the prediction themselves.
% \end{remark}

\clearpage
\section{Properties of \titleShorts}
\label{sec:properties}

In this section we show that, when a predictor has error $\epsilon$ and its predicted probabilities satisfy a $\lambda$-regularity condition, most predictions will be enclosed by \titleShorts. As the calibration sample size grows, these intervals shrink while their union continues to cover most of the prediction mass; the rate depends on $\epsilon$, $\lambda$, and the calibration sample size.
%In this section we show that when the predictor is accurate ({\it i.e.}, the prediction error $\epsilon$ is small) and the true probabilities are well-spread (satisfying a $\lambda$-regularity condition), most predictions can be enclosed in \titleShorts that shrink as the dataset grows. 
%In other words, as we collect more data, we can construct \titleShorts that (i) remain valid (covering the true probability), (ii) become shorter in length, and (iii) cover most of the prediction mass. 
%We can thus achieve both reliability and precision, with the rate at which the intervals tighten depending on the prediction error and the sample size.
%we can guarantee that most predictions are contained within intervals whose length shrinks to zero as the sample size $n$ increases. Moreover, the rate of the length of intervals depends explicitly on the prediction error $\epsilon$ and the available sample size $n$.
% \begin{assumption}[Scaling of the prediction model]\label{asp:scaling-learning}
% There exists constant $\alpha, C > 0$ such that for any $\delta > 0$, with probability at least $1-\delta$,
% \begin{align*}
%     \E_\mu\Big[\Big|\estModel(Y) - \PP(Y = 1|X)\Big|\Big] \leq C \cdot \log(1/\delta) \cdot m^{-\alpha}.
% \end{align*}
% \end{assumption}
Our result requires a regularity condition ensuring that the predicted probabilities are sufficiently spread across $[0,1]$.
%, rather than concentrating on a few isolated values. 
%The following definition formalizes this by requiring that the distribution of predictions assigns at least a $\lambda$-fraction of the Lebesgue measure to every measurable set.


\begin{definition}[$\lambda$-regularity]
We say that a random variable $Z$ over $[0,1]$ is $\lambda$-\emph{regular} if $\PP(Z \in A) \geq \lambda \cdot \mu_B(A)$ for any measurable set $A$, where $\mu_B$ is the Lebesgue measure on $[0,1]$.
\end{definition}

%With this regularity condition in hand, 
%The following theorem shows that when the predictor is accurate (small $\epsilon$) and $\lambda$-regular, Algorithm~\ref{alg:calibration_population} produces \titleShorts whose lengths shrink at a rate determined by both the prediction error and the sample size, while their union still cover most of the prediction mass. Consequently, the intervals become increasingly tight around the predicted value and effectively approaching a point prediction, which implies consistency.


\begin{theorem}%[Consistency of prediction by \titleShorts]
[Finite-sample width guarantee for \titleShorts]
\label{thm:consistency-prediction}
% Fix $m,n \in \Z_+$ and $\delta_0,\delta_1 > 0$. Suppose there are $n$ data points $(x_i,y_i), i=1, \dots, n$ i.i.d. sampled from the ground-truth distribution $\mu$, and a prediction model $\estModel$ trained on $m$ data from the same distribution. 

Consider $\{(x_i,y_i)\}_{i=1}^{m+n}$  i.i.d.\ data points such that 
$n \geq 2^{-1/2} \lambda^{-1} \log(3n/\delta) $.
%hoeffding, replaced by chernoff:$n \geq \lambda^{-2} \log^{2}(3n/\delta)$.
Use $m$ data points to construct the predictor $\estModel: \X \to [0,1]$ and $n$ data points to produce \titleShorts $\intervalset_p$  via Algorithm~\ref{alg:calibration_population}. 
Suppose $X \sim \mu$, $\estModel(X)$ is $\lambda$-regular, and
\begin{align*}
    \E_\mu\Big[\Big| \estModel (X) - \PP(Y = 1 \mid X)\Big|\Big] \leq \epsilon,
\end{align*}
for some $\epsilon>0$.

% \Xuelin{deal with $x_i$ not from $[0,1]$ and discrete. Also do te next $m$ data needs to be sampled from the same dist? talk about how $\epsilon$ is connected to the rate for and give some examples of convergence rate of common estimators (for example $\epsilon=C/n^{-1/4}$ }
For any $\delta,\delta_1 > 0$, with probability at least $1-\delta$ over the randomness of $\{(x_i,y_i)\}_{i=1}^n$,
\begin{align}
    \PP_{x \sim \mu}\Big(\exists [a,b] \in \intervalset_p ~s.t.~ |b-a| \leq C(n,\epsilon,\delta,\delta_1) ~\text{and}~ p(x) \in [a,b] \Big) \geq 1-\delta_1,
\end{align}
where 
\begin{align}\label{eq:width_function}
    C(n,\epsilon,\delta,\delta_1) = c_0 \log(1/\delta_1) \cdot \Big(\Big(\frac{\log(n/\delta)}{\lambda n}\Big)^{1/3}  + \frac{\epsilon}{\delta_1}\Big)
\end{align}
for a universal constant $c_0$.
% \begin{align}\label{eq:width_function}
%     C(n,\epsilon,\delta,\delta_1) \leq 8\log(3/\delta_1) \cdot \Big(2\Big(\frac{2\log(3n/\delta)}{\lambda n}\Big)^{1/3}  + \frac{2\epsilon}{\delta_1}\Big).
% \end{align}

% Under Assumption~\ref{asp:scaling-learning} holds with constant $\alpha > 0$, Algorithm~\ref{alg:calibration_population} with input $\{(x_i,y_i)\}_{i=m+1}^{m+n}, \estModel$ returns a set of intervals $\I$ satisfying: for any $\delta>0$, with probability at least $1-\delta$ 
% \begin{align*}
%     1-C\cdot(e^{-\beta m} + e^{-\gamma n})
% \end{align*} 
% for $x$ sampled from $\mu$, with probability at least $1-\delta_1$ there exists an interval $[a,b]$ in $\I$ such that $\estModel \in [a,b]$ and
% \begin{align*}
%     |b-a| \leq 
% \end{align*}
\end{theorem}

\begin{proofsketch}
    First, we construct an initial collection of intervals of controlled length and desired coverage by applying Lemma~\ref{lem:dist-regularity} to the distribution of $p(X)$. That is, there exist intervals $\{[a_j,b_j]\}$ where the conditional mean prediction $\E[p(X) \mid p(X)\in [a_j, b_j]]$ is bounded away from its own endpoints. Then we transfer this property from the conditional mean prediction $\E[p(X)\mid p(X)\in [a_j, b_j]]$ to the conditional label probability $\E[Y \mid p(X)\in [a_j, b_j]]$ (by setting $U=p^*(X)$ and $V=p(X)$ in Lemma~\ref{lem:conditional-dist-shift}). Upon discarding a small amount of mass, the retained intervals satisfy the desired \titleShort property. We then show that these intervals can be produced by Algorithm~\ref{alg:calibration_population}. Define $A_1$ to be the event that every interval contains sufficiently many calibration samples, and $A_2$ the event that the empirical label average close to the true conditional label probability. Under $A_1\cap A_2$, the empirical estimate is sufficiently far from both endpoints for the algorithm to accept. By Lemma~\ref{lem:sample-size-interval} and Hoeffding's inequality, $A_1\cap A_2$ occurs with probability at least $1-\delta$. We note that Eq.~\eqref{eq:width_function} serves as an upper bound on the widths, while obtaining adaptive bounds that exploit additional model properties or parametric structure is an open direction for future work. For the full proof see
Appendix~\ref{proof:thm:consistency-prediction}.
\end{proofsketch}




% \begin{figure}[ht]
%   \centering
%   \includegraphics[width=0.72\linewidth]{Figures/thm52_width_mean_boxplot.pdf}
%   \caption{Demonstration of width shrinkage. Implementation details see~\ref{app:verification-consistency-prediction-setup}}
%   \label{fig:thm52_box_rate}
% \end{figure}

\begin{figure}[htbp]
    \centering
    % First Subfigure
    \begin{subfigure}[b]{0.52\textwidth}
        \centering
        \includegraphics[width=\textwidth]{Figures/thm52_width_mean_boxplot.pdf} 
  
        \label{fig:thm52_box_rate_a}
    \end{subfigure}
    \hfill 
    \begin{subfigure}[b]{0.45\textwidth}
        \centering
        \includegraphics[width=\textwidth]{Figures/thm52_width_vs_rate.pdf} 
        
        \label{fig:thm52_box_rate_2}
    \end{subfigure}
    
    \caption{Demonstration of width shrinkage. For implementation details see Appendix~\ref{app:verification-consistency-prediction-setup}.}
    \label{fig:thm52_box_rate}
\end{figure}


\subsection{Supporting lemmas}

As shown in the proof sketch, Theorem~\ref{thm:consistency-prediction} relies on three supporting lemmas.  We present these lemmas here to illustrate the key arguments needed to establish the theorem.

\begin{lemma}[Sample size in each interval]\label{lem:sample-size-interval}
Let $\nu$ be a $\lambda$-regular distribution over $[0,1]$ and let $[a_\iteratorInterval,b_\iteratorInterval], \iteratorInterval \in [\numInterval]$ are be intervals such that $|b_\iteratorInterval - a_\iteratorInterval| \geq c$ for all $\iteratorInterval \in [\numInterval]$. For any $n \in Z_+$, let $z_1,\dots,z_n$ be i.i.d. samples from $\nu$, then with probability at least 
% hoeffding's: $1-\numInterval\exp\Big(-\lambda^2c^2n/2\Big)$
$1-J\exp\Big(-\lambda cn/8\Big)$
, we have that
\begin{align*}
    \sumin \mathbbm{1}(z_i \in [a_j,b_j]) \geq \lambda cn/2
\end{align*}
holds for all $j \in [J]$.
\end{lemma}

\begin{proofsketch}
Write the sample count in each interval as a sum of Bernoulli variables and lower-bound their expectations using the $\lambda$-regularity condition. The result follows from a lower-tail Chernoff bound. The full proof can be found in 
Appendix~\ref{proof_lem:sample-size-interval}.
\end{proofsketch}










    
\begin{lemma}[Partition of bounded conditional mean]\label{lem:dist-regularity}
For any $\epsilon, \delta > 0$, and any distribution $\nu$ over $[0,1]$, there exist disjoint intervals $[a_i,b_i], i \in [L]$ such that:
\begin{enumerate}
    \item \emph{(Well-centered conditional means)}  
    The conditional mean of $Z$ on $[a_i,b_i]$ lies away from the interval boundaries by at least $\epsilon$ (except for the endpoints $0$ or $1$). That is, for any $i \in [L]$, 
    \begin{align}\label{eq:regular-condition}
        a_i + \epsilon \cdot \mathbbm{1}(a_i \neq 0) \leq \E_{Z \sim \nu}[Z \mid Z \in [a_i,b_i]] \leq b_i - \epsilon \cdot \mathbbm{1}(b_i \neq 1).
    \end{align}
    \item 
    \emph{(Coverage guarantee)}  
    The union of these \titleShorts covers all but a $\delta$-fraction of the probability mass. That is, \begin{align}\label{eq:regular-condition-coverage}
        \PP_{Z \sim \nu}\Big(Z \notin \cup_{i \in [L]} [a_i,b_i]\Big) \leq \delta.
    \end{align}
    % $\PP\Big(Z \notin \cup_{i \in [K]} [a_i,b_i]\Big) \leq \delta$
    \item \emph{(Bounded interval lengths)}  
    Each interval has a bounded length. That is, for any $i \in [L]$,
    \begin{align}\label{eq:regular-condition-length}
        4\epsilon \leq \Big|b_i - a_i\Big| \leq 56 \epsilon \cdot \log(3/\delta).
    \end{align} %  16 \epsilon \cdot \log(3/\delta) + 24 \epsilon.\log(3/\delta)
\end{enumerate}
\end{lemma}

% \begin{proofsketch}
% Partition $[0,1]$ into cells of width $4\epsilon$.  Between neighboring
% cells, draw an edge in the direction indicated by the cell's
% conditional mean or whenever the neighboring cell has comparable mass.
% Every cell has an outgoing edge.  Vertices at which a right-directed
% and a left-directed edge meet serve as anchors, and the directed chains
% around successive anchors define disjoint intervals.  Truncating each
% chain after a logarithmic number of cells gives the length bound.

% The two boundary edges of each interval imply the interior-mean
% condition: either the boundary-cell mean points inward, or the adjacent
% cell has enough mass to pull the aggregate mean inward.  In each omitted
% gap, failure of the mass-comparison condition forces cell masses to
% decrease geometrically away from a retained interval; summing these
% geometric bounds yields the coverage claim.

% This lemma does not establish the PICPI property.  For $Z=p(X)$, it
% places the conditional mean score in the interval with margin;
% Lemma~\ref{lem:conditional-dist-shift} is then used to transfer that
% margin to the conditional outcome probability. 
% \end{proofsketch}



\begin{proof}[Proof Sketch]
We provide a constructive proof. Discretize $[0,1]$ into $K = \lfloor1/(4\epsilon)\rfloor$ subintervals of length $4\epsilon$ and form a directed graph $G=(V,E)$ on $V = [K]$. We draw a directed edge $i \to j$ ($|i-j|=1$) if either the conditional mean of $Z$ in bin $i$ $\E[Z \mid Z \in B_i]$ is closer to the neighboring endpoint on the side of bin $j$, or the probability mass of $Z$ in bin $i$ is not substantially greater than that of bin $j$ (i.e., $\PP(Z\in B_i)\le 3\PP(Z\in B_j)$). By construction, every vertex has at least one outgoing edge, and there exists  bidirectional edges. 
Group consecutive bidirectional edges into maximal blocks. Partition each block into groups of two or three vertices. Notice that the edges between blocks split into a leftward chain followed by a rightward chain. Attach them to the first and last groups of the neighboring blocks, truncated at a radius of $\lceil\log(3/\delta)\rceil$. This would give the desired intervals.
We then show these intervals satisfy the stated properties. 
The boundedness of widths follows from truncation of the chain. 
The conditional mean property comes from the edge construction and total expectation. Finally, outside these intervals, the absence of reverse edges implies exponential mass decay, bounding total uncovered probability mass by $\delta$. For the full proof see
Appendix~\ref{proof_lem:dist-regularity}.
\end{proof}



\begin{lemma}[Bounds under conditional distribution shift]\label{lem:conditional-dist-shift}
Let $U,V$ be random variables jointly over $[0,1]$, such that $\E[|U-V|] \leq \epsilon > 0$. Then for any $\delta > 0$, integer $L$ and intervals $[a_i,b_i], i = 1,\dots, L$ with disjoint interiors such that $\PP\Big(V \in \cup_{j \in [L]} [a_j,b_j]\Big) > 1- \delta$, 
there exists a subset $\{i_1,\dots,i_\numInterval\}$ of $L$ such that
\begin{enumerate}
    \item \emph{(Conditional distribution on selected intervals)} Conditioning on a selected interval, the discrepancy of the expected value between $U$ and $V$ is controlled by $2\epsilon/\delta$.
 That is, for any $\iteratorInterval \in [\numInterval]$, 
    \begin{align}\label{eq:distribution-shift}
        \Big|\E\Big[U\mid V \in [a_{i_\iteratorInterval},b_{i_\iteratorInterval}]\Big] - \E\Big[V \mid V \in [a_{i_\iteratorInterval},b_{i_\iteratorInterval}]\Big]\Big| \leq 2\epsilon/\delta.
    \end{align}
    \item \emph{(Coverage guarantee)}
    The intervals not selected carry at most a $\delta$-fraction of the probability mass. That is,
\begin{align}\label{eq:distribution-shift-coverage}
        \PP\Big(V \in \cup_{\iteratorInterval \in [L] \setminus \{i_1,\dots,i_\numInterval\}} [a_\iteratorInterval,b_\iteratorInterval]\Big) \leq \delta.
    \end{align}
    Consequently, the union of the selected intervals has probability mass at least $1-2\delta$.
\end{enumerate}
\end{lemma}




\begin{proofsketch}
By applying the triangle inequality to each interval $w$ violating Eq.~\eqref{eq:distribution-shift}, the conditional mean absolute deviation satisfies $\mathbb{E}[\vert{}U - V\vert{} \mid V \in [a_w, b_w]] > 2\epsilon/\delta$. Accounting for endpoint overlaps ($\sum_w \mathbbm{1}(V \in [a_w, b_w]) \leq 2$), we have $2\epsilon \ge \sum_w \mathbb{E}[|U-V| \mathbbm{1}(V \in [a_w, b_w])] > \mathbb{P}(V \in \bigcup_w [a_w, b_w]) \cdot (2\epsilon/\delta)$, yielding total violating mass strictly less than $\delta$. The full proof is provided in Appendix~\ref{proof_lem:conditional-dist-shift}.
\end{proofsketch}

\begin{remark}[Extension to regression] 
\label{rmk:extension_regression}
If $p$ is a regression predictor intended to approximate the conditional mean
$
%m(x) :=
\mathbb{E}[Y\mid X=x],
$ Definition~\ref{def:picpi} becomes 
\[
\mathbb{E}\!\left[Y\mid p(X)\in I_j\right]\in I_j.
\]
%A collection of intervals $\mathcal{I}_p=\{I_j\}_{j=1}^J$, with $I_j\subseteq[a,b]$, is called a collection of PICPIs with respect to $p$ if, for every $j\in[J]$ such that $\mathbb{P}(p(X)\in I_j)>0$, \[ \mathbb{E}\!\left[Y\mid p(X)\in I_j\right]\in I_j. \]
%Each interval simultaneously specifies a subset of predicted conditional means and certifies that the true average outcome over that subset lies in the same interval.
A binary outcome $Y$ can be viewed as a special case where
\[
\mathbb{E}\!\left[Y\mid p(X)\in I_j\right]
=
\mathbb{P}\!\left(Y=1\mid p(X)\in I_j\right).
\]
For bounded outcomes where $Y\in [a,b]$ almost surely, the constructions in Section~\ref{sec:def-alg} still apply by replacing the empirical positive-outcome frequency with the empirical average of $Y$.
%Hoeffding-based
The finite-sample argument in \Cref{thm:consistency-prediction} holds via an affine rescaling of $Y$ and $p(X)$ from $[0,1]$ to $[a,b]$. For unbounded outcomes, the definition remains valid whenever the relevant
conditional expectations exist, but the distribution-free concentration
guarantees developed here do not apply without additional tail or moment assumptions. 
\end{remark}



\iffalse
\begin{claim}[Nearly-regularity]
\label{cla:kl-regularity}
Let $\mu_B$ denote Lebesgue measure on $[0,1]$. Let $X^*$ and $X$ be random variables supported on $[0,1]$ with laws $\nu^*$ and $\nu$, respectively. Assume:
\begin{itemize}
  \item $X^*$ is $\lambda$-regular, i.e. $\nu^*(A) \ge \lambda \mu_B(A)$ for every measurable $A\subseteq[0,1]$.
  \item $\nu^*\ll \nu$ and $D_{\mathrm{KL}}(\nu^*\|\nu)\le \varepsilon$.
\end{itemize}
Then for every $t\in(0,1)$ there exists a measurable set $B_t\subseteq[0,1]$ such that:
\begin{itemize}
  \item $\mu_B(B_t)\le \dfrac{\varepsilon}{\lambda \log(1/t)}$,
  \item for every measurable $A\subseteq [0,1]\setminus B_t$, we have $\nu(A)\ge t\lambda\,\mu_B(A)$.
\end{itemize}
Equivalently, $X$ is $(t\lambda)$-regular on the complement of $B_t$, where the exceptional set has Lebesgue measure at most $\varepsilon/(\lambda\log(1/t))$.

In particular, with $t=\tfrac12$, there exists $B\subseteq[0,1]$ with $\mu_B(B)\le \varepsilon/(\lambda\log 2)$ such that $\nu(A)\ge (\lambda/2)\mu_B(A)$ for all measurable $A\subseteq[0,1]\setminus B$.
\end{claim}

\begin{proof}
Since $X^*$ is $\lambda$-regular, $\nu^*$ dominates $\mu_B$ with constant $\lambda$. In particular, $\nu^*\ll \mu_B$ and the Radon--Nikodym derivative $f^* := \frac{d\nu^*}{d\mu_B}$ exists and satisfies $f^*(x)\ge \lambda$ for $\mu_B$-a.e.\ $x$.

Because $\nu^*\ll \nu$, the Radon--Nikodym derivative $g := \frac{d\nu^*}{d\nu}$ exists $\nu$-a.e.\ and the reverse KL divergence is
\begin{align*}
D_{\mathrm{KL}}(\nu^*\|\nu)
&~:= \int_{[0,1]} \log\!\left(\frac{d\nu^*}{d\nu}\right)\, d\nu^*
 = \int_{[0,1]} \log(g)\, d\nu^*.
\end{align*}
Also $\nu\ll \mu_B$ on the support of $\nu^*$ (and hence on $[0,1]$ after redefining on a null set), so we may write $f:=\frac{d\nu}{d\mu_B}$. Then wherever both densities exist and $f(x)>0$,
\begin{align*}
g(x)=\frac{d\nu^*/d\mu_B}{d\nu/d\mu_B}=\frac{f^*(x)}{f(x)}.
\end{align*}

Fix $t\in(0,1)$ and define the exceptional set
\begin{align*}
B_t := \{x\in[0,1]: f(x) < t\lambda\}.
\end{align*}
On $B_t$, we have $f^*(x)\ge \lambda$ and $f(x)<t\lambda$, hence
\begin{align*}
\frac{f^*(x)}{f(x)} \ge \frac{\lambda}{t\lambda}=\frac1t,
\qquad\text{so}\qquad
\log\!\left(\frac{f^*(x)}{f(x)}\right) \ge \log\!\left(\frac1t\right).
\end{align*}
Using this lower bound inside the reverse KL integral yields
\begin{align*}
\varepsilon
\;\ge\; D_{\mathrm{KL}}(\nu^*\|\nu)
&~= \int_{[0,1]} f^*(x)\log\!\left(\frac{f^*(x)}{f(x)}\right)\, d\mu_B(x)\\
&~\ge \int_{B_t} f^*(x)\log\!\left(\frac{f^*(x)}{f(x)}\right)\, d\mu_B(x)\\
&~\ge \int_{B_t} f^*(x)\log\!\left(\frac1t\right)\, d\mu_B(x)\\
&~= \log\!\left(\frac1t\right)\int_{B_t} f^*(x)\, d\mu_B(x)\\
&~= \log\!\left(\frac1t\right)\,\nu^*(B_t)\\
&~\ge \log\!\left(\frac1t\right)\,\lambda\,\mu_B(B_t),
\end{align*}
where the last step uses $\lambda$-regularity of $\nu^*$ applied to the set $B_t$. Rearranging gives
\begin{align*}
\mu_B(B_t)\le \frac{\varepsilon}{\lambda\log(1/t)}.
\end{align*}

Finally, for any measurable $A\subseteq [0,1]\setminus B_t$, we have $f(x)\ge t\lambda$ for $\mu_B$-a.e.\ $x\in A$, hence
\begin{align*}
\nu(A)
&~= \int_A f(x)\, d\mu_B(x)
\;\ge\; \int_A t\lambda\, d\mu_B(x)
~= t\lambda\,\mu_B(A).
\end{align*}
This proves that $\nu$ is $(t\lambda)$-regular on $[0,1]\setminus B_t$ with the stated bound on $\mu_B(B_t)$.
\end{proof}
\fi

\section{Inference via \titleShorts}\label{sec:tasks}

We now show how to use \titleShorts as building blocks for prediction sets and inference in two tasks: interval prediction and label set prediction.


\subsection{Interval prediction for probabilities}\label{sec:task_1}

In this section we show how \titleShorts can be used to construct intervals for probabilistic prediction with coverage guarantees.
%in a sense that we make formal. 
Assume that we are given a prediction model $p$ and \titleShorts $\intervalset_p$.
Define functions $l,u: [0,1] \to [0,1]$ such that the interval $[l(z),u(z)]$ represents a \titleShort that contains $z$. 
In other words, $[l(z),u(z)]\in\intervalset_p$ and $z \in [l(z) ,u(z)]$.
If there is no such interval, then let $[l(z),u(z)] = [0,1]$. 
Define the expansion operator $E^\epsilon$ as follows: for any set $A \subset [0,1]$, $E^\epsilon(A) = \{b \in [0,1]: d(b,A) \leq \epsilon\}$ where $d$ is the Euclidean distance. 
For $x\in\X$, we obtain a prediction sets based on \titleShorts $\intervalset_p$ by:
\begin{align}\label{eq:task-predicting_true_prob}
    C_{\tp}(x;\epsilon) = E^\epsilon\Big(\Big[l(p(x)),u(p(x))\Big]\Big).
\end{align}
This can be viewed as an $\epsilon$-expanded version of a \titleShort evaluated at $p(x)$.

The following theorem establishes on how well $C_{\tp}(x;\epsilon)$ captures the true conditional probability $\PP(Y=1\mid X=x) $. Specifically, the coverage error is upper-bounded by the fluctuation of true probabilities within each interval and the uncovered prediction mass. By Lemma~\ref{lem:dist-regularity} the second term is at most $\delta$.


% \begin{theorem}[Marginal coverage of the true probability] \label{thm:tp}
% Let $p^*(x) = \PP(Y=1\mid X= x)$ denote the true probability that $Y=1$ given $x$. Assume we have disjoint \titleShorts $\{\interval_\iteratorInterval\}_{\iteratorInterval=1}^{\numInterval}$. Let $\J$ be a random variable following distribution $\PP(\J = j) = \PP(p(X) \in \interval_\iteratorInterval)$. For any $\epsilon > 0$, it holds that 
% \begin{align*} % w/ underbrace
%      \PP\Big(p^*(X) \notin C_{\tp}(X;\epsilon)\Big) \leq 
%     % \epsilon^{-2} \cdot \E\underbrace{\Big[\Var\Big(p^*(X) \mid p(X) \in \interval_\J\Big)\Big]}_{\text{``in-bin" variation}}. 
%      \epsilon^{-2} \cdot \E\underbrace{\Big[\Var\Big(p^*(X) \mid p(X) \in \interval_\J\Big)\Big]}_{\text{``in-bin" variation}} + \underbrace{\Big(1-\sum_{\iteratorInterval=1}^\numInterval\PP(p(X) \in \interval_\iteratorInterval)\Big)}_{\text{predictions not covered by } \intervalset_p}.
% \end{align*}
% \end{theorem}


% The following theorem establishes how well $C_{\tp}(x;\epsilon)$ captures
% the true conditional probability $\PP(Y=1\mid X=x)$. Specifically, the
% coverage error is upper-bounded by the fluctuation of true probabilities
% within each interval, weighted by the probability that the prediction
% falls in that interval. For predictions outside the union of the
% \titleShorts, the reported interval is $[0,1]$.

\begin{theorem}[Marginal coverage of the true probability]\label{thm:tp}
Let $p^*(x)=\PP(Y=1\mid X=x)$ denote the true probability that $Y=1$
given $x$. Assume we have disjoint \titleShorts
$\{\interval_j\}_{j=1}^{\numInterval}$, and write
$m_j=\PP(p(X)\in\interval_j)$. For any $\epsilon>0$, it holds that
\begin{align*}
  \PP\Big(p^*(X)\notin C_{\tp}(X;\epsilon)\Big) \leq\epsilon^{-2}
 \underbrace{\sum_{j:m_j>0}m_j
 \Var\Big(p^*(X)\mid p(X)\in\interval_j\Big)}_{\text{``in-bin'' variation}}
 +\underbrace{\Big(1-\sum_{j=1}^{\numInterval}m_j\Big)}_{
 \text{predictions not covered by }\intervalset_p}.
\end{align*}
\end{theorem}

The proof can be found in Appendix~\ref{proof_thm:tp}.
When the predictor and interval collection are constructed from data that is independent of the test observation, the theorem applies conditionally on those data whenever the resulting intervals satisfy its hypotheses.


The bound in Theorem~\ref{thm:tp} can be sharpened when additional structure is available. The following remark shows that if the predictor $p$ approximates the true probability $p^*$ uniformly or satisfies a Lipschitz-type condition, the in-bin variation term admits explicit upper bounds in terms of the interval lengths.

% \begin{remark}[Corollaries under smoothness conditions]
% If $\|p^* - p\|_\infty \leq \epsilon$, then we have
% \begin{align*}
%     \PP\Big(p^*(X) \notin C_{\tp}(X;\epsilon)\Big) \leq &~
%     \epsilon^{-2} \cdot \E\Big[\Var\Big(p^*(X) \mid p(X) \in \interval_\J\Big)\Big] 
%     + \Big(1-\sum_{\iteratorInterval=1}^\numInterval\PP(p(X) \in \interval_\iteratorInterval)\Big)\\
%     \leq &~ \epsilon^{-2} \cdot \E\Big[(|\interval_\J| + \epsilon)^2/4\Big] 
%     +\Big(1-\sum_{\iteratorInterval=1}^\numInterval\PP(p(X) \in \interval_\iteratorInterval)\Big).
% \end{align*}
% If $|p^*(x) - p^*(y)| \leq L \cdot |p(x) - p(y)|$, then we have
% \begin{align*}
%     \PP\Big(p^*(X) \notin C_{\tp}(X;\epsilon)\Big) \leq &~
%     \epsilon^{-2} \cdot \E\Big[\Var\Big(p^*(X) \mid p(X) \in \interval_\J\Big)\Big]
%     +\Big(1-\sum_{\iteratorInterval=1}^\numInterval\PP(p(X) \in \interval_\iteratorInterval)\Big) \\
%     \leq &~ \epsilon^{-2} \cdot \E\Big[L^2|\interval_\J|^2/4\Big]
%     +\Big(1-\sum_{\iteratorInterval=1}^\numInterval\PP(p(X) \in \interval_\iteratorInterval)\Big).
% \end{align*}
% \end{remark}

% The bound in Theorem~\ref{thm:tp} can be sharpened when additional
% structure is available. The following remark shows that if the predictor
% $p$ approximates the true probability $p^*$ uniformly, or if $p^*$
% satisfies a Lipschitz-type condition with respect to $p$, the in-bin
% variation term admits explicit upper bounds in terms of the interval
% lengths.

\begin{remark}[Corollaries under smoothness conditions]
Under the assumptions of Theorem~\ref{thm:tp}, suppose first that
$\|p^*-p\|_\infty\leq\eta$ for some $\eta\geq0$.
For every positive-mass interval $\interval_j=[a_j,b_j]$, conditioning
on $p(X)\in\interval_j$ gives
\begin{align*}
 p^*(X)\in[\max\{0,a_j-\eta\},\min\{1,b_j+\eta\}]
 \qquad\text{almost surely}.
\end{align*}
The length of this range is at most $|\interval_j|+2\eta$.
Since a random variable supported on an interval of length $d$ has
variance at most $d^2/4$, Theorem~\ref{thm:tp} yields
\begin{align*}
 \PP\Big(p^*(X)\notin C_{\tp}(X;\epsilon)\Big) &~\leq\epsilon^{-2}\sum_{j:m_j>0}m_j
 \Var\Big(p^*(X)\mid p(X)\in\interval_j\Big)
 +\Big(1-\sum_{j=1}^{\numInterval}m_j\Big)\\
 &~\leq\frac{1}{4\epsilon^2}\sum_{j:m_j>0}
 m_j(|\interval_j|+2\eta)^2
 +\Big(1-\sum_{j=1}^{\numInterval}m_j\Big).
\end{align*}
Here $\epsilon>0$ is the expansion radius in
Eq.~\eqref{eq:task-predicting_true_prob}.

Alternatively, suppose that
$|p^*(x)-p^*(x')|\leq L\cdot|p(x)-p(x')|$ for all $x,x'\in\X$.
When $p(x),p(x')\in\interval_j$, their difference is at most
$|\interval_j|$, so the conditional range of $p^*(X)$ has length
at most $L|\interval_j|$. Applying the same variance bound gives
\begin{align*}
 \PP\Big(p^*(X)\notin C_{\tp}(X;\epsilon)\Big)&~\leq\epsilon^{-2}\sum_{j:m_j>0}m_j
 \Var\Big(p^*(X)\mid p(X)\in\interval_j\Big)
 +\Big(1-\sum_{j=1}^{\numInterval}m_j\Big)\\
 &~\leq\frac{L^2}{4\epsilon^2}\sum_{j:m_j>0}
 m_j|\interval_j|^2
 +\Big(1-\sum_{j=1}^{\numInterval}m_j\Big).
\end{align*}
\end{remark}


\iffalse
\begin{remark}
Consider a set of disjoint intervals $I_1,\dots,I_n$. To estimate the upper bound $\E\Big[\Var\Big(p^*(X) \mid p(X) \in \interval_\J\Big)\Big]$, we may use 
\begin{align*}
    \frac{1}{n}\sum_{i=1}^n \frac{\sum_{k=1}^n \mathbbm{1}(p(X_i) \cap p(X_k)) \cdot p^*(X_k)^2 - \Big(\sum_{i=1}^n \mathbbm{1}(p(X_i) \cap p(X_k)) \cdot p^*(X_k)\Big)^2}{\sum_{j=1}^n \mathbbm{1}(p(X_i) \cap p(X_k))}
\end{align*}
where $\mathbbm{1}(p(X_i) \cap p(X_k))$ represents that $p(X_i), p(X_k)$ lie in the same \titleShort. 
% The intervals $\{\interval_\iteratorInterval\}_{\iteratorInterval=1}^{\numInterval}$ should form a partition of $[0,1]$ and may be computed from another set of data $Y_i$.\Xuelin{we should say that in the definition if we require it.}
\end{remark}
\fi

\subsection{Label set prediction for multi-class classification}\label{sec:task_2}

%\subsubsection{Two-sided certificate}

We now use \titleShorts for label set prediction in classification with $\numClass$ classes. In particular, we establish three guarantees: (i) an exact finite-sample certificate for class-conditional coverage, providing the tightest worst-case miscoverage bound available from the interval information (Theorem~\ref{thm:mc_error}); (ii) population size-optimality, showing that our interval selection minimizes expected label-set size among deterministic unions of the
given classwise intervals satisfying the prescribed class-conditional coverage guarantees (Theorem~\ref{thm:mc_optimality}); and (iii) a finite-sample same-sample guarantee, showing that the empirically feasible procedure retains valid class-conditional coverage when the same calibration data are used to construct the intervals and estimate their masses (Theorem~\ref{thm:mc_same_sample}). 



Let $p: \X \rightarrow [0,1]^\numClass$ be a prediction model that outputs class probabilities, where its $\iteratorClass$-th entry 
$p_\iteratorClass(x)$ predicts ${\PP}(Y(X)=\iteratorClass \mid X=x)$.
For each class $\iteratorClass\in[\numClass]$, we construct disjoint \titleShorts
\[
\intervalset_{p_\iteratorClass}
=
\big\{[a_{\iteratorClass,j},b_{\iteratorClass,j}]\big\}_{j=1}^{J_\iteratorClass}
\]
for the one-vs-rest problem $1\{Y(X)=\iteratorClass\}$ using the interval-construction procedure in Section~\ref{sec:def-alg}.
Suppose we select a subset of interval indices $S_\iteratorClass\subseteq [J_\iteratorClass]$ for class $g$ (by some selection rules of our choice), and denote their collection as $S=(S_1,\ldots,S_\numClass)$.
%For any collection $S=(S_1,\ldots,S_\numClass)$ with subset of interval indices $S_\iteratorClass\subseteq [J_\iteratorClass]$. 
For any given $S$, we define the classwise acceptance region as the union of intervals in $S_\iteratorClass$:
\[
A_\iteratorClass(S_\iteratorClass)
:=
\bigcup_{j\in S_\iteratorClass}[a_{\iteratorClass,j},b_{\iteratorClass,j}].
\]
Given input $x$, we predict a set of labels which collects the classes whose predicted probabilities fall within their respective acceptance regions:
\begin{align}\label{eq:mc_acceptance}
C_{\mc}(x;S)
=
\Bigl\{\iteratorClass\in[\numClass]:
p_\iteratorClass(x)\in A_\iteratorClass(S_\iteratorClass)\Bigr\}.
\end{align}
For example, $S_\iteratorClass$ can be constructed by a simple upper-endpoint threshold rule with threshold $\tau$:
\begin{align}\label{eq:classification_thresholding_rule}
S_\iteratorClass(\mcThreshold)
:=
\Bigl\{j\in[J_\iteratorClass]: b_{\iteratorClass,j}\ge \mcThreshold\Bigr\},
\qquad
S(\mcThreshold)
:=
\bigl(S_1(\mcThreshold),\ldots,S_\numClass(\mcThreshold)\bigr),
\end{align}
so that
\begin{align}\label{eq:task-classification}
    % C_{\mc}(x;\mcThreshold)
    % =
    C_{\mc}\bigl(x;S(\mcThreshold)\bigr)
    =
    \Bigl\{\iteratorClass\in[\numClass]:
    p_\iteratorClass(x)\in \bigcup_{j:\,b_{\iteratorClass,j}\ge \mcThreshold}
    [a_{\iteratorClass,j},b_{\iteratorClass,j}]
    \Bigr\}.
\end{align}



For notational convenience, define
\[
m_{\iteratorClass,j}
:=
\PP\bigl(p_\iteratorClass(X)\in [a_{\iteratorClass,j},b_{\iteratorClass,j}]\bigr),
\qquad
\pi_\iteratorClass
:=
\PP\bigl(Y(X)=\iteratorClass\bigr),
\]
\[
t_\iteratorClass
:=
1-\sum_{j=1}^{J_\iteratorClass} m_{\iteratorClass,j},
\qquad
\intervalMaxLen_{p_\iteratorClass}
:=
\max_{1\le j\le J_\iteratorClass}\bigl(b_{\iteratorClass,j}-a_{\iteratorClass,j}\bigr).
\]
Throughout, we assume $\pi_\iteratorClass>0$ for every class $\iteratorClass\in[\numClass]$.




% two-sided 
The following theorem gives an exact class-conditional coverage certificate for any rule of the form $C_{\mc}(x;S)$.

\begin{theorem}[Class-conditional coverage of the true label]\label{thm:mc_error}
Fix a class $\iteratorClass\in[\numClass]$ and a subset $S_\iteratorClass\subseteq [J_\iteratorClass]$.
It holds that 
\begin{align}
\PP\Bigl(Y(X)\notin C_{\mc}(X;S)\mid Y(X)=\iteratorClass\Bigr)
\le
\Gamma_{\iteratorClass}(S_\iteratorClass),
\end{align}
where
\[
\Gamma_{\iteratorClass}(S_\iteratorClass)
:=
\min\Bigl\{
\Gamma_{\iteratorClass}^{\mathrm{up}}(S_\iteratorClass),
\Gamma_{\iteratorClass}^{\mathrm{low}}(S_\iteratorClass)
\Bigr\},
\]
and 
\[
\Gamma_{\iteratorClass}^{\mathrm{up}}(S_\iteratorClass)
:=
\frac{t_\iteratorClass+\sum_{j\notin S_\iteratorClass} m_{\iteratorClass,j} b_{\iteratorClass,j}}{\pi_\iteratorClass},
\qquad
\Gamma_{\iteratorClass}^{\mathrm{low}}(S_\iteratorClass)
:=
1-\frac{\sum_{j\in S_\iteratorClass} m_{\iteratorClass,j} a_{\iteratorClass,j}}{\pi_\iteratorClass},
\]



Moreover, this bound is exact in the following sense: for fixed
\[
\pi_\iteratorClass,\quad
t_\iteratorClass,\quad
\bigl\{m_{\iteratorClass,j},a_{\iteratorClass,j},b_{\iteratorClass,j}\bigr\}_{j=1}^{J_\iteratorClass}.
\]
the quantity $\Gamma_{\iteratorClass}(S_\iteratorClass)$ is the worst-case class-conditional miscoverage
over all one-vs-rest laws satisfying the \titleShorts definition
\[
\PP\Bigl(Y(X)=\iteratorClass \,\big|\, p_\iteratorClass(X)\in[a_{\iteratorClass,j},b_{\iteratorClass,j}]\Bigr)
\in [a_{\iteratorClass,j},b_{\iteratorClass,j}]
\quad \text{for all }j.
\]
\end{theorem}

For the proof see
Appendix~\ref{proof_thm:mc_error}.

% While the coverage certificate ensures statistical validity, our secondary goal is efficiency: we want the predicted label set to be as small as possible. Building on the \titleShorts, we now derive an interval selection rule that achieves size-optimality by minimizing the expected number of predicted labels.


\paragraph{Population-optimal label set.}
Ideally, we want our predicted label set to be as small as possible while still guaranteeing that the true label is caught within our error budget. To formalize this trade-off, let $\alpha = (\alpha_1, \ldots, \alpha_\numClass) \in [0,1]^\numClass$ denote a vector of target class-conditional miscoverage levels specified by the user.
For each class $\iteratorClass \in [\numClass]$, we find the optimal subset of prediction intervals to keep by minimizing the expected size of the prediction set, subject to the error guarantee. Formally, we define the population-optimal retained interval set as:
\begin{align}
S_\iteratorClass^\star(\alpha_\iteratorClass)
\in
\arg\min_{S\subseteq[J_\iteratorClass]}
\sum_{j\in S} m_{\iteratorClass,j}
\quad \text{subject to} \quad
\Gamma_{\iteratorClass}(S)\le \alpha_\iteratorClass.
\label{eq:mc_population_opt}
\end{align}

The resulting population-optimal label set simply collects all classes whose predicted probabilities fall into these optimally selected intervals:
%The corresponding population-optimal label set is
\[
C_{\mc}^\star(x;\alpha)
:=
C_{\mc}\bigl(x;S^\star(\alpha)\bigr),
\qquad
S^\star(\alpha)
:=
\bigl(
S_1^\star(\alpha_1),\ldots,S_\numClass^\star(\alpha_\numClass)
\bigr).
\]
Because
\[
\Gamma_{\iteratorClass}(S)
=
\min\Bigl\{
\Gamma_{\iteratorClass}^{\mathrm{up}}(S),
\Gamma_{\iteratorClass}^{\mathrm{low}}(S)
\Bigr\},
\]
the feasible set in \eqref{eq:mc_population_opt} is the union of two simpler feasible sets.
Equivalently, one may solve
\begin{align}
S_\iteratorClass^{\mathrm{keep}}(\alpha_\iteratorClass)
\in
\arg\min_{S\subseteq[J_\iteratorClass]}
&\sum_{j\in S} m_{\iteratorClass,j}
\label{eq:mc_keep_problem}\\
\text{subject to}\quad
&\sum_{j\in S} a_{\iteratorClass,j}m_{\iteratorClass,j}
\ge
(1-\alpha_\iteratorClass)\pi_\iteratorClass,
\nonumber
\end{align}
and
\begin{align}
S_\iteratorClass^{\mathrm{drop}}(\alpha_\iteratorClass)
\in
\arg\min_{S\subseteq[J_\iteratorClass]}
&\sum_{j\in S} m_{\iteratorClass,j}
\label{eq:mc_drop_problem}\\
\text{subject to}\quad
&t_\iteratorClass+\sum_{j\notin S} m_{\iteratorClass,j} b_{\iteratorClass,j}
\le
\alpha_\iteratorClass \pi_\iteratorClass,
\nonumber
\end{align}
and then keep whichever feasible solution has the smaller objective value.

\begin{algorithm}[t] 
\caption{Constructing population-optimal label set} %two-sided 
\label{alg:mc_population_optimal}
\begin{algorithmic}[1]
\Require Disjoint \titleShorts $\intervalset_{p_\iteratorClass}=\{[a_{\iteratorClass,j},b_{\iteratorClass,j}]\}_{j=1}^{J_\iteratorClass}$, population interval masses $\{m_{\iteratorClass,j}\}$, class priors $\{\pi_\iteratorClass\}$, and targets $\alpha=(\alpha_1,\ldots,\alpha_\numClass)$
\Ensure Population-optimal label set $C_{\mc}^\star(\cdot;\alpha)$
\For{$\iteratorClass=1$ to $\numClass$}
    \State Compute $t_\iteratorClass = 1-\sum_{j=1}^{J_\iteratorClass} m_{\iteratorClass,j}$.
    \State Solve the keep-side problem \eqref{eq:mc_keep_problem}; if infeasible, assign objective value $+\infty$.
    \State Solve the drop-side problem \eqref{eq:mc_drop_problem}; if infeasible, assign objective value $+\infty$.
    \State Let $S_\iteratorClass^\star(\alpha_\iteratorClass)$ be the feasible solution with smaller objective value.
\EndFor
\State \Return $C_{\mc}^\star(x;\alpha)=\Bigl\{\iteratorClass\in[\numClass]: p_\iteratorClass(x)\in \bigcup_{j\in S_\iteratorClass^\star(\alpha_\iteratorClass)} [a_{\iteratorClass,j},b_{\iteratorClass,j}]\Bigr\}$.
\end{algorithmic}
\end{algorithm}

\begin{theorem}[Size-optimality at fixed guaranteed coverage]\label{thm:mc_optimality}
Let $C_{\mc}^\star(x;\alpha)$ be defined by \eqref{eq:mc_population_opt}.
Then, for every class $\iteratorClass\in[\numClass]$,
\[
\PP\Bigl(Y(X)\notin C_{\mc}^\star(X;\alpha)\mid Y(X)=\iteratorClass\Bigr)
\le
\alpha_\iteratorClass.
\]
Moreover,
\[
\E\Bigl[\bigl|C_{\mc}^\star(X;\alpha)\bigr|\Bigr]
=
\min\Bigl\{
\E\bigl[|C_{\mc}(X;S)|\bigr]:
\Gamma_{\iteratorClass}(S_\iteratorClass)\le \alpha_\iteratorClass
\text{ for all }\iteratorClass\in[\numClass]
\Bigr\}.
\]
\end{theorem}


For the proof see
Appendix~\ref{proof_thm:mc_optimality}.

\paragraph{Empirically feasible version of the algorithm.}
The population-optimal algorithm presented above is not directly implementable because the interval masses $m_{\iteratorClass,j}$ and the class priors $\pi_\iteratorClass$ are unknown.
To obtain a practical algorithm we thus replace them by one-sided confidence bounds computed from the same calibration set that is already used to construct the \titleShorts.
Let
\[
D_{\mathrm{cal}}=\bigl\{(X_k,Y_k)\bigr\}_{k=1}^{n_{\mathrm{cal}}}
\]
be the calibration set, and let
\[
\widehat{\intervalset}_{p_\iteratorClass}
=
\bigl\{
\widehat I_{\iteratorClass,j}
=
[\widehat a_{\iteratorClass,j},\widehat b_{\iteratorClass,j}]
\bigr\}_{j=1}^{\widehat J_\iteratorClass}
\]
be the resulting disjoint \titleShorts for class $\iteratorClass$.
Let $\mathcal{A}_{p_\iteratorClass}$ denote the finite candidate interval family searched by Algorithm~\ref{alg:mc_population_optimal} for class $\iteratorClass$, and define
\[
N_{\mathrm{cand}}
:=
\sum_{\iteratorClass=1}^{\numClass}
\bigl|\mathcal{A}_{p_\iteratorClass}\bigr|.
\]
For each selected interval, define the empirical mass
\[
\widehat m_{\iteratorClass,j}
:=
\frac{1}{n_{\mathrm{cal}}}
\sum_{k=1}^{n_{\mathrm{cal}}}
\mathbbm{1}\Bigl(
p_\iteratorClass(X_k)\in \widehat I_{\iteratorClass,j}
\Bigr),
\qquad
\widehat\pi_\iteratorClass
:=
\frac{1}{n_{\mathrm{cal}}}
\sum_{k=1}^{n_{\mathrm{cal}}}
\mathbbm{1}(Y_k=\iteratorClass).
\]
For confidence levels $\delta_m,\delta_\pi\in(0,1)$, let
\[
\varepsilon_m
:=
\sqrt{
\frac{\log(2N_{\mathrm{cand}}/\delta_m)}{2n_{\mathrm{cal}}}
},
\qquad
\varepsilon_\pi
:=
\sqrt{
\frac{\log(2\numClass/\delta_\pi)}{2n_{\mathrm{cal}}}
}.
\]
We then define
\[
\underline m_{\iteratorClass,j}
:=
\bigl(\widehat m_{\iteratorClass,j}-\varepsilon_m\bigr)_+,
\qquad
\underline\pi_\iteratorClass
:=
\bigl(\widehat\pi_\iteratorClass-\varepsilon_\pi\bigr)_+,
\qquad
\overline\pi_\iteratorClass
:=
\bigl(\widehat\pi_\iteratorClass+\varepsilon_\pi\bigr)\wedge 1.
\]
For any subset $S_\iteratorClass\subseteq [\widehat J_\iteratorClass]$, define
\[
\widehat\Gamma_{\iteratorClass}^{\mathrm{up}}(S_\iteratorClass)
:=
\begin{cases}
\displaystyle
\frac{
1-\sum_{j\in S_\iteratorClass}\underline m_{\iteratorClass,j}
-\sum_{j\notin S_\iteratorClass}(1-\widehat b_{\iteratorClass,j})\underline m_{\iteratorClass,j}
}{\underline\pi_\iteratorClass},
& \underline\pi_\iteratorClass>0,\\[3ex]
+\infty,
& \underline\pi_\iteratorClass=0,
\end{cases}
\]
\[
\widehat\Gamma_{\iteratorClass}^{\mathrm{low}}(S_\iteratorClass)
:=
1-
\frac{
\sum_{j\in S_\iteratorClass}\widehat a_{\iteratorClass,j}\underline m_{\iteratorClass,j}
}{\overline\pi_\iteratorClass},
\qquad
\widehat\Gamma_{\iteratorClass}(S_\iteratorClass)
:=
\min\Bigl\{
\widehat\Gamma_{\iteratorClass}^{\mathrm{up}}(S_\iteratorClass),
\widehat\Gamma_{\iteratorClass}^{\mathrm{low}}(S_\iteratorClass)
\Bigr\}.
\]
This yields the empirical relaxation
\[
\widehat S_\iteratorClass^\star(\alpha_\iteratorClass)
\in
\arg\min_{S\subseteq[\widehat J_\iteratorClass]}
\sum_{j\in S}\widehat m_{\iteratorClass,j}
\quad \text{subject to}\quad
\widehat\Gamma_{\iteratorClass}(S)\le \alpha_\iteratorClass,
\]
and the resulting feasible label set
\[
\widehat C_{\mc}^\star(x;\alpha)
:=
\Bigl\{\iteratorClass\in[\numClass]:
p_\iteratorClass(x)\in
\bigcup_{j\in \widehat S_\iteratorClass^\star(\alpha_\iteratorClass)}
\widehat I_{\iteratorClass,j}
\Bigr\}.
\]
The key point is that the same calibration sample may be used both to construct the intervals and to estimate their masses: because Algorithm~\ref{alg:mc_population_optimal} only searches over a fixed finite candidate family, uniform concentration over that family automatically controls the masses of the finally selected intervals as well.

\begin{theorem}[Fully feasible same-sample certificate]\label{thm:mc_same_sample}
Suppose Algorithm~\ref{alg:mc_population_optimal} is run for each class $\iteratorClass\in[\numClass]$ with classwise failure budgets summing to $\delta_{\mathrm{PI}}$, so that with probability at least $1-\delta_{\mathrm{PI}}$ all returned intervals in
\[
\widehat{\intervalset}_{p_\iteratorClass}
=
\bigl\{
\widehat I_{\iteratorClass,j}
=
[\widehat a_{\iteratorClass,j},\widehat b_{\iteratorClass,j}]
\bigr\}_{j=1}^{\widehat J_\iteratorClass},
\qquad
\iteratorClass\in[\numClass],
\]
are genuine disjoint \titleShorts simultaneously.
For any $S=(S_1,\ldots,S_\numClass)$ with $S_\iteratorClass\subseteq [\widehat J_\iteratorClass]$, define
\[
\widehat C_{\mc}(x;S)
=
\Bigl\{
\iteratorClass\in[\numClass]:
p_\iteratorClass(x)\in
\bigcup_{j\in S_\iteratorClass}\widehat I_{\iteratorClass,j}
\Bigr\}.
\]
Then, with probability at least
\[
1-\delta_{\mathrm{PI}}-\delta_m-\delta_\pi
\]
over the draw of $D_{\mathrm{cal}}$, the following holds simultaneously for all classes $\iteratorClass\in[\numClass]$ and all subsets $S_\iteratorClass\subseteq[\widehat J_\iteratorClass]$:
\[
\PP\Bigl(
Y(X)\notin \widehat C_{\mc}(X;S)
\ \Big|\ 
Y(X)=\iteratorClass,\ D_{\mathrm{cal}}
\Bigr)
\le
\widehat\Gamma_{\iteratorClass}(S_\iteratorClass).
\]
In particular, whenever $\widehat S_\iteratorClass^\star(\alpha_\iteratorClass)$ is feasible for every class,
\[
\PP\Bigl(
Y(X)\notin \widehat C_{\mc}^\star(X;\alpha)
\ \Big|\ 
Y(X)=\iteratorClass,\ D_{\mathrm{cal}}
\Bigr)
\le
\alpha_\iteratorClass
\qquad
\text{for all }\iteratorClass\in[\numClass].
\]
\end{theorem}
The proof can be found in 
Appendix~\ref{proof_thm:mc_same_sample}.
\section{Experiments}\label{sec:experiment}


% Calibration error for bin $B_j$ is calculated as 
% \[
% (\text{\{average of $p(x_i)$ when evaluation data $p(x_i)\in B_j$\}} - \text{mid point of }B_j)^2
% \]


We compare \titleShort-based methods with other interval-prediction methods for predicting probabilities (Task 1) and multiclass classification (Task 2). In both tasks, point predictors are fitted on a training split, intervals (or label sets) are constructed using a separate calibration split, and the resulting intervals (or label sets) are evaluated on a test split. 

\subsection{Task 1: Interval prediction for probabilities.}

% \noindent\textbf{Task 1: Interval prediction for probabilities.} 
As described in Section~\ref{sec:task_1}, the goal is to produce an interval to predict $\PP (Y=1\mid X=x)$. Each baseline produces a set of intervals $\intervalset$. That is, for a predictor $\hat{p}$, we obtain the interval $\interval\in\intervalset$ containing predicted $\hat p(x)$ for an input $x$. 


The data-generating processes (DGPs) use 1-dimensional and 20-dimensional covariates with binary targets. Details of the target probabilities and label-generating mechanisms are provided in Appendix~\ref{app:task1_setup}.



\subsection{Task 2: Label set prediction for multi-class classification.}

% \noindent\textbf{Task 2: Label set prediction for multi-class classification.}
As described in Section~\ref{sec:task_2}, the goal is to produce a label set. For each class, interval-based baselines produce an interval set $\intervalset^{(g)}$, from which we select the interval containing $\hat p_g(x)$. 









% The threshold $\tau$ is calibrated on the calibration split using the true-label upper scores
% \[
% S_i = U_{Y_i}(X_i), \qquad i=1,\dots,n_{\mathrm{cal}},
% \]
% and is chosen as their empirical lower $\alpha$-quantile,
% \[
% \tau
% =
% S_{(k)},
% \qquad
% k=\left\lceil (n_{\mathrm{cal}}+1)\alpha \right\rceil,
% \]
% where $S_{(1)} \le \cdots \le S_{(n_{\mathrm{cal}})}$ are the ordered calibration scores.


For DGPs, target labels are drawn as $Y \mid X=x \sim \mathrm{Categorical}(q(x))$. The Linear-Gaussian DGP uses isotropic Gaussian covariates $X \sim \mathcal{N}(0, I_{10})$ with sparse linear softmax probabilities $q(x) = \mathrm{softmax}(Wx + b)$. The Nonlinear-Mixture DGP generates non-Gaussian, multimodal covariates via a four-component latent-factor mixture. Its class probabilities scale the linear logits by a bounded covariate-dependent function $s(x)$ such that $q(x) = \mathrm{softmax}\left((Wx + b)/{s(x)}\right)$. Details are provided in Appendix~\ref{app:task2_setup}.


\subsection{Baselines}
We compare the following baselines: (i) \titleShort constructs intervals from the calibration split using the same underlying procedure as our method, with an $\epsilon$-expansion for Task~1 and classwise thresholding of upper endpoints for Task~2. 
(ii) \textbf{Simultaneous confidence intervals} construct pointwise intervals by applying the Delta method to a Wald confidence region for the parameters of a parametric probability model. 
(iii) \textbf{Split conformal prediction} forms symmetric intervals around the fitted probability using a calibration quantile of absolute residuals, providing marginal coverage guarantees. 
(iv) \textbf{Calibration-based intervals} discretize $[0,1]$ into bins and retain bins whose mass-weighted calibration error is comparable to that of \titleShort, assigning each prediction to the closest retained bin when necessary. 
(v) \textbf{Fixed-width binning} discretizes $[0,1]$ into fixed-width bins and reports the bin containing the predicted probability. (vi) \textbf{Conformal classification}, used only for Task~2, constructs prediction sets by including classes whose predicted probabilities exceed a calibration-based threshold rather. This is not an interval-based method, but has marginal coverage of the true label.
We extend their formulation details and metrics in Appendix~\ref{app:baselines} and~\ref{app:metrics}, respectively.





\subsection{PICPIs with empirical self-consistency}
\label{sec:alg_emp}
\begin{algorithm}[tb] % changed form {algorithm*} (two columns) 
   \caption{ \titleShort Construction (empirical mode with self-consistency on observed data)}
   \label{alg:calibration_empirical} 

  % \small
\begin{algorithmic}[1]
\Procedure{\textsc{Calibration}}{data points $\{(x_i,y_i)\}_{i=1}^\numData,$ predictor $ p$} \Comment{Obtain a set of \titleShorts with respect to the empirical distribution}
\State Compute predictions $p_i\leftarrow p(x_i)$ and sort $p_i$ in ascending order. Let $p_i'$ and $y_i'$ to denote the sorted prediction and outcome respectively.
\State Let $t_0 \leftarrow 0$.

\For {$j = 1,2,\dots,$}
\State Find $\min t_j \in (t_{j-1},1] $ subject to \Comment{find the next endpoint on a grid of [0,1]}
\begin{align*}
    \frac{\sum_{i=1}^n y_i' \mathbbm{1}(p_i' \in (t_{j-1}, t_j])}{\sum_{i=1}^n \mathbbm{1}(p_i' \in (t_{j-1}, t_j])} \in &~ (t_{j-1}, t_j].\\
    \frac{\sum_{i=1}^n y_i' \mathbbm{1}(p_i' \in (t_{j}, 1])}{\sum_{i=1}^n \mathbbm{1}(p_i' \in (t_{j}, 1])} \in &~ (t_{j}, 1].
\end{align*}
\State Let $\interval_j \leftarrow (t_{j-1}, t_j]$.
\If {$t_j = 1$}
\State Let $\numInterval \leftarrow j$. {\bf Break}.
\EndIf
\EndFor
\State {\bf Return} $\{\interval_j\}_{j=1}^\numInterval$
\EndProcedure
\\
\Procedure{\textsc{Inference}}{covariate input $x,$ predictor $ p$} 
\Comment{Find the \titleShort that contains $p(x)$}
\State $\{\interval_j\}_{j=1}^\numInterval \leftarrow \textsc{Calibration}\left(\{(x_i,y_i)\}_{i=1}^n, p\right)$
\State Let $\interval$ be the interval among $\{\interval_j\}_{j=1}^\numInterval$ that contains $p(x)$
\State {\bf Return} $\interval$
\EndProcedure
\end{algorithmic}
\end{algorithm}




We now present Algorithm~\ref{alg:calibration_empirical}, the empirical variant of algorithm used in experiments in this section. It returns \titleShorts with respect to the empirical distribution of the calibration sample. We refer to this property as \emph{empirical self-consistency} to distinguish it from the population validity guaranteed by Theorem~\ref{thm:algorithm_population}. 
% Algorithm~\ref{alg:calibration_empirical} gives a empirical
% alternative of Algorithm~\ref{alg:calibration_population}.
The key difference is that Algorithm~\ref{alg:calibration_empirical} omits the concentration margin (Line~\ref{line:empirical-interval}, Algorithm~\ref{alg:calibration_population}) when evaluating intervals.  
While this omission yields narrower, more practical intervals, the theoretical population validity guaranteed by Theorem~\ref{thm:algorithm_population} no longer applies.
Figure~\ref{fig:emp_mode_diagnostics} illustrates the resulting tradeoff. %\Xuelin{changed here to be more clear on validity} 
%MJ  What does "yet by large still valid" mean?.  Need to be clearer in this paragraph.
Most proposed intervals also satisfy the \titleShort condition on held-out data, whereas the shortest intervals (of length at most $0.03$) are more prone to false acceptance.

The algorithm proceeds as follows: starting at the left
endpoint, it chooses the smallest next endpoint for which the \emph{empirical} label mean of the new bin lies in that bin, and so does the remaining part in $[0,1]$. It then repeats this step on the remainder of $[0,1]$. At inference time, a new
prediction is assigned to its unique bin in the resulting partition.



%old
%Algorithm~\ref{alg:calibration_empirical} serves as an empirical approximation to Algorithm~\ref{alg:calibration_population}. The key difference is that Algorithm~\ref{alg:calibration_empirical} omits the concentration margin in Line~\ref{line:empirical-interval} when evaluating intervals. \Xuelin{changed:} We observe in practice that the empirical version yields narrower intervals than the population version, while a large fraction still satisfies the population-level \titleShort property.
%This is because the concentration margin constitutes a uniform high-probability correction to ensure population \titleShort validity, though at a cost of lower resolution. The empirical mean may exhibit sufficient concentration despite the absence of explicit finite-sample theoretical guarantees. However, the empirical version may produce false-positive intervals that are overly narrow (length $\le$ 0.03). We demonstrate this in Figure~\ref{fig:emp_mode_diagnostics}. 







%We recommend considering this tradeoff when choosing between two alternatives for down-stream applications.

%By omitting this correction, the empirical version admits a narrower candidate intervals and still are valid when the data distribution is not adversarial.
%This is because the worst-case correction in population version via concentration margin may lead to the exclusion of intervals that are in fact well-calibrated. 

%Those narrower intervals, despite omitting the margin, may still be valid. 
% By omitting this margin, the empirical version admits a richer set of candidate intervals, which can better adapt to the observed data and yield tighter or more informative intervals.
%This could reflect a bias-variance trade-off: 


% two separate figures:

% \begin{figure}[htbp]
%     \centering
%     \includegraphics[width=\textwidth]{Figures/heldout_picpi_pass_rate_by_interval_length.pdf} 
%     \caption{Demonstration of empirical version (Algorithm~\ref{alg:calibration_empirical}). \textbf{Left:} more calibration data leads to narrower intervals. \textbf{Right:} the proportion of proposed intervals satisfying the \titleShort property on a held-out dataset. While most intervals satisfy \titleShort property, yet the narrowest intervals (length$<$ 0.03) at the bottom of each panel risk lower \titleShort property pass rate. 50 repetitions. Setup is the same as in \ref{app:task1_setup} besides changing calibration set sizes. }
%     \label{fig:emp_mode_diagnostics}
% \end{figure}


% \begin{figure}[htbp]
%     \centering
%     \includegraphics[width=.4\textwidth]{Figures/heldout_picpi_failure_rate_by_method.pdf} 
%     \caption{Demonstration of empirical mode and population mode \titleShort property pass rate. 50 repetitions. Setup is the same as in \ref{app:task1_setup} besides changing calibration set size. }
%     \label{fig:emp_mode_diagnostics_rate}
% \end{figure}


\begin{figure}[tb]
    \centering
    % First Subfigure
    \begin{subfigure}[b]{0.36\textwidth}
        \centering
        \includegraphics[width=\textwidth]{Figures/heldout_picpi_failure_rate_by_method.pdf} 
  
        
    \end{subfigure}
    \hfill 
    \begin{subfigure}[b]{0.6\textwidth}
        \centering
        \includegraphics[width=\textwidth]{Figures/heldout_picpi_pass_rate_by_interval_length.pdf} 
        
        
    \end{subfigure}
    
    \caption{Demonstration of empirical version (Algorithm~\ref{alg:calibration_empirical}). \textbf{Left:} Proportion of proposed intervals satisfying the \titleShort property on a held-out dataset. \textbf{Middle:} With more calibration data, the empirical version admits narrower intervals. \textbf{Right:} While most intervals satisfy the \titleShort property, the narrowest intervals (length $\leq 0.03$) at the bottom exhibit lower pass rates. Their pass rate follows a U-shaped pattern as the amount of calibration data increases: initially, additional data allow the procedure to admit increasingly narrow and more challenging intervals, leading to a lower pass rate; with sufficiently large calibration samples, however, the increased information improves calibration even for these narrow intervals, and the pass rate recovers.
     For a formal presentation of the setup studied here, see Appendix~\ref{app:task1_setup_emp_mode_diagnostics}. }
    \label{fig:emp_mode_diagnostics}
\end{figure}





\begin{figure}[!ht]
\centering
\includegraphics[width=\linewidth]{artifacts/task1_interval_visualization.pdf}
\caption{Visualization under mildly miscalibrated model. %mild model misspecification. 
\titleShort adapts its interval endpoints to local prediction bins to account for deviations between the fitted and target probabilities, while maintaining substantially shorter intervals than split conformal prediction. The simultaneous confidence interval is less responsive to bin-specific deviations from calibration, while the binning-based baselines have poor coverage of the target probability. 
%Detailed set up is deferred to Section~\ref{sec:experiment}.
}
\label{fig:task1_interval_visualization} % univariate
\end{figure}



\begin{figure}[ht]
\centering
\includegraphics[width=\linewidth]{artifacts/task1_misspecified_point_estimator_decision_tree.pdf}
\caption{Task~1 demonstration with a misspecified decision-tree point estimator. \titleShorts adapt to biased, piecewise-constant predictions by adjusting interval endpoints and widening where calibration reveals instability or mismatch. The smooth curve shows the true conditional probability $p^*(x)$, while the step function shows the misspecified decision-tree estimate $\hat p(x)$. The dashed step curves mark the lower and upper endpoints of the \titleShort interval, and the shaded boxes show the resulting interval assigned at each $x$.}
\label{fig:task1_misspecified_tree}
\end{figure}


\subsection{Task 1 Results}




Figures~\ref{fig:task1_interval_visualization} and~\ref{fig:task1_misspecified_tree}, together with Table~\ref{tab:task1_multivariate_mc_summary}, show that \titleShort achieves a favorable balance among length, calibration, and coverage measures. In the univariate visualization, \titleShort adapts to the local score partition and tracks the nonlinear target probability $p^*(x)$ more faithfully than the two binning baselines, while split conformal is visibly conservative. 

In the misspecified decision-tree example of Figure~\ref{fig:task1_misspecified_tree}, the \titleShorts remain responsive when the fitted predictor is structurally biased: the decision tree produces a piecewise-constant \(\hat p(x)\) that fails to track the nonlinear true probability curve, but \titleShort adapts its interval endpoints around these score plateaus and widens in regions where the calibration data reveal instability or mismatch.

Quantitatively, Table~\ref{tab:task1_multivariate_mc_summary} shows that \titleShort attains the lowest calibration errors while maintaining substantially shorter intervals than the simultaneous confidence interval and split conformal prediction, and it achieves much higher true-probability coverage than the calibration-based and fixed-width binning baselines. It avoids the over-conservatism of conformal-style wide intervals while retaining far better coverage and calibration behavior than binning alternatives. 
In particular,
its calibration errors indicate that the reported intervals are well centered around the true probabilities and the its binning follows the calibration structure of predictor. 

% Thus \titleShort is the only method in this comparison that is simultaneously well calibrated, covers the  prediction space, and retains substantial predictive coverage.







\begin{table}[ht]
\centering
\scriptsize

\begin{tabular}{lccccc}

\toprule
Method & Avg. Length & Rel. ECE$_\text{mid}$ &  Rel. ECE$_\text{pred}$  & Coverage over [0,1] & Coverage of $p^*$ \\
\midrule
Simultaneous CI & 0.41 $\pm$ 0.01 & 6.39 & 3.25 & -- & 100.00\% $\pm$ 0.00\% \\
PICPI & 0.11 $\pm$ 0.01 & 1.00 & 1.00 & 100.00\% $\pm$ 0.00\% & 79.47\% $\pm$ 4.68\% \\
Split conformal prediction & 0.87 $\pm$ 0.01 & 27.02 & 2.43 & -- & 100.00\% $\pm$ 0.00\% \\
Calibration-based interval & 0.02 $\pm$ 0.00 & 1.16 & 1.01 & 60.78\% $\pm$ 11.49\% & 28.07\% $\pm$ 3.25\% \\
Fixed-width binning & 0.02 $\pm$ 0.00 & 1.05 & 1.11 & 100.00\% $\pm$ 0.00\% & 29.28\% $\pm$ 3.34\% \\
\bottomrule

\end{tabular}
\caption{Comparison of interval length, calibration errors, and coverage across methods. \titleShort achieves substantially shorter intervals than simultaneous and conformal methods while maintaining full coverage of the prediction space and improved coverage of the target probability.
$\mathrm{ECE_{mid}}$ measures if interval midpoint is representative of the true probabilities of the predictions assigned to it.
$\mathrm{ECE_{pred}}$ measures if average prediction agree with the true probability among the points receiving this interval. It evaluates how well the interval grouping preserves the calibration of $p(X)$. Definitions of the metrics are supplemented in Appendix~\ref{app:metrics}. Covariate dimension $d=20$. }% $100$ replications in the $d=20$ logistic set up.
% Relative ECE is ECE divided by the smallest mean ECE across baselines.
\label{tab:task1_multivariate_mc_summary}
\end{table}



\subsection{Task 2 Results}


\begin{figure}[ht]
\centering
\includegraphics[width=\linewidth]{Figures/task2_multiclass_dgp_sweep.pdf}
\caption{Task~2 results, averaged over $1000$ replications. First row: linear-Gaussian DGP. Second row: nonlinear-mixture softmax DGP. Each replication uses disjoint splits of size $n_{\mathrm{train}}=n_{\mathrm{cal}}=n_{\mathrm{eval}}=2000$, with $d=10$ covariates and $\numClass=4$ classes. \titleShort maintains near-target worst-class coverage in both settings. This advantage is more pronounced under the more complex nonlinear-mixture DGP, where competing baselines show substantial degradation in worst-class coverage. Implementation details are given in Appendix~\ref{app:task2_setup}.}
\label{fig:task2}
\end{figure}

Across both the linear-Gaussian and nonlinear-mixture softmax DGPs, \titleShort delivers the most stable worst-class coverage relative to prediction-set size, as illustrated in the first column of Figure~\ref{fig:task2}. This advantage becomes increasingly pronounced as the complexity of the DGP grows. The second and third columns of Figure~\ref{fig:task2} demonstrate that under the nonlinear-mixture DGP, the performance of competing methods deteriorates substantially. In contrast, \titleShort continues to track the target coverage closely while maintaining the smallest prediction sets. Note that the average class coverage shown in the second column is distinct from marginal coverage--standard conformal classification, despite its marginal coverage guarantees, suffers from poor worst-class coverage. Conversely, \titleShort consistently yields class-wise coverage that is lower-bounded by the target level. These empirical observations align exactly with Theorems~\ref{thm:mc_error}, \ref{thm:mc_optimality}, and \ref{thm:mc_same_sample}, suggesting that \titleShort is robust to miscalibration resulting from nonlinear covariate structures and model misspecification.

\section{Related Work}\label{sec:related_work}

\paragraph{Conformal prediction.}
Conformal prediction constructs set-valued predictors $C(X)$ with finite-sample marginal validity,
\[
\mathbb{P}\Big(Y \in C(X)\Big) \ge 1-\alpha,
\]
under exchangeability \citep{vovk2005algorithmic,shafer2008tutorial}. Split, cross-, and jackknife-style variants---including the jackknife+, jackknife-after-bootstrap, conformalized quantile regression, and conformal predictive distributions---show that the framework can be layered on top of a broad range of base learners \citep{barber2021predictive,kim2020predictive,romano2019conformalized,vovk2015cross,vovk2018cross}. A stronger target is conditional coverage,
\[
\mathbb{P}\Big(Y \in C(X)\mid X=x\Big) \ge 1-\alpha \quad \text{for a.e. }x,
\]
but distribution-free versions of this guarantee are impossible except through vacuous sets in the worst case \citep{lei2014distribution,barber2021limits}. Recent work therefore studies weaker conditional notions obtained by coarsening the conditioning event: Mondrian or group-conditional validity conditions on a prespecified partition or label \citep{vovk2003mondrian,vovk2012conditional,martinezgil2024groups}; multivalid coverage asks for simultaneous approximate validity over a rich family of subpopulations \citep{jung2023multivalid,gibbs2025conformal}; and recent analyses clarify sample- and training-conditional targets and when they can be attained under additional assumptions \citep{duchi2025few,liang2025algorithmic}. Related extensions handle covariate shift, label shift, online adaptation, and broader forms of nonexchangeability \citep{tibshirani2019conformal,gibbs2021adaptive,podkopaev2021distribution,barber2022beyondexchangeability, berta2024classifier}.

The \titleShort target,
\[
\mathbb{P}\Big(Y=1 \mid p(X)\in I\Big)\in I,
\]
belongs to this ``nearly conditional'' regime, but the conditioning variable is the model score itself rather than the full feature vector. This distinction is central for our results: the score is exactly the quantity used downstream for thresholding, classification, and decision-making, so conditioning on $p(X)$ avoids the impossibility of full covariate-conditional validity while remaining aligned with the deployed predictor. Theorem~\ref{thm:tp} converts this score-conditional statement into a conformal-style guarantee for the latent mean $p^\ast(X)$, and Theorem~\ref{thm:consistency-prediction} supplies an efficiency statement: under $\lambda$-regularity and $L_1$ score error $\varepsilon$, for all but a $\delta_1$ fraction of score values there exists a \titleShort of width of order $(\log n/n)^{1/3}+\varepsilon/\delta_1$. This complements a separate line of work that studies the length, volume, or oracle optimality of conformal regions. For density-level-set predictors, \citet{lei2013distribution} establish finite-sample validity together with asymptotic efficiency relative to oracle level sets; \citet{lei2018distribution} show that split conformal regression inherits consistency and can adapt interval length to heteroscedasticity; \citet{chernozhukov2021distributional} and \citet{izbicki2020flexible,izbicki2022cdsplit} obtain approximate conditional validity and convergence to oracle conditional intervals or highest-density regions when conditional distribution estimators are consistent; and \citet{yang2025selection,lebars2025volume} quantify finite-sample width losses relative to oracle choices. Relative to these papers, Theorem~\ref{thm:consistency-prediction} studies the shrinking object most relevant for our problem: not a response interval around $Y$, but a score-conditional calibration interval that certifies the meaning of the model's own probability output.

\paragraph{Calibration.}
For a binary score $p(X)\in[0,1]$, perfect calibration means
\[
\mathbb{P}\Big(Y=1\mid p(X)=v\Big)=v,
\]
for all relevant $v$ \citep{dawid1982wellcalibrated,foster1998asymptotic}. In practice this pointwise condition is usually assessed through coarser summaries: reliability diagrams compare empirical frequencies to predicted probabilities on bins of the score, and scalar criteria such as ECE aggregate these discrepancies across bins \citep{gneiting2007strictly,vaicenavicius2019evaluating,carrell2022calibration,blasiok2024smooth}. A large post-hoc calibration literature then modifies the original score via parametric or nonparametric maps---Platt scaling, isotonic regression, Bayesian binning, temperature scaling, beta calibration, and Dirichlet calibration---to improve this global agreement \citep{platt1999probabilistic,zadrozny2002transforming,niculescu2005predicting,naeini2015obtaining,guo2017calibration,kull2017beta,kull2019beyond}. Our objective is different: we do not alter $p$, but instead seek data-adaptive subsets of the score range on which the conditional mean outcome is certified to lie inside the same interval. In this sense, classical conformal prediction controls the random event $\{Y\in C(X)\}$, whereas \titleShort applies conformal-style uncertainty quantification to the conditional mean outcome along the score axis.

This places \titleShort between descriptive calibration diagnostics and conformal uncertainty quantification. Multicalibration asks for simultaneous approximate calibration over many subgroups and is closely tied to fairness constraints \citep{hebert2018multicalibration,kleinberg2016inherent,pleiss2017fairness}. \titleShort can be read as a one-dimensional, score-indexed analogue of this idea: the conditioning sets are subsets of the score axis rather than externally specified demographic groups, and the guarantee is interval-valued rather than an additive calibration error bound. Vovk and collaborators developed self-calibrating forecasters, Venn predictors, and Venn--Abers methods that output calibrated probabilities or probability intervals under exchangeability \citep{vovk2003self,vovk2014vennabers}. More recently, \citet{gupta2020distribution} give distribution-free confidence intervals for calibrated binary probabilities, \citet{podkopaev2021distribution} study such statements under structured shift, and \citet{van2024self} recalibrate the models with conformal guarantees. These works are the closest precedents to our calibration view. The key distinction is that \titleShort produces the conditioning interval and the uncertainty statement simultaneously through the self-consistency condition, which is exactly the structure exploited in our theoretical results.

\section{Conclusions}
\label{sec:conclusions}

In summary, we have introduced \titleShorts as a framework for local, data-driven uncertainty quantification that bridges the gap between calibration and conformal prediction. By conditioning on the model’s predicted probabilities, \titleShorts provide intervals that are both interpretable and adaptive to the precision of the underlying predictor. Unlike traditional calibration methods, which rely on fixed bins or post-hoc adjustments, our approach identifies regions of the prediction space where the model is self-certified to be reliable, yielding uncertainty statements that are directly actionable without modifying the base predictor. We show that, under suitable regularity conditions, the width of these intervals shrinks with increasing sample size, enabling increasingly precise risk characterization. 

Practically, \titleShorts offer a tool for probabilistic prediction and multi-class classification with respective coverage guarantees. Overall, our results highlight that leveraging the model’s own risk stratification can produce rigorous, locally adaptive uncertainty quantification, providing an alternative to conventional recalibration or fully conditional prediction frameworks. This opens up directions on extending \titleShorts to structured settings inherently require prediction-conditioned uncertainty quantification, such as representation-based inference (where the predictions acts as a low-dimensional summary), covariate shift, and distribution shift (where uncertainty guarantees may need to hold conditionally on prediction-induced strata that capture regions of varying model reliability). On the theoretical side, a direction for future work is to develop adaptive rates on the width bound that exploit additional properties or parametric structure of the prediction model. 


% The problem is difficult because $\PP \left(\PP\left(Y = 1|p(X) \in [q-l, q + l]\right) \in [q-l, q + l] \right)$ is not monotone in $l$.

% By concentration and union bound, we have
% \begin{align*}
%     \PP\left(\PP(Y = 1|p(X) \in \I_j) \in \I_j,~\forall \I_j\right) \geq 1-\delta
% \end{align*}
% and therefore with probability at least $1-\delta$:
% \begin{align*}
%     \PP(Y = 1|p(X) \in \I_{x_{\test}}) \in \I_{x_{\test}}.
% \end{align*}
% However, it doesn't guarantee
% \begin{align*}
%     \PP(Y = 1|x_{\test}) \in \I_{x_{\test}}.
% \end{align*}


\clearpage
\bibliographystyle{\xuelinbibstyle}
\bibliography{ref}

\clearpage
\appendix
\section{Proofs}

% section 3
\subsection{Proof of Theorem~\ref{thm:algorithm_population}}
\label{proof_thm:algorithm_population}

\begin{proof}

For each discretized interval $[a, b]$ with empirical count $N_{[a,b]} = \sum_i \mathbbm{1}\{p(x_i) \in [a,b]\} > 0$, we apply Hoeffding's inequality and a union bound to derive that with probability at least $1-\delta$:
\begin{align*}
    |\hat{\mu}_{[a,b]} - \mu_{[a,b]}| \leq 2\sqrt{\frac{\log(|\mathcal{I}|/\delta)}{N_{[a,b]}}},
\end{align*}
where $\mu_{[a,b]} = \PP(Y=1 \mid p(X) \in [a, b])$,
\begin{align*}
    \hat{\mu}_{[a,b]} = \frac{\sum_{i=1}^n y_i \cdot \mathbbm{1}(p(x_i) \in [a,b])}{\sum_{i=1}^n \mathbbm{1}(p(x_i) \in [a,b])}
\end{align*}
and $|\mathcal{I}| = K^2$ is the number of candidate intervals.

On the event that all these bounds hold simultaneously, Line~\ref{line:empirical-interval} implies that
\begin{align*}
    \mu_{[a,b]} \in [a,b].
\end{align*}
This confirms that with probability at least $1-\failureProb$, $\intervalset_p$ is a set of \titleShorts.

We note that a relaxed alternative for Line~\ref{line:empirical-interval} is to add indicator functions at the boundary endpoints 0 and 1. That is, 
{
\begin{align}
    & \frac{\sum_{i=1}^n y_i \cdot \mathbbm{1}(p(x_i) \in [a,b])}{\sum_{i=1}^n \mathbbm{1}(p(x_i) \in [a,b])} \nonumber\\
    &\quad \in \left[a + 2\sqrt{\frac{\log (\discreteNum^2/\delta)}{\sum_{i=1}^n \mathbbm{1}(p(x_i) \in [a,b])}} \mathbbm{1}\{a=0\} ,b - 2\sqrt{\frac{\log (\discreteNum^2/\delta)}{\sum_{i=1}^n \mathbbm{1}(p(x_i) \in [a,b])}}\mathbbm{1}\{b=1\}\right].
\end{align}} 
This formulation may be preferable in edge cases where $Y=0$ or $Y=1$ almost surely. The proof still holds since $\mu_{[a,b]}\in[0,1]$. For simplicity, we proceed with the simpler formulation in the main text.
\end{proof}





% section 4
\subsection{Proof of Theorem~\ref{thm:consistency-prediction}}\label{proof:thm:consistency-prediction}

\begin{proof}
Let 
\begin{align*}
    \epsilon_\stat = &~ 2\Big(\frac{2\log(3n/\delta)}{\lambda n}\Big)^{1/3}, \\
    \epsilon_1 = &~ \epsilon_\stat + \frac{4\epsilon}{\delta_1},\\
    K = &~ \max\left\{1, \left\lfloor\frac{1}{4\epsilon_1}\right\rfloor\right\}. 
\end{align*} 
If $K=1$, the interval $[0,1]$ is returned by the algorithm and covers all prediction values, proving the result directly with $C(n,\epsilon,\delta,\delta_1)=1$. For $K \geq 2$, the proof consists of three steps. 

\paragraph{Step 1. } 
We show that there exists $J \in \Z_+$ and disjoint intervals $[a_j,b_j], j \in [J]$ such that:
\begin{enumerate}
    \item For any $j \in [J]$, 
    \begin{align}\label{eq:consistency-existence-calibration}
        a_j + \epsilon_\stat \cdot \mathbbm{1}(a_j \neq 0) \leq \PP[Y=1\mid \estModel(X) \in [a_j,b_j]] \leq b_j - \epsilon_\stat \cdot \mathbbm{1}(b_j \neq 1);
    \end{align}
    \item 
    \begin{align}\label{eq:consistency-existence-coverage}
        \PP\Big(\estModel(X) \notin \cup_{j \in [J]} [a_j,b_j]\Big) \leq \delta_1;
    \end{align}
    \item For any $j \in [J]$,
    \begin{align}\label{eq:consistency-existence-length}%56
        4\epsilon_1 \leq \Big|b_j - a_j\Big| \leq 56\epsilon_1 \cdot \log(6/\delta_1).
    \end{align}
\end{enumerate}
First, applying Lemma~\ref{lem:dist-regularity} with $\nu$ being the distribution of $\estModel(X)$, uncovered mass budget $\delta_1$, and accuracy parameter $\epsilon_1$, we see that there exist disjoint intervals $[a_i,b_i], i \in [L]$ that satisfy 
\begin{align*}
    \PP\Big(\estModel(X) \notin \cup_{i \in [L]} [a_i,b_i]\Big) \leq \delta_1/2,
\end{align*}
Eq.~\eqref{eq:consistency-existence-length}, and
\begin{align}\label{eq:regularity-pm}
    a_i + \epsilon_1 \cdot \mathbbm{1}(a_i \neq 0) \leq \E[\estModel(X)\mid\estModel(X) \in [a_i,b_i]] \leq b_i - \epsilon_1 \cdot \mathbbm{1}(b_i \neq 1).
\end{align}
By Lemma~\ref{lem:conditional-dist-shift}, there exist $i_1,\dots,i_J \in [L]$ such that for any $j \in [J]$, 
\begin{align}\label{eq:conditional-dist-shift-pm}
    \Big|\PP(Y=1\mid\estModel(X) \in [a_{i_j},b_{i_j}])  - \E[\estModel(X)\mid\estModel(X) \in [a_{i_j},b_{i_j}]] \Big| \leq 4\epsilon/\delta_1 
\end{align}
and
\begin{align*}
    \PP\Big(\estModel(X) \in \cup_{j' \notin \{i_1,\dots,i_J\}} [a_{j'},b_{j'}]\Big) \leq \delta_1/2.
\end{align*}
Combining Eq.~\eqref{eq:conditional-dist-shift-pm} and Eq.~\eqref{eq:regularity-pm}, we have
\begin{align*}
    \MoveEqLeft{ \PP(Y = 1 \mid\estModel(X) \in [a_{i_j},b_{i_j}])} \\
    & \leq  \E[\estModel(X)\mid\estModel(X) \in [a_{i_j},b_{i_j}]] + 4\epsilon/\delta_1 \\
    & \leq  b_{i_j} - \epsilon_1 \cdot \mathbbm{1}(b_{i_j} \neq 1) + 4\epsilon/\delta_1 \\
    & =  b_{i_j} - \epsilon_\stat \cdot \mathbbm{1}(b_{i_j} \neq 1), 
\end{align*}
where the last step follows from the definition of $\epsilon_1$, and
\begin{align*}
    \MoveEqLeft{ \PP(Y = 1\mid\estModel(X) \in [a_{i_j},b_{i_j}])} \\
    & \geq  \E[\estModel(X)\mid\estModel(X) \in [a_{i_j},b_{i_j}]] - 4\epsilon/\delta_1\\
    & \geq   a_{i_j} + \epsilon_1 \cdot \mathbbm{1}(a_{i_j} \neq 0) - 4\epsilon/\delta_1\\
    & =  a_{i_j} + \epsilon_\stat \cdot \mathbbm{1}(a_{i_j} \neq 0) \tag{if $a_{i_j} = 0$ then trivially $a_{i_j} \leq \PP(Y=1\mid\estModel(X) \in [a_{i_j},b_{i_j}])$}.
\end{align*}
This establishes 
\begin{align*}
    a_{i_j} + \epsilon_\stat \cdot \mathbbm{1}(a_{i_j} \neq 0) \leq \PP(Y=1\mid\estModel(X) \in [a_{i_j},b_{i_j}]) \leq b_{i_j} - \epsilon_\stat \cdot \mathbbm{1}(b_{i_j} \neq 1),~\forall j \in [J].
\end{align*}
Furthermore,
\begin{align*}
    \PP\Big(\estModel(X) \notin \cup_{j \in [J]} [a_{i_j},b_{i_j}]\Big) \leq &~ \PP\Big(\estModel(X) \notin \cup_{i \in [L]} [a_i,b_i]\Big) + \PP\Big(\estModel(X) \in \cup_{j' \notin \{i_1,\dots,i_J\}} [a_{j'},b_{j'}]\Big) \\
    \leq &~ \delta_1.
\end{align*}
In sum, the intervals $[a_{i_1},b_{i_1}],\dots,[a_{i_J},b_{i_J}]$ satisfy Eq.~\eqref{eq:consistency-existence-calibration}-\eqref{eq:consistency-existence-length} (where we view $[a_{i_j},b_{i_j}]$ as $[a_{j},b_{j}]$).

\paragraph{Step 2. } 
We show that with probability at least $1-\delta$ over the randomness of $\{(x_i,y_i)\}_{i=1}^n$, the intervals in Step 1 will be included in the set of intervals $\I$ output by Algorithm~\ref{alg:calibration_population}. 
For simplicity of notation, we let $[a_{1},b_{1}],\dots,[a_J,b_J]$ denote the intervals found in Step 1; i.e. $[a_{1},b_{1}],\dots,[a_J,b_J] = [a_{i_1},b_{i_1}],\dots,[a_{i_J},b_{i_J}]$. Let $z_i = p(x_{m+i})$ be a shorthand for the predicted probability.
Define events
\begin{align*}
    A_1 = &~ \Big\{\sumin \mathbbm{1}(z_i \in [a_j,b_j]) \geq 2\lambda \epsilon_1 n,\forall j \in [J]\Big\}\\
    A_2 = &~ \Big\{\Big|\frac{\sum_{i=1}^n y_i \cdot \mathbbm{1}(z_i \in [a_j,b_j])}{\sum_{i=1}^n \mathbbm{1}(z_i \in [a_j,b_j])} - \PP(Y=1\mid \estModel(X) \in [a_j,b_j])\Big| \leq \epsilon_\stat/2, \forall j \in [J]\Big\}.
\end{align*}
Next, we show that under events $A_1, A_2$, the intervals $[a_{1},b_{1}],\dots,[a_J,b_J]$ are included in $\I$ by Algorithm~\ref{alg:calibration_population}. 
Indeed, we have
\begin{align*}
    \frac{\sum_{i=1}^n y_i \cdot \mathbbm{1}(z_i \in [a_j,b_j])}{\sum_{i=1}^n \mathbbm{1}(z_i \in [a_j,b_j])} \leq &~ \PP(Y=1\mid\estModel(X) \in [a_j,b_j]) + \epsilon_\stat/2\\
    \leq &~ b_{j} - \epsilon_\stat \cdot \mathbbm{1}(b_{j} \neq 1) + \epsilon_\stat/2\\
    \leq &~  b_{j} - \epsilon_\stat/2 \cdot \mathbbm{1}(b_{j} \neq 1)\\
    \leq &~ b_j - 2\sqrt{\frac{\log (2K^2/\delta)}{\sum_{i=1}^n \mathbbm{1}(z_i \in [a_j,b_j])}},
\end{align*}
where the last inequality is due to
\begin{align*}
    2\sqrt{\frac{\log (2K^2/\delta)}{\sum_{i=1}^n \mathbbm{1}(z_i \in [a_j,b_j])}} \leq &~ 2\sqrt{\frac{2\log(3n/\delta)}{2\lambda \epsilon_1 n}}\\
    \leq &~ \Big(\frac{2\log(3n/\delta)}{\lambda n}\Big)^{1/3} = \epsilon_\stat/2,
\end{align*}
under event $A_1$. 
Similarly,
\begin{align*}
    \frac{\sum_{i=1}^n y_i \cdot \mathbbm{1}(z_i \in [a_j,b_j])}{\sum_{i=1}^n \mathbbm{1}(z_i \in [a_j,b_j])} \geq a_j + 2\sqrt{\frac{\log (2K^2/\delta)}{\sum_{i=1}^n \mathbbm{1}(z_i \in [a_j,b_j])}}.
\end{align*}
Therefore, the interval $[a_j,b_j]$ will be included in $\I$ under $A_1 \cap A_2$. 

Next, we show $\PP(A_1\cap A_2)\ge 1-\delta$. Since 
$
    \Big|b_j - a_j\Big| \geq 4\epsilon_1,
$ 
applying Lemma~\ref{lem:sample-size-interval} with $c=4\epsilon_1$, 
we have 
\begin{align*}
    \PP(A_1) \geq &~ 1-J\exp\Big(-\frac{\lambda\epsilon_1 n}{2}\Big)\\
    \geq &~ 1- n\exp\Big(-\frac{\lambda\epsilon_\stat n}{2}\Big)\\
    \geq &~ 1- n\exp\Big(-\log(3n/\delta)\Big) \tag{$n \geq \lambda^{-1} \log(3n/\delta) / \sqrt{2}$}\\
    \geq &~ 1-\delta/3 .
\end{align*}

Under $A_1$, by Hoeffding's inequality and a union bound, we have
\begin{align*}
    \PP(A_2\mid A_1) \geq &~ 1- 2 J\exp\Big(-\epsilon_\stat^2 \cdot\sum_{i=1}^n \mathbbm{1}(z_i \in [a_j,b_j])/2\Big)\\
    \geq &~ 1- 2 n\exp\Big(-\epsilon_\stat^2 \cdot \lambda \epsilon_1 n\Big)\\
    \geq &~ 1- 2 n\exp\Big(-\epsilon_\stat^3 \cdot \lambda  n\Big)\\
    \geq &~ 1- 2 n\exp\Big(- \frac{16\log(3n/\delta)}{\lambda n } \cdot \lambda  n\Big)\\
    \geq &~ 1- 2 n\exp\Big(- 16\log(3n/\delta)\Big)\\
    \geq &~ 1- 2 n (3n/\delta)^{- 16}\\
    \geq &~ 1- \frac{2 \delta^{16}}{3^{16} n^{15}}\\
    \geq &~ 1-\delta/6.
\end{align*}
Combining, we have
\begin{align*}
    \PP(A_1 \cap A_2) \geq 1 - \delta.
\end{align*}

\paragraph{Step 3. } We show that the intervals $[a_{1},b_{1}],\dots,[a_J,b_J]$ in Step 2 satisfy the property
\begin{align*}
    \PP_{x \sim \mu}\Big(\exists [a_j,b_j]  ~s.t.~ |b_j-a_j| \leq C(n,\epsilon,\delta,\delta_1) ~\text{and}~ p(x) \in [a_j,b_j] \Big) \geq 1-\delta_1,
\end{align*}
where 
\begin{align*}%56
    C(n,\epsilon,\delta,\delta_1) &= \log(6/\delta_1) \cdot 56\Big(2\Big(\frac{2\log(3n/\delta)}{\lambda n}\Big)^{1/3}  + \frac{4\epsilon}{\delta_1}\Big).
\end{align*}
Define event
\begin{align*}
    A_3 = \Big\{\exists [a_j,b_j]  ~s.t.~ |b-a| \leq C(n,\epsilon,\delta,\delta_1) ~\text{and}~ p(x) \in [a_j,b_j] \Big\}.
\end{align*}
By Eq.~\eqref{eq:consistency-existence-length}, for any $j \in [J]$,
\begin{align*}
    |b_j-a_j| \leq C(n,\epsilon,\delta,\delta_1).
\end{align*}
It follows from Eq.~\eqref{eq:consistency-existence-coverage} that
\begin{align*}
    \PP_{x \sim \mu}\Big(\exists [a_j,b_j]  ~s.t.~ |b_j-a_j| \leq C(n,\epsilon,\delta,\delta_1) ~\text{and}~ p(x) \in [a_j,b_j] \Big) = &~ \PP\Big(\estModel(X) \in \cup_{j \in [J]} [a_j,b_j]\Big) \\
    \geq &~ 1- \delta_1.
\end{align*}
\end{proof}

\subsection{Proof of Lemma~\ref{lem:sample-size-interval}}\label{proof_lem:sample-size-interval}
\begin{proof}
Fix an interval $[a_j, b_j]$ for some $j \in [J]$.
Define $w_i := \mathbbm{1}(z_i \in [a_j,b_j])$ for $i \in [n]$. Then $\{w_i\}_{i=1}^n$ are i.i.d.\ Bernoulli random variables with mean
\begin{align*}
    \E[w_i] 
    = \PP(z_i \in [a_j,b_j])
    \ge \lambda (b_j-a_j)
    \ge \lambda c. 
\end{align*}
Let $S_j := \sumin w_i$. The total expected sample count satisfies $\E[S_j] \ge \lambda cn$. 
By the lower-tail multiplicative Chernoff bound,
$$\mathbb{P}(S_j \le (1-\delta)\E[S_j]) \le \exp(-\delta^2 \E[S_j]/2).$$ 
Applying it with $\delta = 1/2$,  
\begin{align*}
    \PP\Big(\sumin w_i < \lambda cn/2\Big) 
    \le &~ \PP\Big(S_j \le \frac{\E[S_j]}{2}\Big) \\
    \le &~ \exp\Big(-\frac{\E[S_j]}{8}\Big) \\
    \le &~ \exp\Big(-\frac{\lambda cn}{8}\Big).
\end{align*}
By a union bound over all $j \in [J]$, with probability at least $1-J\exp\Big(-\frac{\lambda cn}{8}\Big)$,
\begin{align*}
    \sumin \mathbbm{1}(z_i \in [a_j,b_j]) \geq \lambda cn/2
\end{align*}
holds for all $j \in [J]$. This completes the proof.
\end{proof}

\subsection{Proof of Lemma~\ref{lem:dist-regularity}}\label{proof_lem:dist-regularity}


\begin{proof}
We present a constructive proof based on discretization of $[0,1]$. To simplify notation, we abbreviate $\E_{Z \sim \nu}$ as $\E$ and $\PP_{Z \sim \nu}$ as $\PP$.


\paragraph{Step 1: Construct the graph.}
Consider $[(i-1)/\discreteNum, i/\discreteNum], i = 1,2, \dots,\discreteNum$ with $\discreteNum = \lfloor 1/(4\epsilon) \rfloor$. Without loss of generality assume $\discreteNum\ge2$.
Construct a directed graph $G = (V,E)$ where $V = [\discreteNum]$ and let the edges be defined as follows: $1 \ra 2, \discreteNum-1 \la \discreteNum$, where there is a directed edge $i \ra j$ iff both of the following conditions hold:
\begin{itemize}
    \item $|i-j| = 1$, and
    \item Either \begin{align}\label{eq:edge-condition-1}
   \PP(Z\in [(i-1)/\discreteNum, i/\discreteNum]) \text{ and } \Big(\E\Big[Z\mid Z \in [(i-1)/\discreteNum, i/\discreteNum]\Big] - \frac{i-0.5}{\discreteNum}\Big) \cdot (j-i) \ge 0 ,
\end{align}
or 
\begin{align}\label{eq:edge-condition-2}
    \PP\Big(Z \in [(i-1)/\discreteNum, i/\discreteNum]\Big) \leq 3 \cdot \PP\Big(Z \in [(j-1)/\discreteNum, j/\discreteNum]\Big).
\end{align}
\end{itemize}

Eq.~\eqref{eq:edge-condition-1} can be interpreted as: if $\E\Big[Z\mid Z \in [(i-1)/\discreteNum, i/\discreteNum]\Big]$ is closer to $(i-1)/\discreteNum$, then draw a directed edge $i-1 \la i$; if $\E\Big[Z\mid Z \in [(i-1)/\discreteNum, i/\discreteNum]\Big]$ is closer to $i/\discreteNum$, then draw a directed edge $i \ra i+1$. It directs an edge toward the side of the bin on which its conditional mean lies.

Eq.~\eqref{eq:edge-condition-2} can be interpreted as: we draw a directed edge to $i \ra j$ (for $j = i-1$ or $i+1$) if the probability mass on $[(i-1)/\discreteNum, i/\discreteNum]$ is not substantially greater than that of $[(j-1)/\discreteNum, j/\discreteNum]$. It directs an edge toward a neighboring bin whose mass is not substantially smaller.

By construction, it is obvious that every vertex has at least one outgoing edge since for positive-mass vertex Eq.~\eqref{eq:edge-condition-1} holds either for $j = i-1$ or $j=i+1$, and for zero-mass vertex Eq.~\eqref{eq:edge-condition-2} holds for both neighbors.


There is at least one bidirectional adjacent pair. The existence of such pair follows from a continuity argument: since 
$1 \ra 2, \discreteNum-1 \la \discreteNum$, and every vertex has at least one outgoing edge, let $k^*\in \{2,\dots,\discreteNum\}$ be the smallest integer such that there exists an edge $k^* -1 \la k^*$. Then $k^* -1 \ra k^*$ because every vertex has at least one outgoing edge, and if the outgoing edge for $k^*-1$ is $k^*-1 \rightarrow k^*-2$, then  $k^*$ would not be the smallest vertex as defined. This means that $k^* - 1 \leftrightarrows k^*$ forms a bidirectional pair, thus confirming the existence.

Now, denote the collection of $H$ maximal blocks of consecutive vertices connected by bidirectional edges as
\begin{align*}
    \{D_i: D_i = \{l_i, l_i+1, \dots, r_i\} , \text{ where }l_i\leftrightarrows l_i+ 1\leftrightarrows\ldots\leftrightarrows r_i, ~i\in[H] \}
\end{align*}
Then we have 
\begin{align*}
    1\le l_1 < r_1 < l_2 < r_2 < \ldots < l_H \le \discreteNum.
\end{align*}

\paragraph{Step 2: Construct intervals from the graph.}
We now construct intervals, where we first subdivide each block into short groups, and then assign the vertices between consecutive blocks to one of the two neighboring groups.

We now divide each bidirectional block $D_i$ into $m_i:=\lfloor(r_i- l_i+1)/2\rfloor$ consecutive groups, whereas every group has two vertices, except that the last has three when
$r_i- l_i+1$ is odd. 
That is, denote the $j$-th group as
$\{\alpha_{i,j},\dots,\beta_{i,j}\}$, with endpoints
\begin{align}\label{eq:block-core-endpoints}
 \alpha_{i,j}&= l_i+2(j-1),\nonumber\\
 \beta_{i,j}&=
 \begin{cases}
   l_i+2j-1,&j<m_i,\\
  r_i,&j=m_i.
 \end{cases}
\end{align}

Notice that for any $i\in[H-1]$, all integer vertices between two consecutive bidirectional blocks must be connected separately by a leftwards chain and a rightwards chain (upon merging duplicated edges) as one move from left to right. That is,
\begin{align*}
    r_i \la \cdots \la w_i \ra \cdots \ra l_{i+1},
\end{align*}
for some $w_i
\in \{r_i, \dots, l_{i+1}\}$. 
(Otherwise, there exists another bidirectional block between $r_i$ and $l_{i+1}$.) We call $w_i$ as an intervening source vertex, where we assign it the block on its left when constructing intervals. Denote the endpoints of constructed intervals as $\{u_i, v_i\}, i\in [H]$.
Define the cut points and the ranges assigned to the blocks by
\begin{align}\label{eq:block-assignment}
 \omega_0&=0,\qquad\omega_H=\discreteNum,\qquad
 \omega_i=\min\{w_i, l_{i+1}-1\}\quad(i\in[H-1]),\nonumber\\
 u_i&=\omega_{i-1}+1,\qquad v_i=\omega_i\quad(i\in[H]).
\end{align}
They satisfy
\begin{align*}
 u_i\leq l_i<r_i\leq v_i,
 \qquad u_{i+1}=v_i+1\quad(i\in[H-1]),
\end{align*}
and all edges in each assigned range have the form
\begin{align}\label{eq:block-with-chains}
 u_i\ra\cdots\ra l_i\leftrightarrows\cdots
 \leftrightarrows r_i\la\cdots\la v_i.
\end{align}



Truncate the chains in
Eq.~\eqref{eq:block-with-chains} after at most $\lceil \log(3/\delta)\rceil$ vertices by setting
\begin{align}\label{eq:truncated-block-endpoints}
 s_i=u_i\vee( l_i-\lceil \log(3/\delta)\rceil),\qquad
 t_i=v_i\wedge(r_i+\lceil \log(3/\delta)\rceil).
\end{align}
Attach the retained left chain to the first group and the retained
right chain to the last group. That is, set
\begin{align*}
%\label{eq:retained-group-endpoints}
 s_{i,j}&=
 \begin{cases}
  s_i,&j=1,\\
  \alpha_{i,j},&j>1,
 \end{cases}
 &
 t_{i,j}&=
 \begin{cases}
  \beta_{i,j},&j<m_i,\\
  t_i,&j=m_i.
 \end{cases}
\end{align*}
Define intervals
\begin{align}\label{eq:constructed-intervals}
 I_{i,j}=[a_{i,j},b_{i,j}]
 =[(s_{i,j}-1)/\discreteNum,t_{i,j}/\discreteNum],
 \qquad i\in[H],\ j\in[m_i].
\end{align}
An example illustration of the intervals (in boxes, without truncation) as: 

\begin{align}\label{eq:endpoint_illustration}
  \cdots \la v_{i-1} \ra \underbrace{\boxed{ u_i\ra\cdots\ra l_i\leftrightarrows l_i + 1 } }_{\text{attach chain to first block pair}}\leftrightarrows \underbrace{\boxed{ l_i + 2 \leftrightarrows l_i + 3} }_{\text{grouped by 2}} \leftrightarrows   \cdots 
 \leftrightarrows \underbrace{\boxed{ r_i - 1 \leftrightarrows  r_i\la\cdots\la v_i}}_{\text{attach chain to last block pair}}  \ra \cdots .
\end{align}

Truncation could be illustrated as:

\begin{align}\label{eq:truncation_illustration}
 \cdots \la v_{i-1} \ra  u_i\ra\cdots\ra \underbrace{\boxed{l_i - \lceil \log(3/\delta)\rceil\ra \cdots\ra l_i\leftrightarrows l_i + 1 } }_{\text{attach truncated chain}}\leftrightarrows \underbrace{\boxed{ l_i + 2 \leftrightarrows l_i + 3} }_{\text{grouped by 2}} \leftrightarrows   \cdots 
 .
\end{align}




 
% Its
% two endpoint edges point inward:
% \begin{align*}
% %\label{eq:inward-endpoint-edges}
%  s_{i,j}\ra s_{i,j}+1,
%  \qquad t_{i,j}-1\la t_{i,j}.
% \end{align*} 
Remove any intervals of zero probability. Relabel those remaining as
$I_i=[a_i,b_i]$, $i\in[L]$, in increasing order. 
%Then $L\leq\sum_{i=1}^H m_i\leq\lfloor K/2\rfloor$.



% for each $i = 1,2,\dots,L$, let $u_i$ be the smallest integer such that $u_i > v_{i-1}$ (take $k_0 = v_0 = 0$) and there exists a rightwards chain $u_i \ra  \cdots \ra  k_{i}-1$ (if there exists no such integer, then $u_i = k_{i}-1$), and $v_i$ be the greatest integer such that  $v_i < k_{i+1}$ ($k_{L+1} = K$) and there exists a leftwards chain $k_i \la \cdots \la v_i$ (if there exists no such chain, then $v_i = k_i$). Therefore, the edges between $u_i$ and $v_i$ are:
% \begin{align*}
%     u_i \ra  \cdots \ra  k_{i} \leftrightarrows k_i + 1 \la \cdots \la v_i
% \end{align*}
% Let interval $\cI_i = [a_i,b_i]$ be defined by
% \begin{align}
%     \cI_i = \Big[\frac{u_i \vee (k_i - \lceil\log(3/\delta)\rceil) -1}{K}, \frac{v_i \wedge (k_i + \lceil\log(3/\delta)\rceil)}{K}\Big],
% \end{align}
% i.e.,
% \begin{align*}
%     a_i = \frac{u_i \vee (k_i - \lceil\log(3/\delta)\rceil) -1}{K}, ~~ b_i = \frac{v_i \wedge (k_i + \lceil\log(3/\delta)\rceil)}{K}.
% \end{align*}







\paragraph{Step 3: Establish theoretical properties.}


We show that the intervals $[a_i,b_i], i\in[L]$ satisfy the desired properties. 
By construction each interval contains between $1$ and $2\lceil \log(3/\delta)\rceil+3$ vertices, thus Eq.~\eqref{eq:regular-condition-length} is satisfied.
To show Eq.~\eqref{eq:regular-condition}, Fix $i \in [L]$. 
% Let $s_i = u_i \vee (k_i - \log(3/\delta))$ and $t_i = v_i \wedge (k_i + \log(3/\delta))$. 
By $s_i \geq u_i, t_i \leq v_i$ and the definition of $u_i,v_i$ (after relabeling), there exist edges $s_i \ra s_i + 1$ and $t_i - 1 \la t_i$. We will show that $s_i \ra s_i + 1$ implies that $\E[Z\mid Z \in [a_i,b_i]] \geq a_i + \epsilon \cdot \mathbbm{1}(a_i \neq 0)$ by considering the following three cases of $s_i \ra s_i + 1$:

\paragraph{Case 1: $s_i = 1$. } Then $a_i=0$ and 
trivially
$
    \E[Z\mid Z \in [a_i,b_i]] \geq a_i + \epsilon \cdot \mathbbm{1}(a_i \neq 0).
$



\paragraph{Case 2:  }
The edge is drawn by Eq.~\eqref{eq:edge-condition-1}, which is
\begin{align*}
    \E\Big[Z\mid Z \in [(s_i-1)/\discreteNum, s_i/\discreteNum]\Big] \geq  \frac{s_i-0.5}{\discreteNum} .
\end{align*}
Then we have
\begin{align*}
    \E[Z\mid Z \in [a_i,b_i]] = &~ \sum_{j=s_i}^{t_i} \frac{\PP(Z \in [(j-1)/\discreteNum, j/\discreteNum])}{\PP(Z \in [a_i,b_i])} \cdot \E[Z\mid Z \in [(j-1)/\discreteNum, j/\discreteNum]]\\
    \geq &~ \E\Big[Z\mid Z \in [(s_i-1)/\discreteNum, s_i/\discreteNum]\Big] \\
    \geq &~ \frac{s_i-0.5}{\discreteNum} \\
    \geq &~ a_i + \epsilon.
\end{align*}
The first step comes from the law of total expectation and the fact that $[(j-1)/\discreteNum, j/\discreteNum], j=s_i,\dots, t_i$ forms a partition of $[a_i,b_i]$; the second step comes from $\E[Z\mid Z \in [(j-1)/\discreteNum, j/\discreteNum]] \geq \E\Big[Z\mid Z \in [(s_i-1)/\discreteNum, s_i/\discreteNum]\Big]$ for all $j \geq s_i$; the third step comes from Eq.~\eqref{eq:edge-condition-1}; and the final step comes from
$a_i = (s_i-1)/\discreteNum,$ $ \epsilon \le 1/(4\discreteNum)$, so $ a_i + \epsilon \le (s_i-3/4)/\discreteNum \le (s_i - 0.5)/\discreteNum$.



\paragraph{Case 3: } 
The edge is drawn by Eq.~\eqref{eq:edge-condition-2}, which is
\begin{align*}
    \PP\Big(Z \in [(s_i-1)/\discreteNum, s_i/\discreteNum]\Big) \leq 3 \cdot \PP\Big(Z \in [s_i/\discreteNum, (s_i+1)/\discreteNum]\Big).
\end{align*}
Under Case 3, we have
\begin{align*}
    \E[Z\mid Z \in [a_i,b_i]] = &~ \sum_{j=s_i}^{t_i} \frac{\PP(Z \in [(j-1)/\discreteNum, j/\discreteNum])}{\PP(Z \in [a_i,b_i])} \cdot \E[Z \mid Z \in [(j-1)/\discreteNum, j/\discreteNum]]\\
    \geq &~ \frac{\PP(Z \in [(s_i-1)/\discreteNum, s_i/\discreteNum])}{\PP(Z \in [a_i,b_i])} \cdot \E[Z \mid Z \in [(s_i-1)/\discreteNum, s_i/\discreteNum]] \\
    &~ + \left(1-\frac{\PP(Z \in [(s_i-1)/\discreteNum, s_i/\discreteNum])}{\PP(Z \in [a_i,b_i])}\right) \cdot \E[Z \mid Z \in [s_i/\discreteNum, (s_i+1)/\discreteNum]] \\
    \geq &~ 3/4 \cdot (s_i-1)/\discreteNum + 1/4 \cdot s_i/\discreteNum\\
    \ge &~ a_i + \epsilon. 
\end{align*}

The first step comes from the law of total expectation and $[(j-1)/\discreteNum, j/\discreteNum], j=s_i,\dots, t_i$ forms a partition of $[a_i,b_i]$; the second step comes from pulling the first term out $j=s_i$ and then using $\E[Z\mid Z \in [(j-1)/\discreteNum, j/\discreteNum]] \geq \E\Big[Z\mid Z \in [s_i/\discreteNum, (s_i+1)/\discreteNum]\Big]$ for all $j \geq s_i+1$; the third step comes from Eq.~\eqref{eq:edge-condition-2}; the last step comes from $a_i = (s_i-1)/\discreteNum,$ $ \epsilon \le 1/(4\discreteNum)$, so $ a_i + \epsilon \le (s_i-3/4)/\discreteNum$.


%In sum, $s_i \ra s_i + 1$ implies that $\E[Z \mid Z \in [a_i,b_i]] \geq a_i + \epsilon \cdot \mathbbm{1}(a_i \neq 0)$. 
Repeating similar arguments we can show $t_i - 1 \la t_i $ implies that $\E[Z\mid Z \in [a_i,b_i]] \leq b_i - \epsilon \cdot \mathbbm{1}(b_i \neq 1)$. Therefore, Eq.~\eqref{eq:regular-condition} is confirmed. 


Finally, we establish Eq.~\eqref{eq:regular-condition-coverage} by bounding the probability of the cells removed when the
chains were truncated. 
Define $q_i:=\PP( Z \in [(i-1)/\discreteNum, i/\discreteNum])$. Recall that as demonstrated, on the left chain assigned to block $D_i$,
the reverse edge is absent at every step. In particular, for each
$d\in\{1,\dots, l_i-u_i\}$, the edge
$ l_i-d+1\ra l_i-d$ does not exist and thus Eq.\eqref{eq:edge-condition-2} is violated. Therefore, 
\begin{align*}
 q_{ l_i-d+1}>3q_{ l_i-d}.
\end{align*}
Iterating this inequality along the chain gives
\begin{align*}
 q_{ l_i-d}<3^{-d}q_{ l_i}.
\end{align*}
We then have 
\begin{align*}    \sum_{d=\lceil\log(3/\delta)\rceil+1}^{ l_i-u_i}q_{ l_i-d} < \sum_{d=\lceil\log(3/\delta)\rceil+1}^{ l_i-u_i}3^{-d} q_{ l_i} .
\end{align*}
The analogous argument on the right chain ($d\in\{1,\dots,v_i-r_i\}$) gives
\begin{align*}
 q_{r_i+d}&<3^{-d}q_{r_i}
\\ \sum_{d=\lceil\log(3/\delta)\rceil+1}^{v_i-r_i}q_{r_i+d} &< \sum_{d=\lceil\log(3/\delta)\rceil+1}^{v_i-r_i}3^{-d}q_{r_i}.
\end{align*}

Combining, the total mass sum of discarded cells on two sides
is at most
\begin{align*}
 \sum_{d=\lceil\log(3/\delta)\rceil+1}^{ l_i-u_i}q_{ l_i-d}
 +\sum_{d=\lceil\log(3/\delta)\rceil+1}^{v_i-r_i}q_{r_i+d} &\leq (q_{ l_i}+q_{r_i})\sum_{d=\lceil\log(3/\delta)\rceil+1}^{\infty}3^{-d}\\
 &=\frac12\,3^{-\lceil\log(3/\delta)\rceil}(q_{ l_i}+q_{r_i}) \tag{sum of geometric series}.
\end{align*}

 Recall that each point can at most belong to two of the intervals,
\begin{align*}
 \sum_{i=1}^H(q_{ l_i}+q_{r_i})\leq2,
\end{align*}
by union bound,
\begin{align*}
 \PP\Big(Z\notin\cup_{i\in[L]}[a_i,b_i]\Big)
 &\leq\frac12\,3^{-\lceil\log(3/\delta)\rceil}\sum_{i=1}^H(q_{ l_i}+q_{r_i})\\
 &\leq3^{-\lceil\log(3/\delta)\rceil}\leq e^{-\lceil\log(3/\delta)\rceil}\leq\delta/3\leq\delta.
\end{align*}

This proves Eq.~\eqref{eq:regular-condition-coverage} and completes
the proof.

\end{proof}

\subsection{Proof of Lemma~\ref{lem:conditional-dist-shift}}
\label{proof_lem:conditional-dist-shift}

\begin{proof}
Let $j_1,j_2,\dots,j_M$ be all integers in $[L]$ such that Eq.~\eqref{eq:distribution-shift} is not satisfied. If $M=0$, then all intervals satisfy Eq.~\eqref{eq:distribution-shift}, and choosing $\{i_1,\dots,i_J\} = [L]$ completes the proof.

Otherwise, for any $w \in \{j_1,j_2,\dots,j_M\}$, we have $\PP(V \in [a_{w},b_{w}]) > 0$ and
\begin{align*}
    \frac{2\epsilon}{\delta} < &~ \Big|\E\Big[U\mid V \in [a_{w},b_{w}]\Big] - \E\Big[V\mid V \in [a_{w},b_{w}]\Big]\Big|\\
    \leq &~ \E\Big[|U-V| \mid V \in [a_{w},b_{w}]\Big] \tag{by the triangle inequality}\\
    = &~ \frac{1}{\PP(V \in [a_{w},b_{w}])} \cdot \E\Big[|U-V| \cdot \mathbbm{1}(V \in [a_{w},b_{w}])\Big].
\end{align*}
Each point belongs to at most two intervals, so  $\sum_{w \in \{j_1,\dots,j_M\}} \mathbbm{1}(V \in [a_w,b_w]) \leq 2$. It follows that
\begin{align*}
    2\epsilon \geq &~ 2\E[|U-V|]\\
    \geq &~ \E\left[|U-V| \sum_{w \in \{j_1,\dots,j_M\}} \mathbbm{1}(V \in [a_{w},b_{w}])\right]\\
    = &~ \sum_{w \in \{j_1,j_2,\dots,j_M\}}\E\Big[|U-V| \cdot \mathbbm{1}(V \in [a_{w},b_{w}])\Big]\\
    > &~ \sum_{w \in \{j_1,j_2,\dots,j_M\}}\PP(V \in [a_{w},b_{w}]) \cdot \frac{2\epsilon}{\delta}\\
    \geq &~ \PP\Big(V \in \cup_{w \in \{j_1,j_2,\dots,j_M\}} [a_w,b_w]\Big) \cdot \frac{2\epsilon}{\delta}.
\end{align*}
Dividing both sides by $2\epsilon/\delta$ yields
\begin{align*}
    \PP\Big(V \in \cup_{w \in \{j_1,j_2,\dots,j_M\}} [a_w,b_w]\Big) < \delta.
\end{align*}
Setting $\{i_1,\dots,i_J\} := [L] \setminus \{j_1,\dots,j_M\}$, the remaining intervals satisfy Eq.~\eqref{eq:distribution-shift} by construction and Eq.~\eqref{eq:distribution-shift-coverage}.
\end{proof}

% section 5
% \subsection{Proof of Theorem~\ref{thm:tp}} 
% \label{proof_thm:tp}

% \begin{proof}
% Consider any disjoint set of \titleShorts $\{\interval_\iteratorInterval\}_{\iteratorInterval=1}^{\numInterval}$ and let $\epsilon > 0$.
% % , let $\interval_{\numInterval+1} = [0,1]/\cup_{\iteratorInterval=1}^\numInterval\interval_\iteratorInterval$. 
% We have 
% \begin{align*}
%     \PP\Big(p^*(X) \notin C_{\tp}(X;\epsilon)\Big) = &~ \PP\Big(p^*(X) \notin E^\epsilon\Big(\Big[l(p(X)),u(p(X))\Big]\Big)\Big)\\
%     %= &~ \PP\Big(p^*(X) \notin E^\epsilon\Big(\interval_k\Big)\Big) \quad \text{for some }k\\
%     %= &~ \sum_{\iteratorInterval=1}^{\numInterval+1} \PP\Big(p^*(X) \notin E^\epsilon\Big(\interval_k\Big)\mid p(X) \in \interval_\iteratorInterval\Big) \cdot \PP(p(X) \in \interval_\iteratorInterval) \quad \text{for some }k\\
%     % changed = to <= when adding prediction range not covered
%     \le &~ \sum_{\iteratorInterval=1}^{\numInterval} \PP\Big(p^*(X) \notin E^\epsilon(\interval_\iteratorInterval)\mid p(X) \in \interval_\iteratorInterval\Big) \cdot \PP(p(X) \in \interval_\iteratorInterval) 
%     + \Big(1-\sum_{\iteratorInterval=1}^\numInterval\PP(p(X) \in \interval_\iteratorInterval)\Big)\\
%     \leq &~ \sum_{\iteratorInterval=1}^{\numInterval} \frac{\Var\Big(p^*(X) \mid p(X) \in \interval_\iteratorInterval\Big)}{\epsilon^2} \cdot \PP(p(X) \in \interval_\iteratorInterval) 
%     + \Big(1-\sum_{\iteratorInterval=1}^\numInterval\PP(p(X) \in \interval_\iteratorInterval)\Big).
%     % + \PP(p(X) \in \interval_{\numInterval+1})
% \end{align*}
% The penultimate equality follows from the law of total probability. The final inequality follows from Chebyshev's inequality and from the \titleShort definition which implies
% \begin{align*}
%     \E\Big[p^*(X) \mid p(X) \in \interval_\iteratorInterval\Big] \in \interval_\iteratorInterval.
% \end{align*}

% Recalling that $\J$ is a random variable following distribution $\PP(\J = j) = \PP(p(X) \in \interval_\iteratorInterval)$, we may rewrite
% \begin{align*}
%     \sum_{\iteratorInterval=1}^\numInterval \frac{\Var\Big(p^*(X) \mid p(X) \in \interval_\iteratorInterval\Big)}{\epsilon^2} \cdot \PP(p(X) \in \interval_\iteratorInterval)
%     = &~ \epsilon^{-2} \cdot \E\Big[\Var\Big(p^*(X) \mid p(X) \in \interval_\J\Big)\Big] \\
%     = &~ \epsilon^{-2} \cdot \Big(\Var\Big(p^*(X) \Big) - \Var\Big(\E\Big[p^*(X) \mid p(X) \in \interval_\J\Big]\Big)\Big).
% \end{align*}
% Here the penultimate equality follows from the law of total probability and the final equality follows from the law of total variance.
% \end{proof}




\subsection{Proof of Theorem~\ref{thm:tp}}
\label{proof_thm:tp}
\begin{proof}
Let $A_j=\{p(X)\in\interval_j\}$, $j\in[\numInterval]$, and
$A_0=\{p(X)\notin\cup_{j=1}^{\numInterval}\interval_j\}$.
These events form a partition. For each $j$ with $m_j>0$, Eq.~\eqref{eq:task-predicting_true_prob} gives $C_{\tp}(X;\epsilon)=E^\epsilon(\interval_j)$ on $A_j$.
The law of total probability therefore implies
\begin{align}\label{eq:task-predicting_true_prob_decomposition}
 \PP\Big(p^*(X)\notin C_{\tp}(X;\epsilon)\Big)~\leq\sum_{j:m_j>0}m_j
 \PP\Big(p^*(X)\notin E^\epsilon(\interval_j)\mid A_j\Big)
 +\PP(A_0).
\end{align}
For a positive-mass interval, let $\mu_j=\E[p^*(X)\mid A_j]$.
By the tower property and the \titleShort condition,
\begin{align*}
 \mu_j=\E[\E[Y\mid X]\mid A_j]=\E[Y\mid A_j]\in\interval_j.
\end{align*}
If $p^*(X)\notin E^\epsilon(\interval_j)$, its distance from every
point of $\interval_j$, including $\mu_j$, exceeds $\epsilon$.
Chebyshev's inequality conditionally on $A_j$ gives
\begin{align*}
 \PP\Big(p^*(X)\notin E^\epsilon(\interval_j)\mid A_j\Big)
 &\leq\PP\Big(|p^*(X)-\mu_j|>\epsilon\mid A_j\Big)\\
 &\leq\frac{\Var(p^*(X)\mid A_j)}{\epsilon^2}.
\end{align*}
Since $\PP(A_0)=1-\sum_{j=1}^{\numInterval}m_j$, substituting them into Eq.~\eqref{eq:task-predicting_true_prob_decomposition} yields
\begin{align*}
 \PP\Big(p^*(X)\notin C_{\tp}(X;\epsilon)\Big)~\leq\epsilon^{-2}\sum_{j:m_j>0}m_j
 \Var\Big(p^*(X)\mid p(X)\in\interval_j\Big)
 +\Big(1-\sum_{j=1}^{\numInterval}m_j\Big).
\end{align*}
This proves the theorem.
\end{proof}

\subsection{Proof of Theorem~\ref{thm:mc_error}}
\label{proof_thm:mc_error}


\begin{proof}
Fix a class $\iteratorClass\in[\numClass]$ and abbreviate
\[
\pi:=\pi_\iteratorClass,\qquad
t:=t_\iteratorClass,\qquad
J:=J_\iteratorClass,\qquad
S:=S_\iteratorClass,
\]
\[
m_j:=m_{\iteratorClass,j},\qquad
a_j:=a_{\iteratorClass,j},\qquad
b_j:=b_{\iteratorClass,j},
\qquad j=1,\ldots,J.
\]
Define the disjoint events
\[
A_0
:=
\left\{
p_\iteratorClass(X)\notin \bigcup_{j=1}^{J}[a_j,b_j]
\right\},
\qquad
A_j
:=
\bigl\{
p_\iteratorClass(X)\in [a_j,b_j]
\bigr\},
\quad j=1,\ldots,J.
\]
Let
\[
\mu_0
:=
\PP\bigl(Y(X)=\iteratorClass \mid A_0\bigr)\in[0,1],
\qquad
\mu_j
:=
\PP\bigl(Y(X)=\iteratorClass \mid A_j\bigr)\in[a_j,b_j],
\quad j=1,\ldots,J,
\]
where the inclusion $\mu_j\in[a_j,b_j]$ follows from the definition of \titleShorts.
Since the events $A_0,A_1,\ldots,A_J$ partition the sample space,
\begin{align}
\pi
=
\PP\bigl(Y(X)=\iteratorClass\bigr)
=
t\mu_0+\sum_{j=1}^{J} m_j\mu_j.
\label{eq:mc_partition_identity}
\end{align}

Now,
\begin{align}
&~\PP\Bigl(Y(X)\notin C_{\mc}(X;S),\,Y(X)=\iteratorClass\Bigr)
\nonumber\\
=
&~t\mu_0+\sum_{j\notin S} m_j\mu_j,
\label{eq:mc_miscoverage_numerator}
\end{align}
because class $\iteratorClass$ is omitted precisely on the uncovered region $A_0$ and on those intervals $A_j$ with $j\notin S$.
Dividing Eq.~\eqref{eq:mc_miscoverage_numerator} by $\pi$ gives
\begin{align}
\PP\Bigl(Y(X)\notin C_{\mc}(X;S)\mid Y(X)=\iteratorClass\Bigr)
=
\frac{t\mu_0+\sum_{j\notin S} m_j\mu_j}{\pi}.
\label{eq:mc_conditional_error_exact}
\end{align}

We first prove the upper-endpoint bound.
Since $\mu_0\le 1$ and $\mu_j\le b_j$ for every $j\notin S$,
\[
t\mu_0+\sum_{j\notin S} m_j\mu_j
\le
t+\sum_{j\notin S} m_j b_j.
\]
Combining with Eq.~\eqref{eq:mc_conditional_error_exact} yields
\[
\PP\Bigl(Y(X)\notin C_{\mc}(X;S)\mid Y(X)=\iteratorClass\Bigr)
\le
\Gamma_{\iteratorClass}^{\mathrm{up}}(S_\iteratorClass).
\]

Next we prove the lower-endpoint bound.
By Eqs.~\eqref{eq:mc_partition_identity} and \eqref{eq:mc_miscoverage_numerator},
\[
t\mu_0+\sum_{j\notin S} m_j\mu_j
=
\pi-\sum_{j\in S} m_j\mu_j.
\]
Since $\mu_j\ge a_j$ for all $j\in S$,
\[
t\mu_0+\sum_{j\notin S} m_j\mu_j
\le
\pi-\sum_{j\in S} m_j a_j.
\]
Combining again with Eq.~\eqref{eq:mc_conditional_error_exact} gives
\[
\PP\Bigl(Y(X)\notin C_{\mc}(X;S)\mid Y(X)=\iteratorClass\Bigr)
\le
\Gamma_{\iteratorClass}^{\mathrm{low}}(S_\iteratorClass).
\]
Therefore
\[
\PP\Bigl(Y(X)\notin C_{\mc}(X;S)\mid Y(X)=\iteratorClass\Bigr)
\le
\Gamma_{\iteratorClass}(S_\iteratorClass).
\]

It remains to prove exactness.
Define
\[
U:=t+\sum_{j\notin S} m_j b_j,
\qquad
L:=\pi-\sum_{j\in S} m_j a_j.
\]
We show that the miscoverage numerator in Eq.~ \eqref{eq:mc_miscoverage_numerator} can attain $\min\{U,L\}$.

\emph{Case 1: $U\le L$.}
Set
\[
\mu_0=1,
\qquad
\mu_j=b_j \ \text{ for all } j\notin S.
\]
Then the miscoverage numerator equals $U$.
To make Eq.~\eqref{eq:mc_partition_identity} hold, we need to choose $\mu_j\in[a_j,b_j]$ for $j\in S$ so that
\[
\sum_{j\in S} m_j\mu_j=\pi-U.
\]
This is possible because
\[
U\le L=\pi-\sum_{j\in S} m_j a_j
\quad \Longrightarrow \quad
\pi-U\ge \sum_{j\in S} m_j a_j,
\]
and, since any admissible law must satisfy
\[
\pi
\le
t+\sum_{j\notin S} m_j b_j+\sum_{j\in S} m_j b_j
=
U+\sum_{j\in S} m_j b_j,
\]
we also have
\[
\pi-U\le \sum_{j\in S} m_j b_j.
\]
Hence $\pi-U$ lies in the interval
\[
\left[\sum_{j\in S} m_j a_j,\ \sum_{j\in S} m_j b_j\right],
\]
so such a choice of $\{\mu_j\}_{j\in S}$ exists.
Therefore equality is attained with value
\[
\frac{U}{\pi}
=
\Gamma_{\iteratorClass}^{\mathrm{up}}(S_\iteratorClass)
=
\Gamma_{\iteratorClass}(S_\iteratorClass).
\]

\emph{Case 2: $L<U$.}
Set
\[
\mu_j=a_j \ \text{ for all } j\in S.
\]
Then the retained intervals contribute exactly $\sum_{j\in S} m_j a_j$, so by Eq.~\eqref{eq:mc_partition_identity} the miscoverage numerator must equal $L$.
It remains to show that one can choose $\mu_0\in[0,1]$ and $\mu_j\in[a_j,b_j]$ for $j\notin S$ so that
\[
t\mu_0+\sum_{j\notin S} m_j\mu_j=L.
\]
The attainable range of the left-hand side is
\[
\left[\sum_{j\notin S} m_j a_j,\ t+\sum_{j\notin S} m_j b_j\right].
\]
Because any admissible law must satisfy
\[
\pi\ge \sum_{j=1}^{J} m_j a_j,
\]
we have
\[
L
=
\pi-\sum_{j\in S} m_j a_j
\ge
\sum_{j\notin S} m_j a_j.
\]
And by assumption,
\[
L<U=t+\sum_{j\notin S} m_j b_j.
\]
Hence $L$ lies in the attainable range, so equality is attained with value
\[
\frac{L}{\pi}
=
\Gamma_{\iteratorClass}^{\mathrm{low}}(S_\iteratorClass)
=
\Gamma_{\iteratorClass}(S_\iteratorClass).
\]

This proves exactness in both cases.
\end{proof}

\subsection{Proof of Theorem~\ref{thm:mc_optimality}}
\label{proof_thm:mc_optimality}



\begin{proof}
By construction, $\Gamma_{\iteratorClass}\bigl(S_\iteratorClass^\star(\alpha_\iteratorClass)\bigr)\le \alpha_\iteratorClass$ for every $\iteratorClass\in[\numClass]$.
Applying Theorem~\ref{thm:mc_error} with $S=S^\star(\alpha)$ yields
\[
\PP\Bigl(Y(X)\notin C_{\mc}^\star(X;\alpha)\mid Y(X)=\iteratorClass\Bigr)
\le
\alpha_\iteratorClass,
\]
which proves the coverage claim.

We now prove size-optimality.
For any $S=(S_1,\ldots,S_\numClass)$,
\begin{align*}
\E\bigl[|C_{\mc}(X;S)|\bigr]
=
\sum_{\iteratorClass=1}^{\numClass}
\PP\Bigl(\iteratorClass\in C_{\mc}(X;S)\Bigr).
\end{align*}
Since the intervals in $\intervalset_{p_\iteratorClass}$ are disjoint for each fixed class $\iteratorClass$,
\[
\PP\Bigl(\iteratorClass\in C_{\mc}(X;S)\Bigr)
=
\sum_{j\in S_\iteratorClass}
\PP\bigl(p_\iteratorClass(X)\in [a_{\iteratorClass,j},b_{\iteratorClass,j}]\bigr)
=
\sum_{j\in S_\iteratorClass} m_{\iteratorClass,j}.
\]
Therefore
\begin{align}
\E\bigl[|C_{\mc}(X;S)|\bigr]
=
\sum_{\iteratorClass=1}^{\numClass}
\sum_{j\in S_\iteratorClass} m_{\iteratorClass,j}.
\label{eq:mc_expected_size}
\end{align}

The objective \eqref{eq:mc_expected_size} separates over classes, and the feasibility constraints
\[
\Gamma_{\iteratorClass}(S_\iteratorClass)\le \alpha_\iteratorClass,
\qquad
\iteratorClass\in[\numClass],
\]
also separate over classes.
Hence the global minimization problem decomposes into $\numClass$ independent classwise problems, and each classwise optimizer is exactly $S_\iteratorClass^\star(\alpha_\iteratorClass)$ by \eqref{eq:mc_population_opt}.
Summing the classwise minima proves the result.
\end{proof}

\subsection{Proof of Theorem~\ref{thm:mc_same_sample}}
\label{proof_thm:mc_same_sample}


\begin{proof}
Let $\mathcal{E}_{\mathrm{PI}}$ denote the event that all intervals returned by Algorithm~1 are genuine disjoint \titleShorts simultaneously.
By assumption,
\[
\PP(\mathcal{E}_{\mathrm{PI}})\ge 1-\delta_{\mathrm{PI}}.
\]

For each class $\iteratorClass\in[\numClass]$ and each candidate interval $I\in\mathcal{A}_{p_\iteratorClass}$, define
\[
Z_k^{(\iteratorClass,I)}
:=
\mathbbm{1}\bigl(p_\iteratorClass(X_k)\in I\bigr),
\qquad
k=1,\ldots,n_{\mathrm{cal}}.
\]
Conditioned on the data-generating distribution, these are i.i.d.\ Bernoulli random variables with mean
\[
m_{\iteratorClass}(I)
:=
\PP\bigl(p_\iteratorClass(X)\in I\bigr).
\]
Hence Hoeffding's inequality gives, for every fixed $(\iteratorClass,I)$,
\[
\PP\left(
\left|
\frac{1}{n_{\mathrm{cal}}}\sum_{k=1}^{n_{\mathrm{cal}}} Z_k^{(\iteratorClass,I)}
-
m_{\iteratorClass}(I)
\right|
>
\varepsilon_m
\right)
\le
2e^{-2n_{\mathrm{cal}}\varepsilon_m^2}.
\]
Applying a union bound over all candidate intervals in all classes yields
\[
\PP\left(
\sup_{\iteratorClass\in[\numClass]}
\sup_{I\in\mathcal{A}_{p_\iteratorClass}}
\left|
\frac{1}{n_{\mathrm{cal}}}\sum_{k=1}^{n_{\mathrm{cal}}} Z_k^{(\iteratorClass,I)}
-
m_{\iteratorClass}(I)
\right|
>
\varepsilon_m
\right)
\le
2N_{\mathrm{cand}}e^{-2n_{\mathrm{cal}}\varepsilon_m^2}
=
\delta_m.
\]
Therefore, with probability at least $1-\delta_m$, we have simultaneously for every class $\iteratorClass$ and every candidate interval $I\in \mathcal{A}_{p_\iteratorClass}$,
\[
\left|
\widehat m_{\iteratorClass}(I)-m_{\iteratorClass}(I)
\right|
\le
\varepsilon_m.
\]
Since each selected interval $\widehat I_{\iteratorClass,j}$ belongs to $\mathcal{A}_{p_\iteratorClass}$ by construction, it follows that
\[
m_{\iteratorClass,j}
:=
\PP\bigl(p_\iteratorClass(X)\in \widehat I_{\iteratorClass,j}\bigr)
\ge
\underline m_{\iteratorClass,j}
\qquad
\text{for all } \iteratorClass \text{ and } j.
\]
Let $\mathcal{E}_m$ denote this event.

Next, for each class $\iteratorClass$, the random variables $\mathbbm{1}(Y_k=\iteratorClass)$ are i.i.d.\ Bernoulli with mean $\pi_\iteratorClass$.
Another Hoeffding bound and a union bound over $\iteratorClass\in[\numClass]$ give
\[
\PP\left(
\sup_{\iteratorClass\in[\numClass]}
\bigl|
\widehat\pi_\iteratorClass-\pi_\iteratorClass
\bigr|
>
\varepsilon_\pi
\right)
\le
2\numClass\,e^{-2n_{\mathrm{cal}}\varepsilon_\pi^2}
=
\delta_\pi.
\]
Hence, with probability at least $1-\delta_\pi$, we simultaneously have
\[
\underline\pi_\iteratorClass
\le
\pi_\iteratorClass
\le
\overline\pi_\iteratorClass
\qquad
\text{for all }\iteratorClass\in[\numClass].
\]
Let $\mathcal{E}_\pi$ denote this event.

By the union bound,
\[
\PP\bigl(\mathcal{E}_{\mathrm{PI}}\cap \mathcal{E}_m\cap \mathcal{E}_\pi\bigr)
\ge
1-\delta_{\mathrm{PI}}-\delta_m-\delta_\pi.
\]
Fix a realization of $D_{\mathrm{cal}}$ in the event
\[
\mathcal{E}
:=
\mathcal{E}_{\mathrm{PI}}\cap \mathcal{E}_m\cap \mathcal{E}_\pi.
\]
Conditioned on this realization, the intervals $\widehat I_{\iteratorClass,j}=[\widehat a_{\iteratorClass,j},\widehat b_{\iteratorClass,j}]$ are now deterministic and genuine \titleShorts, so Theorem~\ref{thm:mc_error} applies to them.

Fix a class $\iteratorClass\in[\numClass]$ and a subset $S_\iteratorClass\subseteq[\widehat J_\iteratorClass]$.
By Theorem~\ref{thm:mc_error},
\begin{align}
&~
\PP\Bigl(
Y(X)\notin \widehat C_{\mc}(X;S)
\ \Big|\
Y(X)=\iteratorClass,\ D_{\mathrm{cal}}
\Bigr)
\nonumber\\
\le&~
\min\Biggl\{
\frac{
\widehat t_\iteratorClass
+\sum_{j\notin S_\iteratorClass}
\widehat m_{\iteratorClass,j}^{\mathrm{pop}}
\widehat b_{\iteratorClass,j}
}{
\pi_\iteratorClass
},
\ 
1-
\frac{
\sum_{j\in S_\iteratorClass}
\widehat m_{\iteratorClass,j}^{\mathrm{pop}}
\widehat a_{\iteratorClass,j}
}{
\pi_\iteratorClass
}
\Biggr\},
\label{eq:mc_empirical_start}
\end{align}
where
\[
\widehat m_{\iteratorClass,j}^{\mathrm{pop}}
:=
\PP\bigl(p_\iteratorClass(X)\in \widehat I_{\iteratorClass,j}\mid D_{\mathrm{cal}}\bigr),
\qquad
\widehat t_\iteratorClass
:=
1-\sum_{j=1}^{\widehat J_\iteratorClass}\widehat m_{\iteratorClass,j}^{\mathrm{pop}}.
\]
Using
\[
\widehat t_\iteratorClass
+
\sum_{j\notin S_\iteratorClass}
\widehat m_{\iteratorClass,j}^{\mathrm{pop}}
\widehat b_{\iteratorClass,j}
=
1-
\sum_{j\in S_\iteratorClass}\widehat m_{\iteratorClass,j}^{\mathrm{pop}}
-
\sum_{j\notin S_\iteratorClass}
(1-\widehat b_{\iteratorClass,j})
\widehat m_{\iteratorClass,j}^{\mathrm{pop}},
\]
together with $\widehat m_{\iteratorClass,j}^{\mathrm{pop}}\ge \underline m_{\iteratorClass,j}$ on $\mathcal{E}_m$, we obtain
\[
\widehat t_\iteratorClass
+
\sum_{j\notin S_\iteratorClass}
\widehat m_{\iteratorClass,j}^{\mathrm{pop}}
\widehat b_{\iteratorClass,j}
\le
1-
\sum_{j\in S_\iteratorClass}\underline m_{\iteratorClass,j}
-
\sum_{j\notin S_\iteratorClass}
(1-\widehat b_{\iteratorClass,j})
\underline m_{\iteratorClass,j}.
\]
Since $\pi_\iteratorClass\ge \underline\pi_\iteratorClass$ on $\mathcal{E}_\pi$, it follows that
\[
\frac{
\widehat t_\iteratorClass
+\sum_{j\notin S_\iteratorClass}
\widehat m_{\iteratorClass,j}^{\mathrm{pop}}
\widehat b_{\iteratorClass,j}
}{
\pi_\iteratorClass
}
\le
\widehat\Gamma_{\iteratorClass}^{\mathrm{up}}(S_\iteratorClass).
\]
Likewise, on $\mathcal{E}_m$ and $\mathcal{E}_\pi$,
\[
\sum_{j\in S_\iteratorClass}
\widehat m_{\iteratorClass,j}^{\mathrm{pop}}
\widehat a_{\iteratorClass,j}
\ge
\sum_{j\in S_\iteratorClass}
\underline m_{\iteratorClass,j}
\widehat a_{\iteratorClass,j},
\qquad
\pi_\iteratorClass\le \overline\pi_\iteratorClass,
\]
so
\[
1-
\frac{
\sum_{j\in S_\iteratorClass}
\widehat m_{\iteratorClass,j}^{\mathrm{pop}}
\widehat a_{\iteratorClass,j}
}{
\pi_\iteratorClass
}
\le
1-
\frac{
\sum_{j\in S_\iteratorClass}
\underline m_{\iteratorClass,j}
\widehat a_{\iteratorClass,j}
}{
\overline\pi_\iteratorClass
}
=
\widehat\Gamma_{\iteratorClass}^{\mathrm{low}}(S_\iteratorClass).
\]
Combining these inequalities with Ineq. \eqref{eq:mc_empirical_start} yields
\[
\PP\Bigl(
Y(X)\notin \widehat C_{\mc}(X;S)
\ \Big|\
Y(X)=\iteratorClass,\ D_{\mathrm{cal}}
\Bigr)
\le
\widehat\Gamma_{\iteratorClass}(S_\iteratorClass).
\]
Since the argument holds simultaneously for every class and every subset on the event $\mathcal{E}$, the theorem follows.

Finally, if $\widehat S_\iteratorClass^\star(\alpha_\iteratorClass)$ is feasible for every class, then
\[
\widehat\Gamma_{\iteratorClass}\bigl(\widehat S_\iteratorClass^\star(\alpha_\iteratorClass)\bigr)
\le
\alpha_\iteratorClass.
\]
Applying the displayed bound above with $S_\iteratorClass=\widehat S_\iteratorClass^\star(\alpha_\iteratorClass)$ gives
\[
\PP\Bigl(
Y(X)\notin \widehat C_{\mc}^\star(X;\alpha)
\ \Big|\
Y(X)=\iteratorClass,\ D_{\mathrm{cal}}
\Bigr)
\le
\alpha_\iteratorClass
\]
for every class $\iteratorClass\in[\numClass]$.
\end{proof}
\clearpage
\section{Implementation Details}
\subsection{Baselines}
\label{app:baselines}
\begin{enumerate}
\item \textbf{\titleShort.} We construct \titleShorts from the calibration split. For Task~1, We adapt the procedure of Eq.~\eqref{eq:task-predicting_true_prob} for an $\epsilon$-expansion of \titleShorts.  For Task~2, we use Eq.~\eqref{eq:mc_acceptance}. That is, include class if its predicted probability lies in the an acceptance region by the subset solver in Algorithm~\ref{alg:mc_population_optimal}.

\item \textbf{Simultaneous confidence interval.}
Suppose a point estimator $p_{\theta}$ is parametrized by $\theta \in \mathbb{R}^d$. 
Given covariate value $x$, denote $p_{\theta^*}(x)$ as the true conditional probability when correctly specified. The simultaneous confidence interval for $p_{\theta^*}(x)$ is computed from a Wald confidence region for $\theta^*$ (using the chi-square critical value).
Let $\hat{\theta}$ be the fitted parameter and let $\widehat V$ be an estimate of its covariance matrix.
Let
\[
p_\theta(x) = g^{-1}\!\bigl(z(x)^\top \theta\bigr),
\]
where $z(x)$ is the feature vector (for example, $z(x) = (1~x^\top)^\top$) and $g^{-1}$ is the inverse link function. 
The simultaneous confidence interval for $p_{\theta^*}(x)$ is
\[
\mathrm{SCI}_{1-\alpha}(x)
=
\left[
\inf_{\theta \in \mathcal C_{1-\alpha}} p_\theta(x),
\;
\sup_{\theta \in \mathcal C_{1-\alpha}} p_\theta(x)
\right],
\]
where 
\[
\mathcal C_{1-\alpha}
=
\left\{
\theta \in \mathbb{R}^d :
(\theta-\hat\theta)^\top \widehat V^{-1}(\theta-\hat\theta)
\le \chi^2_{d,\,1-\alpha}
\right\}.
\]

\item \textbf{Split conformal prediction.} For each test covariate input $x$, the split conformal interval is 
\[
\widehat C_{1-\alpha}(x)
=
\bigl[\hat p(x)-\hat q_{1-\alpha},\ \hat p(x)+\hat q_{1-\alpha}\bigr]\cap[0,1],
\]
where $\hat q_{1-\alpha}$ is an empirical $1-\alpha$ quantile of the residual scores $|Y_i-\hat p(X_i)|$ on $n_\mathrm{cal}$ calibration data points. The coverage satisfies
\begin{align}\label{eq:conformal_tp}
\PP\Big(Y \in \widehat C_{1-\alpha}(x) \mid \{X_i, Y_i\}_{i\in [n_\mathrm{cal}]}\Big) \in \Big[1-\alpha,1-\alpha+(1+n_\mathrm{cal})^{-1}\Big)
\end{align}
\citep{papadopoulos2002inductive}.
% For Task~2, we apply the same construction classwise 
% in a one-vs-rest manner. For each class $g$, we compute residual scores
% \[
% R_i^{(g)} = \bigl|\mathbf{1}\{Y_i=g\} - \hat p_g(X_i)\bigr|
% \]
% We then define the classwise upper score
% \[
% U_g(x) = \min\{\hat p_g(x)+\hat q_{1-\alpha}^{(g)},\,1\},
% \]
% and form the prediction set by thresholding at $\tau$:
% \[
% \widehat \Gamma_\tau(x) = \{g \in [K] : U_g(x) \ge \tau\}.
% \]
% For Task~2, we apply the same construction classwise to obtain upper scores and threshold them at $\tau$.

\item \textbf{Calibration-based interval.}
We discretize $[0,1]$ into bins $B_k=(b_{k-1},b_k]$ with midpoint $m_k=(b_{k-1}+b_k)/2$. 
We retain a subset of bins such that their mass-weighted absolute calibration error is close to that of \titleShort; if that bin is not retained, we use the retained bin whose midpoint is closest to $\hat p(x)$.

\item \textbf{Fixed-width binning.} We discretize $[0,1]$ into fixed bins and use the bin containing $\hat p(x)$ as the reported interval.

\item \textbf{Conformal classification.}
This an additional baseline only for Task~2 (which differs from the others by not relying on intervals of point predictions). Let $\hat p_g(x)$ denote the predicted probability of class $g$. 
Given a test covariate input $x$, the label set includes classes whose predicted probability exceeds a threshold $\hat \tau_\alpha$,
\[
\widehat \Gamma_{1-\alpha}(x)
=
\{g \in [K] : \hat p_g(x) \ge \hat \tau_\alpha\},
\]
where $\hat \tau_\alpha$ is obtained from the lower tail of scores on calibration data:
% \[
% \hat \tau_\alpha
% =
% \operatorname{Quantile}_{\lceil (n_{\mathrm{cal}}+1)\alpha\rceil / n_{\mathrm{cal}}}
% \bigl(\{p_{Y_i}(X_i) : i \in [n_\mathrm{cal}]\}\bigr).
% \]
\[
\hat \tau_\alpha = \lceil\alpha(n_\mathrm{cal} + 1)\rceil \text{ smallest of $\{\hat p_{Y_i}(X_i)\}_{ i\in{[n_\mathrm{cal}]}}$}.
\]
Analogous to Eq.~\eqref{eq:conformal_tp}, it is marginally guaranteed that on the test data 
\[
\PP\Big(Y \in \widehat \Gamma_{1-\alpha}(x)\mid \{X_i, Y_i\}_{i\in [n_\mathrm{cal}]}\Big) \in \Big[1-\alpha,1-\alpha+(1+n_\mathrm{cal})^{-1}\Big).
\]
\end{enumerate}




\paragraph{Interval subset solver for Task 2} For \titleShort, fixed-width binning, calibration-based intervals, and split conformal prediction, the classwise candidate intervals from the calibration set are passed to Algorithm~\ref{alg:mc_population_optimal}. For each class the solver keeps a min-mass subset whose plugin two-sided certificate meets the target \(\alpha\). Simultaneous confidence intervals and conformal classification bypass the solver. For simultaneous CIs we include class \(g\) if the upper endpoint satisfies \(U_g(x)\ge\tau\), with \(\tau\) the \(\lceil\alpha(n_{\mathrm{cal}}+1)\rceil\)-th smallest of \(\{U_{Y_i}(X_i)\}\). For conformal classification, we use the calibration set to select a threshold on point predictions. 



\subsection{Metrics}
\label{app:metrics}
For Task~1, we report the average interval length, calibration errors, the coverage of the prediction space $[0,1]$, and the coverage of the target probability $p^*(x)$ on the test data. For a method that outputs intervals $\widehat C(X_i)=[L_i,U_i]$ on evaluation points $\{X_i\}_{i=1}^{n_{\mathrm{eval}}}$, let $\mathcal I=\{I_1,\dots,I_M\}$ denote the set of distinct reported intervals, where $I_j=[\ell_j,u_j]$. For each interval $I_j$, define its midpoint as $m_j=(\ell_j+u_j)/2$ and let $A_j=\{i:\widehat C(X_i)=I_j\}$ be the set of evaluation points assigned to that interval. We estimate the expected calibration error of the intervals relative to their midpoints by:
\[
\mathrm{ECE_{mid}}
=
\frac{1}{M}\sum_{j=1}^M
\left|
\frac{1}{|A_j|}\sum_{i\in A_j} p^*(X_i)-m_j
\right|.
\]
We also estimate the expected calibration error of the intervals relative to the predictions $p(X_i)$ by:
\[
\mathrm{ECE_{pred}}
=
\frac{1}{M}\sum_{j=1}^M
\left|
\frac{1}{|A_j|}\sum_{i\in A_j} p^*(X_i)- \frac{1}{|A_j|}\sum_{i\in A_j} p(X_i)
\right|.
\]
Intuitively, $\mathrm{ECE_{mid}}$ assesses the discrepancy between an interval’s midpoint and the average target probability $p^*(X)$ among points assigned to that interval, capturing how well the interval is centered around the target probabilities.
$\mathrm{ECE_{pred}}$ measures the discrepancy between the average prediction $p(X)$ and the average target probability $p^*(X)$ among points assigned to each interval, representing if grouping align with the underlying predictions.
The reported relative ECE is the mean ECE of each method divided by the smallest mean ECE across all methods.
%Here relative ECE is each method's mean ECE divided by the smallest mean ECE across methods. 

For Task 2, we measure class-conditional error on label set $C(X;\tau)$,
\[
\frac{\sum_{i=1}^{n_{\text{eval}}} q_g(X_i)\mathbf{1}\{g\notin C(X_i;\tau)\}}
     {\sum_{i=1}^{n_{\text{eval}}} q_g(X_i)}.
\] We also report the average prediction-set size.




\subsection{Setup for Task 1}\label{app:task1_setup}

\paragraph{Univariate interval visualization (Figure~\ref{fig:task1_interval_visualization}).}
The covariate is sampled uniformly from a discrete grid $\mathrm{X\_GRID} = \{-2,\ldots,2\}$ with $1000$ equally spaced points, and the noise-free target probability used for evaluation is
\[
p^*(x) = \sigma(2x),
\]
where $\sigma(t) = 1/(1+e^{-t})$.
Observed labels are generated from the noisy logistic model
\[
Y \mid X=x,\xi \sim \mathrm{Bernoulli}\!\left(\sigma(2x + \xi)\right),
\qquad
\xi \sim N(0, 0.1^2).
\]
Thus the marginal conditional label probability is $\PP(Y=1\mid X=x)=\mathbb E_\xi[\sigma(2x+\xi)]$, whereas the reported target-probability coverage uses the noise-free $p^*(x)=\sigma(2x)$ above.
We set 
\[
n_{\mathrm{train}} = 2000,\qquad
n_{\mathrm{cal}} = 2000,\qquad
n_{\mathrm{eval}} = 10000.
\]
The common method settings are $\alpha = 0.1$, PICPI with \texttt{num\_bin}$=50$ in empirical mode, PICPI expansion parameter $\epsilon = 0.02$, and $50$ bins for the calibration-based and fixed-width baselines. The shared point predictor is scikit-learn logistic regression with its default $\ell_2$ penalty and regularization strength $1$.

\paragraph{Misspecified decision-tree illustration (Figure~\ref{fig:task1_misspecified_tree}).}
For the qualitative misspecification figure, the covariate is drawn from
\[
X \sim \mathrm{Unif}[-2,2],
\]
with true probability
\[
p^*(x) = \sigma\!\left(1.2\sin(2.5x) + 0.8x\right),
\]
and labels
\[
Y \mid X=x \sim \mathrm{Bernoulli}(p^*(x)).
\]
The point predictor is a deliberately misspecified decision tree with \texttt{max\_depth = 5} and \texttt{min\_samples\_leaf = 25}, fitted on
\[
n_{\mathrm{train}} = 20000,\qquad
n_{\mathrm{cal}} = 20000
\]
samples. PICPI is built with \texttt{num\_bin}$=100$ candidate grid points in empirical mode and expanded by $\epsilon=0.02$. The curves are plotted on a dense evaluation grid of $600$ evenly spaced points in $[-2,2]$.

\paragraph{Multivariate result table (Table~\ref{tab:task1_multivariate_mc_summary}).}
The covariate is
\[
X \sim N(0, I_{20}),
\]
and the true target probability is
\[
p^*(x) = \sigma(\beta_0 + x^\top \beta),
\qquad
\beta_0 = 0,
\qquad
\beta = (1.6,\,-1.1,\,0.9,\,0,\ldots,0)^\top \in \mathbb{R}^{20}.
\]
Labels are again generated from a noisy logistic model,
\[
Y \mid X=x,\xi \sim \mathrm{Bernoulli}\!\left(\sigma(\beta_0 + x^\top \beta + \xi)\right),
\qquad
\xi \sim N(0, 0.1^2).
\]
As in the univariate experiment, the reported target-probability coverage uses the noise-free $p^*(x)$ above, not the Gaussian-averaged marginal conditional label probability.
Each replication uses
\[
n_{\mathrm{train}} = 2000,\qquad
n_{\mathrm{cal}} = 2000,\qquad
n_{\mathrm{eval}} = 10000,
\]
and the table is averaged over $100$ Monte Carlo replications. The shared point predictor is scikit-learn logistic regression with its default $\ell_2$ penalty. Unless otherwise stated, the method hyperparameters are the same as in the univariate setting: $\alpha = 0.1$, PICPI with \texttt{num\_bin}$=50$ candidate grid points and $\varepsilon_{\mathrm{PICPI}} = 0.02$, and $50$ bins for the two binning baselines.

\paragraph{Split conformal implementation.}
For the Task~1 split conformal baseline, the implementation computes residual scores $R_i=|Y_i-\hat p(X_i)|$ on the calibration split and sets
\[
u=\min\!\left\{\frac{\lceil(1-\alpha)(n_{\mathrm{cal}}+1)\rceil}{n_{\mathrm{cal}}},1\right\},
\qquad
\hat q_{1-\alpha}=\operatorname{Quantile}^{\mathrm{linear}}_u
\bigl(\{R_i\}_{i\in[n_{\mathrm{cal}}]}\bigr).
\]
Here $\operatorname{Quantile}^{\mathrm{linear}}$ denotes the default linear-interpolation convention of \texttt{numpy.quantile}. This numerical implementation uses an interpolated empirical quantile rather than the exact conformal order statistic.

\paragraph{Empirical mode diagnostics (Figure~\ref{fig:emp_mode_diagnostics}).}
\label{app:task1_setup_emp_mode_diagnostics} 
The covariate and noisy-logit DGP, training size $n_{\mathrm{train}}=2000$, and fitted logistic model are the same as for Table~\ref{tab:task1_multivariate_mc_summary}. We compare empirical-mode and population-mode PICPI using $100$ equal-width score bins, with $\delta=0.1$ for population mode. We vary the calibration size over $\{10^3,3\times10^3,10^4,3\times10^4,10^5,3\times10^5,10^6,3\times10^6,10^7\}$, use $n_{\mathrm{eval}}=10^6$, and run $50$ replications for each calibration size. In these diagnostics, the held-out true probability is the Gaussian-averaged marginal probability $\mathbb E_\xi[\sigma(\beta_0+x^\top\beta+\xi)]$.

\subsection{Setup for Task 2}\label{app:task2_setup}


\paragraph{Coverage sweep for two DGPs (Figure~\ref{fig:task2})}
 Both use a $d = 10$ dimensional covariate and $\numClass = 4$ classes. Given the true class-probability vector $q(x) \in [0,1]^{\numClass}$, labels are drawn as
\[
Y \mid X = x \sim \mathrm{Categorical}\big(q(x)\big).
\]
Each replication uses disjoint splits of sizes
\[
n_{\mathrm{train}} = n_{\mathrm{cal}} = n_{\mathrm{eval}} = 2000,
\]
and all curves are averages over $1000$ independent replications. \titleShort uses \texttt{num\_bin}$=100$ equally spaced candidate grid points in empirical mode, while the calibration-based and fixed-width baselines use $100$ bins. Candidate intervals from \titleShort, split conformal prediction, calibration-based intervals, and fixed-width binning are passed to the classwise interval-selection solver (Algorithm~\ref{alg:mc_population_optimal}). For split conformal prediction, each class uses at most $200$ interval centers sampled without replacement from its calibration scores. Simultaneous confidence intervals and conformal classification instead use a calibration threshold, as described in Section~\ref{sec:experiment}. The coverage sweep evaluates all methods at $16$ equally spaced target coverage levels $1 - \alpha \in \{0, 1/15, \ldots, 1\}$.

\paragraph{DGP 1: Linear-Gaussian.}
The covariate is drawn from
\[
X \sim N(0, I_{10}),
\]
and the true class probabilities follow a linear softmax model,
\[
q(x) = \mathrm{softmax}(W x + b),
\]
with a sparse weight matrix
{\small
\[
\setlength{\arraycolsep}{3pt}
W =
\begin{bmatrix}
1.4 & -1.0 & 0.8 & 0 & 0 & 0 & 0 & 0 & 0 & 0 \\
0 & -0.7 & 0 & 1.1 & 0.8 & 0 & 0 & 0 & 0 & 0 \\
0 & 0 & 0.9 & 0 & 0 & -1.2 & 1.0 & 0 & 0 & 0 \\
-0.6 & 0 & 0 & 0 & 0 & 0 & 0 & 0.8 & -0.9 & 1.2
\end{bmatrix}
\]
}%
and an intercept $b = (0.25,\,-0.20,\,0.15,\,-0.20)^\top$.
The multinomial logistic working model is therefore correctly specified.

\paragraph{DGP 2: Nonlinear-mixture.}
The covariate is drawn from a four-component latent-factor mixture with component weights
\[
(0.32,\, 0.24,\, 0.27,\, 0.17).
\]
Equivalently, a latent component label is first drawn as
\[
k \sim \mathrm{Categorical}(0.32,\, 0.24,\, 0.27,\, 0.17).
\]
Conditional on component $k\in\{1,\ldots, 4\}$, an initial draw $\widetilde{X}\in\mathbb{R}^{10}$ is generated as
\[
\widetilde{X} = \mu_k + \Lambda^\top Z + \mathrm{diag}(\sigma_k)\, h(Z)\, S\, \varepsilon,
\qquad
Z \sim N(0, I_2),
\quad
\varepsilon \sim N(0, I_{10}),
\]
where $S = \sqrt{7/\chi^2_7}$ induces Student-$t$ (heavy) tails and $h(z) = 1 + 0.25\tanh(z_1 + 0.45 z_2)$ induces heteroscedasticity in the latent factors. The component-specific location and scale vectors $\mu_k, \sigma_k \in \mathbb{R}^{10}$ are the rows of
{\small
\begin{gather*}
\setlength{\arraycolsep}{3pt}
M =
\begin{bmatrix}
-1.4 & 0.5 & -0.9 & 0.2 & -0.3 & 0.6 & -0.5 & 0.4 & -0.7 & 0.3 \\
-0.2 & -1.0 & 0.7 & -0.6 & 0.8 & -0.4 & 0.3 & -0.5 & 0.6 & -0.8 \\
0.9 & 0.6 & -0.2 & 1.1 & -0.7 & 0.4 & 0.9 & 0.7 & -0.1 & 0.5 \\
1.6 & -0.3 & 1.0 & 0.5 & 0.4 & -0.8 & -1.0 & -0.6 & 0.9 & -0.4
\end{bmatrix},
\\
\setlength{\arraycolsep}{3pt}
V =
\begin{bmatrix}
0.60 & 0.45 & 0.55 & 0.40 & 0.50 & 0.35 & 0.45 & 0.42 & 0.50 & 0.38 \\
0.45 & 0.65 & 0.50 & 0.60 & 0.55 & 0.40 & 0.35 & 0.48 & 0.44 & 0.57 \\
0.55 & 0.50 & 0.70 & 0.45 & 0.60 & 0.50 & 0.55 & 0.52 & 0.46 & 0.50 \\
0.50 & 0.40 & 0.60 & 0.55 & 0.45 & 0.65 & 0.60 & 0.58 & 0.54 & 0.47
\end{bmatrix},
\end{gather*}
}%
and the shared factor-loading matrix is
{\small
\[
\setlength{\arraycolsep}{3pt}
\Lambda =
\begin{bmatrix}
0.95 & -0.55 & 0.35 & 0.45 & -0.30 & 0.40 & -0.20 & 0.32 & -0.28 & 0.24 \\
0.30 & 0.60 & -0.45 & 0.20 & 0.55 & -0.35 & 0.65 & -0.25 & 0.48 & -0.42
\end{bmatrix}.
\]
}%
The covariate $X$ is then obtained from $\widetilde{X}$ by fixed deterministic warps. Starting from $x = \tilde{x} + 0.22\sin(\tilde{x}_1)\,\delta_k$, the coordinates are updated sequentially (each line using the current values) via
\begin{align*}
x_3 &\leftarrow x_3 + 0.28\,(x_1^2 - 1), &
x_5 &\leftarrow x_5 + 0.16\, x_2 x_4, &
x_7 &\leftarrow x_7 + 0.14\, x_3 x_6, \\
x_8 &\leftarrow x_8 + 0.18\sin(x_1 x_3), &
x_9 &\leftarrow x_9 + 0.15\, x_5 x_7, &
x_{10} &\leftarrow x_{10} + 0.12\,(x_8^2 - 1),
\end{align*}
followed by the component offset $x \leftarrow x + o_k$ and the final warps
\begin{align*}
x_6 &\leftarrow x_6 + 0.20\,\mathrm{sgn}(x_1)\sqrt{|x_1| + 10^{-6}}, &
x_2 &\leftarrow x_2 - 0.15\tanh(x_7 - x_5), \\
x_9 &\leftarrow x_9 + 0.18\,\mathrm{sgn}(x_8)\sqrt{|x_8| + 10^{-6}}, &
x_{10} &\leftarrow x_{10} - 0.12\tanh(x_9 - x_6).
\end{align*}
The drift and offset vectors $\delta_k, o_k \in \mathbb{R}^{10}$ are the rows of
{\small
\begin{gather*}
\setlength{\arraycolsep}{3pt}
D =
\begin{bmatrix}
0.30 & -0.20 & 0.10 & 0.00 & 0.12 & -0.08 & 0.05 & 0.07 & -0.06 & 0.04 \\
-0.10 & 0.28 & -0.12 & 0.08 & -0.06 & 0.14 & -0.04 & -0.09 & 0.10 & -0.05 \\
0.08 & -0.12 & 0.26 & -0.16 & 0.10 & -0.06 & 0.18 & 0.11 & -0.04 & 0.09 \\
-0.14 & 0.06 & -0.08 & 0.24 & -0.10 & 0.20 & -0.22 & -0.05 & 0.12 & -0.13
\end{bmatrix},
\\
\setlength{\arraycolsep}{3pt}
O =
\begin{bmatrix}
-0.12 & 0.06 & -0.04 & 0.02 & -0.05 & 0.03 & -0.02 & 0.04 & -0.03 & 0.01 \\
0.05 & -0.10 & 0.08 & -0.06 & 0.09 & -0.04 & 0.03 & -0.02 & 0.05 & -0.07 \\
0.10 & 0.04 & -0.06 & 0.11 & -0.08 & 0.06 & 0.12 & 0.03 & -0.09 & 0.08 \\
0.16 & -0.02 & 0.10 & 0.07 & 0.05 & -0.11 & -0.14 & -0.08 & 0.07 & -0.05
\end{bmatrix}.
\end{gather*}
}%
The resulting marginal covariate law is non-Gaussian, multimodal, and dependent across coordinates.
The true class probabilities follow a softmax model whose linear logits are rescaled by a covariate-dependent scale function,
\begin{gather*}
q(x) = \mathrm{softmax}\!\left(\frac{W x + b}{s(x)}\right),
\\
s(x) = \mathrm{clip}\big(0.86 + 0.22\tanh(0.60 x_1 - 0.35 x_4 + 0.20 x_9),\ 0.65,\ 1.40\big),
\end{gather*}
where the weight matrix $W \in \mathbb{R}^{4 \times 10}$ and intercept $b \in \mathbb{R}^4$ were generated once (by combining i.i.d.\ Gaussian entries with structured rank-one terms drawn from a fixed seed) and held constant across all replications; rounded to two decimals,
{\small
\[
\setlength{\arraycolsep}{3pt}
W =
\begin{bmatrix}
-1.49 & 1.57 & -0.40 & -1.12 & 1.37 & -0.67 & 0.81 & -0.92 & 0.96 & 0.07 \\
-0.10 & 0.32 & -0.75 & -0.45 & 0.29 & -0.14 & 0.54 & -0.37 & -0.30 & -0.10 \\
0.32 & -0.65 & 0.47 & -0.05 & 0.11 & -0.12 & -0.44 & 0.32 & 0.16 & 0.57 \\
1.53 & -1.48 & 1.06 & 1.23 & 0.05 & 1.37 & -0.50 & 0.91 & 0.18 & 0.41
\end{bmatrix}
\]
}%
and $b = (0.30,\, 0.16,\, -0.05,\, -0.41)^\top$.
Because the log-odds are nonlinear in $x$ through $s(x)$, the multinomial logistic working model is misspecified and thus is not expected to be well calibrated under this DGP.


\subsection{Setup for Figure~\ref{fig:thm52_box_rate}}

\label{app:verification-consistency-prediction-setup}

The covariate and noisy-label DGP are the same as in the univariate interval visualization of Section~\ref{app:task1_setup}. We use Algorithm~\ref{alg:calibration_population} with \(100\) bins and \(\delta=0.1\) and \(100\) repetitions.
The noise-free target is $p^*(x)=\sigma(2x)$. We fit an unpenalized logistic-regression estimator once using a training set of size $10,000$, fix an evaluation set of size $10,000$, and vary the calibration-set size over $\{10^3,3\times10^3,10^4,3\times10^4,10^5,3\times10^5,10^6,3\times10^6,10^7\}$. 
\begin{enumerate}
\item fit a logistic-regression estimator \(\hat p\) on the training fold;
\item apply the calibration procedure on the calibration fold to obtain an interval collection \(\mathcal I=\{[a_j,b_j]\}_j\);
\item evaluate the interval system on the independent evaluation fold by computing, for each point, the minimal width among all intervals that contain the target probability \(p^*(X_i)\), with width \(1.0\) assigned if no interval contains \(p^*(X_i)\).
\end{enumerate}
%For each trained model, we also estimate the mean point-estimation error \(\hat\varepsilon=\mathbb E|\hat p(X)-p^\star(X)|\) on an auxiliary Monte Carlo sample of size \(100{,}000\) drawn from the same marginal law of \(X\). 


\section{Additional Algorithms}\label{app:alg}

Algorithm~\ref{alg:max_coverage_intervals} selects a disjoint subset for intervals returned by Algorithm~\ref{alg:calibration_population}. The \textsc{DisjointCalibration} procedure first imposes a maximum length and
then uses weighted interval scheduling to select a non-overlapping subcollection with maximum covered length. 

The \textsc{MaxCoverageIntervals} procedure solves a weighted interval scheduling problem to select a set of non-overlapping intervals that maximizes total coverage, where each interval is weighted by its length. It sorts intervals by end time, precomputes the last compatible interval for each candidate, and then uses dynamic programming to recursively compute the optimal value, followed by a backtracking step to recover the selected intervals.




\begin{algorithm*}[h]
   \caption{Maximum-coverage non-overlapping intervals (weighted interval scheduling)}
   \label{alg:max_coverage_intervals}
   \small
\begin{algorithmic}[1]
\Procedure{\textsc{DisjointCalibration}}{length bound $l$, \titleShorts $\intervalset_p$}
\State
\Comment{Restrict \titleShorts to intervals of length $\le l$ and maximize covered length}
\State $\intervalset_p^{(\le l)} \la \emptyset$
\For{{\bf each} interval $[a,b] \in \intervalset_p$}
    \If{$b - a \le l$}
        \State $\intervalset_p^{(\le l)} \la \intervalset_p^{(\le l)} \cup \{[a,b]\}$
    \EndIf
\EndFor
\If{$\intervalset_p^{(\le l)} = \emptyset$}
    \State {\bf Return} $\emptyset$ \Comment{No interval satisfies the length constraint}
\EndIf
\State Relabel $\intervalset_p^{(\le l)}$ as $\{[s_i, e_i]\}_{i=1}^m$
\State $\mathcal{S}^\star \la \textsc{MaxCoverageIntervals}\big(\{[s_i, e_i]\}_{i=1}^m\big)$

\Comment{Select a disjoint set intervals of maximal coverage}
\State {\bf Return} $\mathcal{S}^\star$
\EndProcedure
\\
\Procedure{\textsc{MaxCoverageIntervals}}{Intervals $\{[s_i, e_i]\}_{i=1}^m$}

\Comment{Find a maximum-length set of non-overlapping intervals}
\State Sort the intervals so that $e_1 \le e_2 \le \cdots \le e_m$
\For{$j \la 1$ {\bf to} $m$}
    \State $w_j \la e_j - s_j$ \Comment{Weight = interval length}
\EndFor
\For{$j \la 1$ {\bf to} $m$}
    \State $p(j) \la \max\{\, i < j : e_i \le s_j \,\}$, or $0$ if no such $i$ exists
    \Comment{Binary search over $\{e_1,\dots,e_{j-1}\}$}
\EndFor
\State Initialize array $\text{OPT}[0..m]$ \Comment{Dynamic Programming}
\State $\text{OPT}[0] \la 0$
\For{$j \la 1$ {\bf to} $m$}
    \State $\text{OPT}[j] \la \max\bigl(w_j + \text{OPT}[p(j)],~\text{OPT}[j-1]\bigr)$\Comment{Maximum coverage using $\{1,2,\dots,j\}$}
\EndFor
\State Initialize selected set $\mathcal{S}^\star \la \emptyset$
\State $j \la m$
\While{$j > 0$}
    \If{$w_j + \text{OPT}[p(j)] > \text{OPT}[j-1]$}
        \State $\mathcal{S}^\star \la \mathcal{S}^\star \cup \{[s_j, e_j]\}$
        \State $j \la p(j)$
    \Else
        \State $j \la j - 1$
    \EndIf
\EndWhile
\State {\bf Return} $\mathcal{S}^\star$
\EndProcedure



\end{algorithmic}
\end{algorithm*}
% Population version
%Algorithm~\ref{alg:calibration_population} finds all \titleShorts that satisfies Definition~\ref{def:self-calibrated}.






\end{document}
